% Informe final del proyecto DeepSpine. Compilar desde docs/informe/:
%   pdflatex informe_final && bibtex informe_final && pdflatex informe_final && pdflatex informe_final
\documentclass[11pt,a4paper]{article}
\usepackage[utf8]{inputenc}
\usepackage[T1]{fontenc}
\usepackage[spanish,es-noshorthands]{babel}
\usepackage{amsmath,amssymb}
\usepackage{graphicx,booktabs,xcolor,caption,url}
\usepackage{tikz,pgfplots}
\usetikzlibrary{positioning,arrows.meta}
\pgfplotsset{compat=1.16}
\usepackage[margin=2.5cm]{geometry}
\usepackage[hidelinks]{hyperref}
% Reduce una ecuacion solo si no cabe en el ancho del texto
\newsavebox{\cajaec}
\newcommand{\ajusta}[1]{\sbox{\cajaec}{$\displaystyle #1$}%
  \ifdim\wd\cajaec>\linewidth\resizebox{\linewidth}{!}{\usebox{\cajaec}}\else\usebox{\cajaec}\fi}
\graphicspath{{../guias/}{figuras/}{../resultados/}{../modelo_matematico/figuras/}}

\title{DeepSpine: clasificación de discos lumbares en resonancia magnética sagital T2\\[4pt]
\large Informe final del proyecto de Visión Artificial}
\author{Jhon Jairo Perdomo Perdomo \and Julián Andrés Silva Morales \and Juan Manuel Tengonó Lasso}
\date{30 de octubre de 2026}

\begin{document}
\maketitle
\begin{abstract}
DeepSpine clasifica los discos intervertebrales L1-L2 a L5-S1 de una resonancia sagital T2 en tres clases: Normal, Degeneración discal y Protrusión. Este informe recoge la preparación del dataset de 35\,007 imágenes, la selección de cuatro modelos, su modelo matemático y su entrenamiento, el MVP de demostración, la simulación de resultados, la aplicación web para el radiólogo, las correcciones hechas a esa aplicación y las pruebas finales. El mejor modelo por validación, una ResNet-50 ajustada, obtuvo una exactitud de 88,9\,\% y un F1 macro de 0,769 sobre las 768 imágenes originales de prueba. Es una herramienta académica, no diagnóstica: las etiquetas las asignaron estudiantes, sin radiólogo.
\end{abstract}
\tableofcontents
\clearpage

% =====================================================================
% Seccion: Preparacion del dataset (Hito 2)
% Fragmento SIN preambulo: se incluye con % =====================================================================
% Seccion: Preparacion del dataset (Hito 2)
% Fragmento SIN preambulo: se incluye con % =====================================================================
% Seccion: Preparacion del dataset (Hito 2)
% Fragmento SIN preambulo: se incluye con \input{secciones/preparacion_dataset}
% Paquetes requeridos: graphicx, booktabs, xcolor, tikz (+ calc no), pgfplots, caption
% Figuras: guia_visual_5rois.png vive en docs/guias/ (ver \graphicspath en el archivo principal)
% Las cifras salen de data/processed/etiquetado/etiquetas_revisadas.csv,
% data/processed/dataset_35k/manifest_35k.csv y de los scripts de scripts/.
% Todo lo que no se pudo verificar se marca con \pendiente{...} (se ve en rojo).
% =====================================================================
\providecommand{\pendiente}[1]{\textcolor{red}{\textbf{[PENDIENTE: #1]}}}
\providecommand{\arch}[1]{\texttt{\detokenize{#1}}}

\section{Preparación del dataset}
\label{sec:dataset}

Esta sección describe, en orden, cómo se pasó del conjunto público original a los 35\,007 recortes (ROI) con los que se entrenaron y compararon los cuatro modelos. La Figura~\ref{fig:flujo} resume el flujo y el número de elementos que pasan por cada etapa. Todas las cifras proceden de los archivos versionados del repositorio; lo que no pudo verificarse se señala como pendiente.

\begin{figure}[t]
\centering
\begin{tikzpicture}[
  box/.style={draw, rounded corners=2pt, align=left, text width=0.60\linewidth, inner sep=3pt, font=\footnotesize},
  num/.style={align=left, text width=0.30\linewidth, font=\footnotesize\bfseries},
  arrow/.style={-{Stealth[length=2mm]}, thick},
  node distance=4.2mm]
  \node[box] (a) {1. Conjunto público (Mendeley), archivos \texttt{.IMA}};
  \node[num, right=3mm of a] {575 carpetas, $\approx$48\,000 archivos};
  \node[box, below=of a] (b) {2. Descarte de series inútiles y conservación de T2 sagital};
  \node[num, right=3mm of b] {solo T2 sagital};
  \node[box, below=of b] (c) {3. Conversión DICOM$\to$PNG y limpieza manual de cortes};
  \node[num, right=3mm of c] {1\,086 cortes, 414 pacientes};
  \node[box, below=of c] (d) {4. Recorte L1--S1 y detección de discos (ROI 128$\times$128)};
  \node[num, right=3mm of d] {5\,385 ROI};
  \node[box, below=of d] (e) {5. Preetiquetado, revisión humana (A/B) y exclusiones};
  \node[num, right=3mm of e] {5\,123 ROI, 399 pacientes};
  \node[box, below=of e] (f) {6. División por paciente 70/15/15};
  \node[num, right=3mm of f] {279 / 60 / 60 pacientes};
  \node[box, below=of f] (g) {7. Aumento dentro de cada partición (balance por conteo)};
  \node[num, right=3mm of g] {35\,007 imágenes};
  \foreach \i/\j in {a/b,b/c,c/d,d/e,e/f,f/g} \draw[arrow] (\i) -- (\j);
\end{tikzpicture}
\caption{Flujo de preparación del dataset y número de elementos en cada etapa.}
\label{fig:flujo}
\end{figure}

\subsection{Origen de los datos}
Se usó el \emph{Lumbar Spine MRI Dataset} (Mendeley Data) \cite{sudirman2019}, el mismo conjunto público empleado en trabajos previos sobre clasificación de hernia discal \cite{sari2024}. Los estudios se adquirieron con un único equipo, un Siemens MAGNETOM Essenza de 1,5~T, y se distribuyen como archivos DICOM con extensión \arch{.IMA}, en una carpeta por paciente. Los cortes tienen un tamaño de píxel de 0,729~mm, un grosor de 4~mm y una separación entre centros de 4,8~mm. El campo \texttt{PatientID} está vacío en todo el corpus y el repositorio solo conserva identificadores de estudio y de serie (los campos de nombre, nacimiento y fecha están vacíos en \arch{meta_dicom_por_slice.csv}). El conjunto descargado contenía 575 carpetas de pacientes. Se eliminaron las que, en lugar de resonancia lumbar, contenían estudios del segmento cervical o de rodillas, y las lumbares con aberraciones o manchas que habrían introducido ruido en los modelos. Tras esa depuración quedaron 414 pacientes con cortes útiles y, tras las exclusiones de calidad de los ROI (Sección~\ref{sec:roi}), los \textbf{399 pacientes} del dataset final. Cada una de las 414 carpetas restantes corresponde a un \texttt{StudyInstanceUID} distinto y no se repite ningún \texttt{StudyTime}; no hay evidencia de pacientes duplicados, aunque con el identificador vacío no puede descartarse del todo.
% VERIFICAR: el sitio de Mendeley describe 515 pacientes; el equipo reporta 575 carpetas descargadas.

\subsection{Depuración de la fuente}
La descarga inicial contenía cerca de 48\,000 registros en formato \arch{.IMA}. Cada paciente trae varias series; solo una parte sirve para este problema. La Tabla~\ref{tab:descartes} lista lo que se descartó y por qué.

\begin{table}[t]
\centering\small
\caption{Series descartadas de cada estudio y motivo.}
\label{tab:descartes}
\begin{tabular}{p{0.30\linewidth}p{0.60\linewidth}}
\toprule
Serie & Motivo del descarte \\
\midrule
\arch{LOCALIZER_0001} & Localizadores de muy baja resolución usados por el técnico para calibrar la mesa y posicionar al paciente; no aportan información anatómica ni de textura. \\
\arch{POSDISP_[4]_T2_TSE_TRA_...} & Vistas sagitales con líneas blancas superpuestas (\emph{overlays}) que marcan los planos transversales; ese artefacto distorsiona descriptores de textura (GLCM/LBP) y confunde a las redes. \\
\arch{T1_TSE_SAG_...}, \arch{T1_TSE_TRA_...} & En T1 la grasa es brillante y el líquido oscuro, de modo que no permite evaluar la hidratación del disco (base de la escala de Pfirrmann); además, mezclar T1 y T2 altera los rangos de brillo. \\
\arch{T2_TSE_TRA_...} (axial) & Cortes axiales fragmentados en lotes de 3 a 5 por nivel; exigirían un emparejamiento espacial adicional por disco, mientras que el plano sagital cubre L1--S1 en una sola imagen. Se deja para una fase futura. \\
\bottomrule
\end{tabular}
\end{table}

El proceso se hizo en cuatro pasos, cada uno con su script (carpeta \arch{scripts/preprocesamiento/}):
\begin{enumerate}
  \item \arch{ingestion.py} recorre las carpetas de pacientes y copia solo las imágenes T2 sagitales, identificadas por el nombre de la serie.
  \item \arch{limpiar_slices.py} conserva el 40\,\% central del volumen sagital (se descarta el 30\,\% inicial y el 30\,\% final), donde el canal medular y la región L1--S1 son visibles.
  \item \arch{convertidor_dcm_a_png.py} convierte a PNG de 8 bits con una sola escala por serie: percentiles 1 y 99,5 calculados sobre los píxeles de tejido. Se decidió no aplicar CLAHE ni otro realce, para no alterar la señal relativa del T2 que sustenta el criterio diagnóstico.
  \item Depuración manual de los PNG: algunos pacientes tenían estudios del segmento cervical o de rodillas, y otros cortes T2 sagitales quedaban oscuros o con aberraciones; se eliminaron.
\end{enumerate}
Tras esta depuración quedaron \textbf{1\,086 cortes sagitales de 414 pacientes} (2,6 cortes por paciente en promedio; entre 1 y 6).

\subsection{Recorte de la región lumbar y de los ROI por disco}
\label{sec:roi}
\paragraph{Recorte L1--S1.} Con \arch{aplicar_recortes.py} se recortó cada corte en altura para abarcar principalmente L1--S1, replicando por \emph{template matching} unos recortes manuales previos (mediana de la puntuación de coincidencia 0,9996; solo el corte 0547\_008 quedó por debajo de 0,99). El recorte se extendió 60~px hacia arriba y hacia abajo hasta el borde de la imagen; con ello el relleno necesario en L1-L2 bajó de 779 de 1\,080 recortes (72\,\%) a cero.

\paragraph{Detector de discos.} Para localizar los discos se entrenó un detector con un conjunto propio en Roboflow (proyecto \emph{Intervertebral Disks}, versión~1): 100 cortes (70 de entrenamiento, 20 de validación y 10 de prueba), con orientación automática (\emph{auto-orient}) y reescalado a 640$\times$640 con bordes negros, sin aumento de datos en Roboflow. Las cinco clases del detector son los niveles L1-L2 a L5-S1. Se ajustó un YOLOv8n de detección de cajas (\arch{yolov8n.pt}) durante 50 épocas, con paciencia 30, lote 16, imagen de 640~px y semilla 42; el aumento en línea de Ultralytics (mosaico, escala, traslación, brillo) es el predeterminado, con reflexión horizontal. La mejor época en validación fue la 49. El detector solo ubica el nivel del disco; no clasifica patología. La Tabla~\ref{tab:detector} resume su desempeño con el umbral de uso (confianza 0,35).

\begin{table}[t]
\centering\small
\caption{Desempeño del detector de discos (corrida \arch{detector_discos_v1_s42}).}
\label{tab:detector}
\begin{tabular}{lrrrrrrr}
\toprule
Partición & Cortes & Precisión & Recall & mAP50 & mAP50-95 & Cortes con 5 niveles & IoU medio \\
\midrule
Validación & 20 & 0,922 & 0,964 & 0,962 & 0,663 & 95\,\% & 0,819 \\
Prueba     & 10 & 0,983 & 0,994 & 0,995 & 0,691 & 100\,\% & 0,843 \\
\bottomrule
\end{tabular}
\end{table}

El nivel más difícil es L1-L2 (en validación: precisión 0,839, recall 0,820 y IoU 0,74); los demás tienen recall de 1,0. Con solo 20 y 10 cortes, un corte mueve varios puntos porcentuales; es una sola corrida con una sola semilla, y la época se eligió con la validación, por lo que sus cifras están algo infladas.

\paragraph{Extracción del ROI.} Con \arch{autocrop_rois.py} se detecta sobre el recorte y se extrae el parche sobre el PNG completo, sumando el desplazamiento. La caja se amplía (margen posterior 0,40, anterior 0,10 y vertical 0,50), se hace cuadrada y se guarda a 128$\times$128~px, una por nivel (L1-L2, L2-L3, L3-L4, L4-L5 y L5-S1). Cada recorte queda registrado en \arch{roi_log.csv} con su caja, su ordenada central, su confianza y el relleno por lado. El resultado fueron \textbf{5\,385 ROI}. El margen anterior de L5-S1 se probó en 0,20 frente a 0,10 y la ganancia fue marginal, por lo que se mantuvo 0,10. La Figura~\ref{fig:rois} muestra los cinco ROI de un corte.

\begin{figure}[t]
\centering
\includegraphics[width=0.9\linewidth]{guia_visual_5rois.png}
\caption{Los cinco ROI (L1-L2 a L5-S1) de un corte sagital T2. \pendiente{confirmar que es la figura que se quiere en el informe o sustituirla por un montaje propio con las cajas del detector}}
\label{fig:rois}
\end{figure}

\paragraph{Exclusiones por calidad.} Se aplicaron reglas objetivas sobre \arch{roi_log.csv} (función \arch{elegibles()} de \arch{scripts/build_dataset.py}): se descarta el corte completo si falta algún nivel, si dos niveles adyacentes se solapan (IoU $>0{,}5$) o si el orden vertical de los centros no crece de L1-L2 a L5-S1; además se excluyen el paciente 0456, el corte 0547\_008 y los ROI con confianza del detector menor que 0,5. En total se excluyeron \textbf{262 ROI}, que afectan a 48 pacientes; 15 pacientes quedaron fuera por completo. Así, los 5\,385 ROI etiquetados quedan en \textbf{5\,123 ROI de 399 pacientes}; los 15 pacientes restantes (de los 414) no conservaron ningún ROI elegible. De los ROI excluidos, 207 eran Protrusión, 38 Normal y 17 Degeneración.

\subsection{Anotación y retiquetado}
\paragraph{Clases.} La guía de anotación (\arch{docs/guias/guia_anotacion_4_clases.md}) define cuatro clases nominales: Normal, Degeneración, Abombamiento y Hernia. Cuando coinciden varias, se aplica la prioridad Hernia $>$ Abombamiento $>$ Degeneración $>$ Normal, que solo sirve para desempatar el etiquetado y no es un orden de gravedad. La guía recoge criterios de uso común y no está validada por un radiólogo. Las etiquetas las asignaron estudiantes, sin radiólogo.

\paragraph{Fusión a tres clases.} En la primera corrida del clasificador semilla, con las cuatro clases (60 imágenes de validación) la exactitud fue 0,65 (IC 95\,\%: 0,52--0,76) y 10 de los 20 errores fueron confusiones Hernia--Abombamiento. Sin un corte axial, la extensión del abultamiento sobre la circunferencia del disco, que es lo que separa ambas entidades en la nomenclatura de consenso \cite{fardon2014}, no puede medirse. Por eso Hernia y Abombamiento se fusionaron en \textbf{Protrusión}, quedando tres clases: Normal, Degeneración discal y Protrusión. Otros autores también reportan confusión entre abultamiento y protrusión \cite{gao2025}. La fusión es una decisión posterior a ver los datos (\emph{post hoc}) y se declara así.

\paragraph{Preetiquetado y revisión.} Etiquetar 5\,385 ROI a mano desde cero era inviable en el plazo, así que se siguió un esquema de aprendizaje activo. Primero se etiquetó a mano una semilla en Roboflow (proyecto \emph{deepspine-seed-350-cls}); su versión~3, la usada para el preetiquetado final, tiene 719 imágenes de entrenamiento y 84 de validación (24 Degeneración, 30 Normal y 30 Protrusión). Con ella se entrenó un clasificador pequeño, YOLOv8s-cls (entrada 224$\times$224, hasta 60 épocas, paciencia 15, lote 32, semilla 42; cuaderno \arch{entrenamiento_clasificador_semilla_v3.ipynb}). En validación alcanzó una exactitud de 0,833 (IC 95\,\%: 0,739--0,898) y un F1 de 0,784 en Degeneración, 0,852 en Normal y 0,857 en Protrusión; con solo 84 imágenes de validación estas cifras son una guía y no una garantía. El modelo preetiquetó los 5\,385 ROI con su clase y su confianza (media 0,946; mediana 1,0) y marcó 2\,265 (42,1\,\%) para revisión prioritaria: 1\,881 por predecir una clase rara (Protrusión) y el resto por confianza menor que 0,7, por margen menor que 0,2 entre las dos primeras clases, o por combinación de motivos.

Después, todos los ROI pasaron por revisión humana en dos grupos, A y B, de 2\,693 y 2\,692 ROI respectivamente. La clase final es la decisión humana. De los 5\,385 ROI, 3\,366 (62,5\,\%) conservaron la clase del modelo y \textbf{2\,019 (37,5\,\%) fueron corregidos}. La Tabla~\ref{tab:reetiquetado} muestra cómo se reorganizó el dataset.

\begin{table}[t]
\centering\small
\caption{Clase propuesta por el preetiquetador frente a la clase final tras la revisión humana (5\,385 ROI).}
\label{tab:reetiquetado}
\begin{tabular}{lrrrrr}
\toprule
 & \multicolumn{3}{c}{Clase final} & & \\
\cmidrule(lr){2-4}
Predicha & Normal & Degeneración & Protrusión & Total & Corregidas \\
\midrule
Normal       & 1\,086 & 32  & 1\,387 & 2\,505 & 56,6\,\% \\
Degeneración & 51     & 256 & 534    & 841    & 69,6\,\% \\
Protrusión   & 6      & 9   & 2\,024 & 2\,039 & 0,7\,\% \\
\midrule
Total final  & 1\,143 & 297 & 3\,945 & 5\,385 & 37,5\,\% \\
\bottomrule
\end{tabular}
\end{table}

La lectura de la Tabla~\ref{tab:reetiquetado} es clara: el preetiquetador, entrenado con una semilla balanceada, \emph{subestimó} la Protrusión. De los 3\,945 ROI de Protrusión finales, 1\,921 (48,7\,\%) llegaron como Normal o Degeneración, mientras que casi nadie de lo que predijo como Protrusión cambió (15 de 2\,039). Dicho de otro modo, la revisión humana fue lo que fijó la distribución de clases del dataset, no el modelo, y la distribución real está muy desbalanceada (Sección~\ref{sec:balance}).

\paragraph{Fiabilidad de las etiquetas.} Cada ROI fue revisado por una sola persona (A o B); 408 de los 414 pacientes tienen ROI revisados por ambos, pero no hay ROI doble. En consecuencia \textbf{no puede calcularse el coeficiente kappa entre anotadores}. Sí se observan diferencias de criterio: A etiquetó Degeneración en el 6,9\,\% de sus ROI y B en el 4,1\,\%, y corrigió el 35,3\,\% frente al 39,7\,\% de B. No es posible separar cuánto de esa diferencia es criterio y cuánto es composición de la muestra.

\subsection{División por paciente}
\label{sec:split}
Cada paciente aporta unos 13 ROI (5\,123/399) de cortes vecinos casi idénticos. Si se dividiera por imagen, un paciente quedaría repartido entre entrenamiento y prueba y el modelo podría reconocer al paciente en lugar de la patología, inflando las métricas (fuga de datos) \cite{brunekreef2024,rumala2023}. Por eso \textbf{cada paciente va entero a una sola partición}.

El reparto es 70/15/15 por pacientes, con semilla 42. El script prueba 20\,000 repartos aleatorios y conserva el que mejor mantiene, en las tres particiones, la proporción de cada clase (con peso triple) y de cada nivel. Se eligieron tres particiones porque la validación sirve para decidir (época, hiperparámetros, umbrales) y la prueba para medir una sola vez al final, con todo ya decidido. La Tabla~\ref{tab:split} muestra el resultado.

\begin{table}[t]
\centering\small
\caption{División por paciente y ROI originales por clase.}
\label{tab:split}
\begin{tabular}{lrrrrrr}
\toprule
Partición & Pacientes & Normal & Degeneración & Protrusión & ROI & \% ROI \\
\midrule
Entrenamiento & 279 & 771   & 194 & 2\,612 & 3\,577 & 69,8 \\
Validación    & 60  & 170   & 44  & 564    & 778    & 15,2 \\
Prueba        & 60  & 164   & 42  & 562    & 768    & 15,0 \\
\midrule
Total         & 399 & 1\,105 & 280 & 3\,738 & 5\,123 & 100 \\
\bottomrule
\end{tabular}
\end{table}

La intersección de pacientes entre particiones es vacía y se comprueba con aserciones en el script. Los pacientes con al menos un ROI de Degeneración son 81, 17 y 22 en entrenamiento, validación y prueba. La docente aceptó esta división, y los 35\,007 registros que exige la rúbrica son las 35\,007 imágenes del dataset.

\subsection{Construcción del dataset de 35\,007 imágenes}
\label{sec:aumento}
La rúbrica exige un dataset balanceado con más de 35\,000 registros. Como los ROI originales son 5\,123, se recurrió a aumento de datos, siempre \emph{después} de dividir y \emph{dentro} de cada partición, para que ninguna copia cruce de partición. La cantidad de copias por ROI se calcula para que cada clase alcance la misma meta por partición: 8\,167 imágenes por clase en entrenamiento y 1\,751 en validación y en prueba, es decir 35\,007 en total.

\paragraph{Transformaciones.} Rotación de $\pm 8^{\circ}$, desplazamiento de hasta $\pm 8\,\%$ del tamaño, zoom entre 0,90 y 1,10, y brillo y contraste de $\pm 10\,\%$, con reflejo en los bordes. No se usan reflexiones (el eje anterior--posterior tiene significado anatómico), deformación elástica ni realce. Cada copia se genera con una semilla derivada de un hash (SHA-256) del nombre del ROI, la semilla global y el índice de la copia, por lo que el proceso es determinista y reproducible.

\begin{table}[t]
\centering\small
\caption{Imágenes finales por partición y clase (mismo número por clase en cada partición), ROI originales y copias por original.}
\label{tab:35k}
\begin{tabular}{llrrrr}
\toprule
Partición & Clase & Originales & Aumentadas & Total & Copias por original \\
\midrule
Entrenamiento & Normal       & 771    & 7\,396 & 8\,167 & 10,6 \\
              & Degeneración & 194    & 7\,973 & 8\,167 & 42,1 \\
              & Protrusión   & 2\,612 & 5\,555 & 8\,167 & 3,1  \\
Validación    & Normal       & 170    & 1\,581 & 1\,751 & 10,3 \\
              & Degeneración & 44     & 1\,707 & 1\,751 & 39,8 \\
              & Protrusión   & 564    & 1\,187 & 1\,751 & 3,1  \\
Prueba        & Normal       & 164    & 1\,587 & 1\,751 & 10,7 \\
              & Degeneración & 42     & 1\,709 & 1\,751 & 41,7 \\
              & Protrusión   & 562    & 1\,189 & 1\,751 & 3,1  \\
\midrule
Total         &              & 5\,123 & 29\,884 & 35\,007 & \\
\bottomrule
\end{tabular}
\end{table}

\begin{figure}[t]
\centering
\begin{tikzpicture}
\begin{axis}[
  ybar, bar width=14pt, width=0.85\linewidth, height=5.6cm,
  symbolic x coords={Normal,Degeneración,Protrusión}, xtick=data,
  ymin=0, ymax=9800, ylabel={Imágenes (entrenamiento)},
  enlarge x limits=0.25, legend style={at={(0.5,1.04)},anchor=south,legend columns=2,font=\footnotesize},
  nodes near coords, every node near coord/.append style={font=\tiny},
  tick label style={font=\footnotesize}, ylabel style={font=\footnotesize}]
\addplot coordinates {(Normal,771) (Degeneración,194) (Protrusión,2612)};
\addplot coordinates {(Normal,8167) (Degeneración,8167) (Protrusión,8167)};
\legend{Originales, Tras el aumento}
\end{axis}
\end{tikzpicture}
\caption{Balance por conteo en entrenamiento: ROI originales y total tras el aumento.}
\label{fig:balance}
\end{figure}

\paragraph{Trazabilidad.} Cada archivo se llama \arch{T2_TSE_SAG__0188_008_L4-L5__orig.png} (original) o \arch{..._aug012.png} (copia), con paciente, corte y nivel. El manifiesto \arch{manifest_35k.csv} registra, para cada una de las 35\,007 imágenes, su partición, clase, paciente, nivel, ROI original, revisor y, en las copias, los parámetros exactos (semilla, ángulo, zoom, desplazamientos, contraste y brillo); \arch{split_pacientes.csv} guarda la partición de cada paciente. Las copias respetan los rangos declarados (ángulo entre $-8^{\circ}$ y $8^{\circ}$, zoom entre 0,90 y 1,10, contraste entre 0,90 y 1,10).

\paragraph{Verificaciones.} El script comprueba que (i) los conjuntos de pacientes de las tres particiones son disjuntos, (ii) toda copia está en la misma partición que su original, (iii) los nombres de archivo son únicos, (iv) todo PNG está íntegro y (v) el total supera los 35\,000.

\subsection{Balance de clases y sesgo por nivel}
\label{sec:balance}
El balance de 1/3 por clase es \emph{por conteo}, no por originales. Entre los 5\,123 ROI originales, la Protrusión es el 73,0\,\% (3\,738), Normal el 21,6\,\% (1\,105) y Degeneración el 5,5\,\% (280); un clasificador que siempre dijera Protrusión acertaría cerca del 73\,\%. Las copias igualan la contribución de cada clase al entrenamiento, pero no añaden anatomía nueva: el tamaño efectivo del dataset es el número de pacientes (399), y la Degeneración, repetida unas 42 veces, solo tiene 42 originales en la partición de prueba, lo que deja un margen de error amplio en su recall. Para que las métricas no se inflen con copias casi idénticas, los resultados principales se calculan sobre los \textbf{originales} de la partición de prueba, y el conjunto completo se reporta aparte. Además, un modelo entrenado con prior uniforme tiende a sobrepredecir las clases raras en la población real.

\begin{table}[t]
\centering\small
\caption{ROI originales elegibles por nivel y clase.}
\label{tab:nivel}
\begin{tabular}{lrrrrr}
\toprule
Nivel & Normal & Degeneración & Protrusión & Total & \% Protrusión \\
\midrule
L1-L2 & 505 & 121 & 401    & 1\,027 & 39,0 \\
L2-L3 & 383 & 94  & 550    & 1\,027 & 53,6 \\
L3-L4 & 172 & 44  & 811    & 1\,027 & 79,0 \\
L4-L5 & 15  & 5   & 1\,005 & 1\,025 & 98,0 \\
L5-S1 & 30  & 16  & 971    & 1\,017 & 95,5 \\
\bottomrule
\end{tabular}
\end{table}

La Tabla~\ref{tab:nivel} muestra un sesgo importante: casi toda la patología se concentra en L4-L5 y L5-S1. Un clasificador que solo adivinara la clase más frecuente de cada nivel acertaría 3\,842 de 5\,123 ROI (75,0\,\%), de modo que el nivel podría actuar como atajo. Por ello los resultados se reportan también por nivel. Esto relativiza cualquier umbral de exactitud: un 95\,\% medido sobre estas etiquetas mide acuerdo con etiquetas de estudiantes y puede alcanzarse en parte por el atajo, y cifras internas altas suelen bajar al evaluar con datos externos \cite{yi2023}. Por eso se acompaña la exactitud con F1 macro y recall por clase, que además es lo que justifica el uso de F1 con clases desbalanceadas \cite{alkubaisi2022}.

\subsection{Viabilidad del dataset}
\begin{itemize}
  \item \textbf{Disponibilidad y licencia.} Los datos son públicos y los scripts reproducen el procesamiento a partir de ellos. Las imágenes del dataset construido no se suben al repositorio: se distribuyen como un archivo comprimido de 405~MB y el repositorio guarda sus manifiestos.
  \item \textbf{Tamaño.} Cumple el requisito de más de 35\,000 registros. El contenido informativo real es menor: 5\,123 ROI de 399 pacientes, con 280 originales de Degeneración.
  \item \textbf{Balance.} Equilibrado por conteo entre las tres clases y en cada partición; desbalanceado en originales y por nivel (Secciones~\ref{sec:aumento} y~\ref{sec:balance}).
  \item \textbf{Reproducibilidad.} Semilla única (42), manifiesto por imagen y verificaciones automáticas.
  \item \textbf{Cómputo.} Los cuatro modelos se entrenaron con este mismo dataset en una GPU de consumo (RTX 3060~Ti), lo que indica que el tamaño es manejable con recursos modestos.
\end{itemize}

\subsection{Limitaciones}
(1) Las etiquetas son de estudiantes, sin radiólogo, con apoyo de un modelo, y no se midió el acuerdo entre anotadores. (2) Solo hay plano sagital; Hernia y Abombamiento no son separables y se fusionaron después de ver los datos. (3) Un único equipo de resonancia de 1,5~T, sin diversidad de escáner. (4) El tamaño efectivo son 399 pacientes; las copias no aportan anatomía nueva. (5) Degeneración cuenta con pocos pacientes (81, 17 y 22 en las tres particiones). (6) Sesgo por nivel y prior uniforme (Sección~\ref{sec:balance}). (7) Con el identificador de paciente vacío no puede descartarse la presencia de un mismo paciente en dos carpetas. (8) El conjunto es una herramienta académica, no diagnóstica.

% Paquetes requeridos: graphicx, booktabs, xcolor, tikz (+ calc no), pgfplots, caption
% Figuras: guia_visual_5rois.png vive en docs/guias/ (ver \graphicspath en el archivo principal)
% Las cifras salen de data/processed/etiquetado/etiquetas_revisadas.csv,
% data/processed/dataset_35k/manifest_35k.csv y de los scripts de scripts/.
% Todo lo que no se pudo verificar se marca con \pendiente{...} (se ve en rojo).
% =====================================================================
\providecommand{\pendiente}[1]{\textcolor{red}{\textbf{[PENDIENTE: #1]}}}
\providecommand{\arch}[1]{\texttt{\detokenize{#1}}}

\section{Preparación del dataset}
\label{sec:dataset}

Esta sección describe, en orden, cómo se pasó del conjunto público original a los 35\,007 recortes (ROI) con los que se entrenaron y compararon los cuatro modelos. La Figura~\ref{fig:flujo} resume el flujo y el número de elementos que pasan por cada etapa. Todas las cifras proceden de los archivos versionados del repositorio; lo que no pudo verificarse se señala como pendiente.

\begin{figure}[t]
\centering
\begin{tikzpicture}[
  box/.style={draw, rounded corners=2pt, align=left, text width=0.60\linewidth, inner sep=3pt, font=\footnotesize},
  num/.style={align=left, text width=0.30\linewidth, font=\footnotesize\bfseries},
  arrow/.style={-{Stealth[length=2mm]}, thick},
  node distance=4.2mm]
  \node[box] (a) {1. Conjunto público (Mendeley), archivos \texttt{.IMA}};
  \node[num, right=3mm of a] {575 carpetas, $\approx$48\,000 archivos};
  \node[box, below=of a] (b) {2. Descarte de series inútiles y conservación de T2 sagital};
  \node[num, right=3mm of b] {solo T2 sagital};
  \node[box, below=of b] (c) {3. Conversión DICOM$\to$PNG y limpieza manual de cortes};
  \node[num, right=3mm of c] {1\,086 cortes, 414 pacientes};
  \node[box, below=of c] (d) {4. Recorte L1--S1 y detección de discos (ROI 128$\times$128)};
  \node[num, right=3mm of d] {5\,385 ROI};
  \node[box, below=of d] (e) {5. Preetiquetado, revisión humana (A/B) y exclusiones};
  \node[num, right=3mm of e] {5\,123 ROI, 399 pacientes};
  \node[box, below=of e] (f) {6. División por paciente 70/15/15};
  \node[num, right=3mm of f] {279 / 60 / 60 pacientes};
  \node[box, below=of f] (g) {7. Aumento dentro de cada partición (balance por conteo)};
  \node[num, right=3mm of g] {35\,007 imágenes};
  \foreach \i/\j in {a/b,b/c,c/d,d/e,e/f,f/g} \draw[arrow] (\i) -- (\j);
\end{tikzpicture}
\caption{Flujo de preparación del dataset y número de elementos en cada etapa.}
\label{fig:flujo}
\end{figure}

\subsection{Origen de los datos}
Se usó el \emph{Lumbar Spine MRI Dataset} (Mendeley Data) \cite{sudirman2019}, el mismo conjunto público empleado en trabajos previos sobre clasificación de hernia discal \cite{sari2024}. Los estudios se adquirieron con un único equipo, un Siemens MAGNETOM Essenza de 1,5~T, y se distribuyen como archivos DICOM con extensión \arch{.IMA}, en una carpeta por paciente. Los cortes tienen un tamaño de píxel de 0,729~mm, un grosor de 4~mm y una separación entre centros de 4,8~mm. El campo \texttt{PatientID} está vacío en todo el corpus y el repositorio solo conserva identificadores de estudio y de serie (los campos de nombre, nacimiento y fecha están vacíos en \arch{meta_dicom_por_slice.csv}). El conjunto descargado contenía 575 carpetas de pacientes. Se eliminaron las que, en lugar de resonancia lumbar, contenían estudios del segmento cervical o de rodillas, y las lumbares con aberraciones o manchas que habrían introducido ruido en los modelos. Tras esa depuración quedaron 414 pacientes con cortes útiles y, tras las exclusiones de calidad de los ROI (Sección~\ref{sec:roi}), los \textbf{399 pacientes} del dataset final. Cada una de las 414 carpetas restantes corresponde a un \texttt{StudyInstanceUID} distinto y no se repite ningún \texttt{StudyTime}; no hay evidencia de pacientes duplicados, aunque con el identificador vacío no puede descartarse del todo.
% VERIFICAR: el sitio de Mendeley describe 515 pacientes; el equipo reporta 575 carpetas descargadas.

\subsection{Depuración de la fuente}
La descarga inicial contenía cerca de 48\,000 registros en formato \arch{.IMA}. Cada paciente trae varias series; solo una parte sirve para este problema. La Tabla~\ref{tab:descartes} lista lo que se descartó y por qué.

\begin{table}[t]
\centering\small
\caption{Series descartadas de cada estudio y motivo.}
\label{tab:descartes}
\begin{tabular}{p{0.30\linewidth}p{0.60\linewidth}}
\toprule
Serie & Motivo del descarte \\
\midrule
\arch{LOCALIZER_0001} & Localizadores de muy baja resolución usados por el técnico para calibrar la mesa y posicionar al paciente; no aportan información anatómica ni de textura. \\
\arch{POSDISP_[4]_T2_TSE_TRA_...} & Vistas sagitales con líneas blancas superpuestas (\emph{overlays}) que marcan los planos transversales; ese artefacto distorsiona descriptores de textura (GLCM/LBP) y confunde a las redes. \\
\arch{T1_TSE_SAG_...}, \arch{T1_TSE_TRA_...} & En T1 la grasa es brillante y el líquido oscuro, de modo que no permite evaluar la hidratación del disco (base de la escala de Pfirrmann); además, mezclar T1 y T2 altera los rangos de brillo. \\
\arch{T2_TSE_TRA_...} (axial) & Cortes axiales fragmentados en lotes de 3 a 5 por nivel; exigirían un emparejamiento espacial adicional por disco, mientras que el plano sagital cubre L1--S1 en una sola imagen. Se deja para una fase futura. \\
\bottomrule
\end{tabular}
\end{table}

El proceso se hizo en cuatro pasos, cada uno con su script (carpeta \arch{scripts/preprocesamiento/}):
\begin{enumerate}
  \item \arch{ingestion.py} recorre las carpetas de pacientes y copia solo las imágenes T2 sagitales, identificadas por el nombre de la serie.
  \item \arch{limpiar_slices.py} conserva el 40\,\% central del volumen sagital (se descarta el 30\,\% inicial y el 30\,\% final), donde el canal medular y la región L1--S1 son visibles.
  \item \arch{convertidor_dcm_a_png.py} convierte a PNG de 8 bits con una sola escala por serie: percentiles 1 y 99,5 calculados sobre los píxeles de tejido. Se decidió no aplicar CLAHE ni otro realce, para no alterar la señal relativa del T2 que sustenta el criterio diagnóstico.
  \item Depuración manual de los PNG: algunos pacientes tenían estudios del segmento cervical o de rodillas, y otros cortes T2 sagitales quedaban oscuros o con aberraciones; se eliminaron.
\end{enumerate}
Tras esta depuración quedaron \textbf{1\,086 cortes sagitales de 414 pacientes} (2,6 cortes por paciente en promedio; entre 1 y 6).

\subsection{Recorte de la región lumbar y de los ROI por disco}
\label{sec:roi}
\paragraph{Recorte L1--S1.} Con \arch{aplicar_recortes.py} se recortó cada corte en altura para abarcar principalmente L1--S1, replicando por \emph{template matching} unos recortes manuales previos (mediana de la puntuación de coincidencia 0,9996; solo el corte 0547\_008 quedó por debajo de 0,99). El recorte se extendió 60~px hacia arriba y hacia abajo hasta el borde de la imagen; con ello el relleno necesario en L1-L2 bajó de 779 de 1\,080 recortes (72\,\%) a cero.

\paragraph{Detector de discos.} Para localizar los discos se entrenó un detector con un conjunto propio en Roboflow (proyecto \emph{Intervertebral Disks}, versión~1): 100 cortes (70 de entrenamiento, 20 de validación y 10 de prueba), con orientación automática (\emph{auto-orient}) y reescalado a 640$\times$640 con bordes negros, sin aumento de datos en Roboflow. Las cinco clases del detector son los niveles L1-L2 a L5-S1. Se ajustó un YOLOv8n de detección de cajas (\arch{yolov8n.pt}) durante 50 épocas, con paciencia 30, lote 16, imagen de 640~px y semilla 42; el aumento en línea de Ultralytics (mosaico, escala, traslación, brillo) es el predeterminado, con reflexión horizontal. La mejor época en validación fue la 49. El detector solo ubica el nivel del disco; no clasifica patología. La Tabla~\ref{tab:detector} resume su desempeño con el umbral de uso (confianza 0,35).

\begin{table}[t]
\centering\small
\caption{Desempeño del detector de discos (corrida \arch{detector_discos_v1_s42}).}
\label{tab:detector}
\begin{tabular}{lrrrrrrr}
\toprule
Partición & Cortes & Precisión & Recall & mAP50 & mAP50-95 & Cortes con 5 niveles & IoU medio \\
\midrule
Validación & 20 & 0,922 & 0,964 & 0,962 & 0,663 & 95\,\% & 0,819 \\
Prueba     & 10 & 0,983 & 0,994 & 0,995 & 0,691 & 100\,\% & 0,843 \\
\bottomrule
\end{tabular}
\end{table}

El nivel más difícil es L1-L2 (en validación: precisión 0,839, recall 0,820 y IoU 0,74); los demás tienen recall de 1,0. Con solo 20 y 10 cortes, un corte mueve varios puntos porcentuales; es una sola corrida con una sola semilla, y la época se eligió con la validación, por lo que sus cifras están algo infladas.

\paragraph{Extracción del ROI.} Con \arch{autocrop_rois.py} se detecta sobre el recorte y se extrae el parche sobre el PNG completo, sumando el desplazamiento. La caja se amplía (margen posterior 0,40, anterior 0,10 y vertical 0,50), se hace cuadrada y se guarda a 128$\times$128~px, una por nivel (L1-L2, L2-L3, L3-L4, L4-L5 y L5-S1). Cada recorte queda registrado en \arch{roi_log.csv} con su caja, su ordenada central, su confianza y el relleno por lado. El resultado fueron \textbf{5\,385 ROI}. El margen anterior de L5-S1 se probó en 0,20 frente a 0,10 y la ganancia fue marginal, por lo que se mantuvo 0,10. La Figura~\ref{fig:rois} muestra los cinco ROI de un corte.

\begin{figure}[t]
\centering
\includegraphics[width=0.9\linewidth]{guia_visual_5rois.png}
\caption{Los cinco ROI (L1-L2 a L5-S1) de un corte sagital T2. \pendiente{confirmar que es la figura que se quiere en el informe o sustituirla por un montaje propio con las cajas del detector}}
\label{fig:rois}
\end{figure}

\paragraph{Exclusiones por calidad.} Se aplicaron reglas objetivas sobre \arch{roi_log.csv} (función \arch{elegibles()} de \arch{scripts/build_dataset.py}): se descarta el corte completo si falta algún nivel, si dos niveles adyacentes se solapan (IoU $>0{,}5$) o si el orden vertical de los centros no crece de L1-L2 a L5-S1; además se excluyen el paciente 0456, el corte 0547\_008 y los ROI con confianza del detector menor que 0,5. En total se excluyeron \textbf{262 ROI}, que afectan a 48 pacientes; 15 pacientes quedaron fuera por completo. Así, los 5\,385 ROI etiquetados quedan en \textbf{5\,123 ROI de 399 pacientes}; los 15 pacientes restantes (de los 414) no conservaron ningún ROI elegible. De los ROI excluidos, 207 eran Protrusión, 38 Normal y 17 Degeneración.

\subsection{Anotación y retiquetado}
\paragraph{Clases.} La guía de anotación (\arch{docs/guias/guia_anotacion_4_clases.md}) define cuatro clases nominales: Normal, Degeneración, Abombamiento y Hernia. Cuando coinciden varias, se aplica la prioridad Hernia $>$ Abombamiento $>$ Degeneración $>$ Normal, que solo sirve para desempatar el etiquetado y no es un orden de gravedad. La guía recoge criterios de uso común y no está validada por un radiólogo. Las etiquetas las asignaron estudiantes, sin radiólogo.

\paragraph{Fusión a tres clases.} En la primera corrida del clasificador semilla, con las cuatro clases (60 imágenes de validación) la exactitud fue 0,65 (IC 95\,\%: 0,52--0,76) y 10 de los 20 errores fueron confusiones Hernia--Abombamiento. Sin un corte axial, la extensión del abultamiento sobre la circunferencia del disco, que es lo que separa ambas entidades en la nomenclatura de consenso \cite{fardon2014}, no puede medirse. Por eso Hernia y Abombamiento se fusionaron en \textbf{Protrusión}, quedando tres clases: Normal, Degeneración discal y Protrusión. Otros autores también reportan confusión entre abultamiento y protrusión \cite{gao2025}. La fusión es una decisión posterior a ver los datos (\emph{post hoc}) y se declara así.

\paragraph{Preetiquetado y revisión.} Etiquetar 5\,385 ROI a mano desde cero era inviable en el plazo, así que se siguió un esquema de aprendizaje activo. Primero se etiquetó a mano una semilla en Roboflow (proyecto \emph{deepspine-seed-350-cls}); su versión~3, la usada para el preetiquetado final, tiene 719 imágenes de entrenamiento y 84 de validación (24 Degeneración, 30 Normal y 30 Protrusión). Con ella se entrenó un clasificador pequeño, YOLOv8s-cls (entrada 224$\times$224, hasta 60 épocas, paciencia 15, lote 32, semilla 42; cuaderno \arch{entrenamiento_clasificador_semilla_v3.ipynb}). En validación alcanzó una exactitud de 0,833 (IC 95\,\%: 0,739--0,898) y un F1 de 0,784 en Degeneración, 0,852 en Normal y 0,857 en Protrusión; con solo 84 imágenes de validación estas cifras son una guía y no una garantía. El modelo preetiquetó los 5\,385 ROI con su clase y su confianza (media 0,946; mediana 1,0) y marcó 2\,265 (42,1\,\%) para revisión prioritaria: 1\,881 por predecir una clase rara (Protrusión) y el resto por confianza menor que 0,7, por margen menor que 0,2 entre las dos primeras clases, o por combinación de motivos.

Después, todos los ROI pasaron por revisión humana en dos grupos, A y B, de 2\,693 y 2\,692 ROI respectivamente. La clase final es la decisión humana. De los 5\,385 ROI, 3\,366 (62,5\,\%) conservaron la clase del modelo y \textbf{2\,019 (37,5\,\%) fueron corregidos}. La Tabla~\ref{tab:reetiquetado} muestra cómo se reorganizó el dataset.

\begin{table}[t]
\centering\small
\caption{Clase propuesta por el preetiquetador frente a la clase final tras la revisión humana (5\,385 ROI).}
\label{tab:reetiquetado}
\begin{tabular}{lrrrrr}
\toprule
 & \multicolumn{3}{c}{Clase final} & & \\
\cmidrule(lr){2-4}
Predicha & Normal & Degeneración & Protrusión & Total & Corregidas \\
\midrule
Normal       & 1\,086 & 32  & 1\,387 & 2\,505 & 56,6\,\% \\
Degeneración & 51     & 256 & 534    & 841    & 69,6\,\% \\
Protrusión   & 6      & 9   & 2\,024 & 2\,039 & 0,7\,\% \\
\midrule
Total final  & 1\,143 & 297 & 3\,945 & 5\,385 & 37,5\,\% \\
\bottomrule
\end{tabular}
\end{table}

La lectura de la Tabla~\ref{tab:reetiquetado} es clara: el preetiquetador, entrenado con una semilla balanceada, \emph{subestimó} la Protrusión. De los 3\,945 ROI de Protrusión finales, 1\,921 (48,7\,\%) llegaron como Normal o Degeneración, mientras que casi nadie de lo que predijo como Protrusión cambió (15 de 2\,039). Dicho de otro modo, la revisión humana fue lo que fijó la distribución de clases del dataset, no el modelo, y la distribución real está muy desbalanceada (Sección~\ref{sec:balance}).

\paragraph{Fiabilidad de las etiquetas.} Cada ROI fue revisado por una sola persona (A o B); 408 de los 414 pacientes tienen ROI revisados por ambos, pero no hay ROI doble. En consecuencia \textbf{no puede calcularse el coeficiente kappa entre anotadores}. Sí se observan diferencias de criterio: A etiquetó Degeneración en el 6,9\,\% de sus ROI y B en el 4,1\,\%, y corrigió el 35,3\,\% frente al 39,7\,\% de B. No es posible separar cuánto de esa diferencia es criterio y cuánto es composición de la muestra.

\subsection{División por paciente}
\label{sec:split}
Cada paciente aporta unos 13 ROI (5\,123/399) de cortes vecinos casi idénticos. Si se dividiera por imagen, un paciente quedaría repartido entre entrenamiento y prueba y el modelo podría reconocer al paciente en lugar de la patología, inflando las métricas (fuga de datos) \cite{brunekreef2024,rumala2023}. Por eso \textbf{cada paciente va entero a una sola partición}.

El reparto es 70/15/15 por pacientes, con semilla 42. El script prueba 20\,000 repartos aleatorios y conserva el que mejor mantiene, en las tres particiones, la proporción de cada clase (con peso triple) y de cada nivel. Se eligieron tres particiones porque la validación sirve para decidir (época, hiperparámetros, umbrales) y la prueba para medir una sola vez al final, con todo ya decidido. La Tabla~\ref{tab:split} muestra el resultado.

\begin{table}[t]
\centering\small
\caption{División por paciente y ROI originales por clase.}
\label{tab:split}
\begin{tabular}{lrrrrrr}
\toprule
Partición & Pacientes & Normal & Degeneración & Protrusión & ROI & \% ROI \\
\midrule
Entrenamiento & 279 & 771   & 194 & 2\,612 & 3\,577 & 69,8 \\
Validación    & 60  & 170   & 44  & 564    & 778    & 15,2 \\
Prueba        & 60  & 164   & 42  & 562    & 768    & 15,0 \\
\midrule
Total         & 399 & 1\,105 & 280 & 3\,738 & 5\,123 & 100 \\
\bottomrule
\end{tabular}
\end{table}

La intersección de pacientes entre particiones es vacía y se comprueba con aserciones en el script. Los pacientes con al menos un ROI de Degeneración son 81, 17 y 22 en entrenamiento, validación y prueba. La docente aceptó esta división, y los 35\,007 registros que exige la rúbrica son las 35\,007 imágenes del dataset.

\subsection{Construcción del dataset de 35\,007 imágenes}
\label{sec:aumento}
La rúbrica exige un dataset balanceado con más de 35\,000 registros. Como los ROI originales son 5\,123, se recurrió a aumento de datos, siempre \emph{después} de dividir y \emph{dentro} de cada partición, para que ninguna copia cruce de partición. La cantidad de copias por ROI se calcula para que cada clase alcance la misma meta por partición: 8\,167 imágenes por clase en entrenamiento y 1\,751 en validación y en prueba, es decir 35\,007 en total.

\paragraph{Transformaciones.} Rotación de $\pm 8^{\circ}$, desplazamiento de hasta $\pm 8\,\%$ del tamaño, zoom entre 0,90 y 1,10, y brillo y contraste de $\pm 10\,\%$, con reflejo en los bordes. No se usan reflexiones (el eje anterior--posterior tiene significado anatómico), deformación elástica ni realce. Cada copia se genera con una semilla derivada de un hash (SHA-256) del nombre del ROI, la semilla global y el índice de la copia, por lo que el proceso es determinista y reproducible.

\begin{table}[t]
\centering\small
\caption{Imágenes finales por partición y clase (mismo número por clase en cada partición), ROI originales y copias por original.}
\label{tab:35k}
\begin{tabular}{llrrrr}
\toprule
Partición & Clase & Originales & Aumentadas & Total & Copias por original \\
\midrule
Entrenamiento & Normal       & 771    & 7\,396 & 8\,167 & 10,6 \\
              & Degeneración & 194    & 7\,973 & 8\,167 & 42,1 \\
              & Protrusión   & 2\,612 & 5\,555 & 8\,167 & 3,1  \\
Validación    & Normal       & 170    & 1\,581 & 1\,751 & 10,3 \\
              & Degeneración & 44     & 1\,707 & 1\,751 & 39,8 \\
              & Protrusión   & 564    & 1\,187 & 1\,751 & 3,1  \\
Prueba        & Normal       & 164    & 1\,587 & 1\,751 & 10,7 \\
              & Degeneración & 42     & 1\,709 & 1\,751 & 41,7 \\
              & Protrusión   & 562    & 1\,189 & 1\,751 & 3,1  \\
\midrule
Total         &              & 5\,123 & 29\,884 & 35\,007 & \\
\bottomrule
\end{tabular}
\end{table}

\begin{figure}[t]
\centering
\begin{tikzpicture}
\begin{axis}[
  ybar, bar width=14pt, width=0.85\linewidth, height=5.6cm,
  symbolic x coords={Normal,Degeneración,Protrusión}, xtick=data,
  ymin=0, ymax=9800, ylabel={Imágenes (entrenamiento)},
  enlarge x limits=0.25, legend style={at={(0.5,1.04)},anchor=south,legend columns=2,font=\footnotesize},
  nodes near coords, every node near coord/.append style={font=\tiny},
  tick label style={font=\footnotesize}, ylabel style={font=\footnotesize}]
\addplot coordinates {(Normal,771) (Degeneración,194) (Protrusión,2612)};
\addplot coordinates {(Normal,8167) (Degeneración,8167) (Protrusión,8167)};
\legend{Originales, Tras el aumento}
\end{axis}
\end{tikzpicture}
\caption{Balance por conteo en entrenamiento: ROI originales y total tras el aumento.}
\label{fig:balance}
\end{figure}

\paragraph{Trazabilidad.} Cada archivo se llama \arch{T2_TSE_SAG__0188_008_L4-L5__orig.png} (original) o \arch{..._aug012.png} (copia), con paciente, corte y nivel. El manifiesto \arch{manifest_35k.csv} registra, para cada una de las 35\,007 imágenes, su partición, clase, paciente, nivel, ROI original, revisor y, en las copias, los parámetros exactos (semilla, ángulo, zoom, desplazamientos, contraste y brillo); \arch{split_pacientes.csv} guarda la partición de cada paciente. Las copias respetan los rangos declarados (ángulo entre $-8^{\circ}$ y $8^{\circ}$, zoom entre 0,90 y 1,10, contraste entre 0,90 y 1,10).

\paragraph{Verificaciones.} El script comprueba que (i) los conjuntos de pacientes de las tres particiones son disjuntos, (ii) toda copia está en la misma partición que su original, (iii) los nombres de archivo son únicos, (iv) todo PNG está íntegro y (v) el total supera los 35\,000.

\subsection{Balance de clases y sesgo por nivel}
\label{sec:balance}
El balance de 1/3 por clase es \emph{por conteo}, no por originales. Entre los 5\,123 ROI originales, la Protrusión es el 73,0\,\% (3\,738), Normal el 21,6\,\% (1\,105) y Degeneración el 5,5\,\% (280); un clasificador que siempre dijera Protrusión acertaría cerca del 73\,\%. Las copias igualan la contribución de cada clase al entrenamiento, pero no añaden anatomía nueva: el tamaño efectivo del dataset es el número de pacientes (399), y la Degeneración, repetida unas 42 veces, solo tiene 42 originales en la partición de prueba, lo que deja un margen de error amplio en su recall. Para que las métricas no se inflen con copias casi idénticas, los resultados principales se calculan sobre los \textbf{originales} de la partición de prueba, y el conjunto completo se reporta aparte. Además, un modelo entrenado con prior uniforme tiende a sobrepredecir las clases raras en la población real.

\begin{table}[t]
\centering\small
\caption{ROI originales elegibles por nivel y clase.}
\label{tab:nivel}
\begin{tabular}{lrrrrr}
\toprule
Nivel & Normal & Degeneración & Protrusión & Total & \% Protrusión \\
\midrule
L1-L2 & 505 & 121 & 401    & 1\,027 & 39,0 \\
L2-L3 & 383 & 94  & 550    & 1\,027 & 53,6 \\
L3-L4 & 172 & 44  & 811    & 1\,027 & 79,0 \\
L4-L5 & 15  & 5   & 1\,005 & 1\,025 & 98,0 \\
L5-S1 & 30  & 16  & 971    & 1\,017 & 95,5 \\
\bottomrule
\end{tabular}
\end{table}

La Tabla~\ref{tab:nivel} muestra un sesgo importante: casi toda la patología se concentra en L4-L5 y L5-S1. Un clasificador que solo adivinara la clase más frecuente de cada nivel acertaría 3\,842 de 5\,123 ROI (75,0\,\%), de modo que el nivel podría actuar como atajo. Por ello los resultados se reportan también por nivel. Esto relativiza cualquier umbral de exactitud: un 95\,\% medido sobre estas etiquetas mide acuerdo con etiquetas de estudiantes y puede alcanzarse en parte por el atajo, y cifras internas altas suelen bajar al evaluar con datos externos \cite{yi2023}. Por eso se acompaña la exactitud con F1 macro y recall por clase, que además es lo que justifica el uso de F1 con clases desbalanceadas \cite{alkubaisi2022}.

\subsection{Viabilidad del dataset}
\begin{itemize}
  \item \textbf{Disponibilidad y licencia.} Los datos son públicos y los scripts reproducen el procesamiento a partir de ellos. Las imágenes del dataset construido no se suben al repositorio: se distribuyen como un archivo comprimido de 405~MB y el repositorio guarda sus manifiestos.
  \item \textbf{Tamaño.} Cumple el requisito de más de 35\,000 registros. El contenido informativo real es menor: 5\,123 ROI de 399 pacientes, con 280 originales de Degeneración.
  \item \textbf{Balance.} Equilibrado por conteo entre las tres clases y en cada partición; desbalanceado en originales y por nivel (Secciones~\ref{sec:aumento} y~\ref{sec:balance}).
  \item \textbf{Reproducibilidad.} Semilla única (42), manifiesto por imagen y verificaciones automáticas.
  \item \textbf{Cómputo.} Los cuatro modelos se entrenaron con este mismo dataset en una GPU de consumo (RTX 3060~Ti), lo que indica que el tamaño es manejable con recursos modestos.
\end{itemize}

\subsection{Limitaciones}
(1) Las etiquetas son de estudiantes, sin radiólogo, con apoyo de un modelo, y no se midió el acuerdo entre anotadores. (2) Solo hay plano sagital; Hernia y Abombamiento no son separables y se fusionaron después de ver los datos. (3) Un único equipo de resonancia de 1,5~T, sin diversidad de escáner. (4) El tamaño efectivo son 399 pacientes; las copias no aportan anatomía nueva. (5) Degeneración cuenta con pocos pacientes (81, 17 y 22 en las tres particiones). (6) Sesgo por nivel y prior uniforme (Sección~\ref{sec:balance}). (7) Con el identificador de paciente vacío no puede descartarse la presencia de un mismo paciente en dos carpetas. (8) El conjunto es una herramienta académica, no diagnóstica.

% Paquetes requeridos: graphicx, booktabs, xcolor, tikz (+ calc no), pgfplots, caption
% Figuras: guia_visual_5rois.png vive en docs/guias/ (ver \graphicspath en el archivo principal)
% Las cifras salen de data/processed/etiquetado/etiquetas_revisadas.csv,
% data/processed/dataset_35k/manifest_35k.csv y de los scripts de scripts/.
% Todo lo que no se pudo verificar se marca con \pendiente{...} (se ve en rojo).
% =====================================================================
\providecommand{\pendiente}[1]{\textcolor{red}{\textbf{[PENDIENTE: #1]}}}
\providecommand{\arch}[1]{\texttt{\detokenize{#1}}}

\section{Preparación del dataset}
\label{sec:dataset}

Esta sección describe, en orden, cómo se pasó del conjunto público original a los 35\,007 recortes (ROI) con los que se entrenaron y compararon los cuatro modelos. La Figura~\ref{fig:flujo} resume el flujo y el número de elementos que pasan por cada etapa. Todas las cifras proceden de los archivos versionados del repositorio; lo que no pudo verificarse se señala como pendiente.

\begin{figure}[t]
\centering
\begin{tikzpicture}[
  box/.style={draw, rounded corners=2pt, align=left, text width=0.60\linewidth, inner sep=3pt, font=\footnotesize},
  num/.style={align=left, text width=0.30\linewidth, font=\footnotesize\bfseries},
  arrow/.style={-{Stealth[length=2mm]}, thick},
  node distance=4.2mm]
  \node[box] (a) {1. Conjunto público (Mendeley), archivos \texttt{.IMA}};
  \node[num, right=3mm of a] {575 carpetas, $\approx$48\,000 archivos};
  \node[box, below=of a] (b) {2. Descarte de series inútiles y conservación de T2 sagital};
  \node[num, right=3mm of b] {solo T2 sagital};
  \node[box, below=of b] (c) {3. Conversión DICOM$\to$PNG y limpieza manual de cortes};
  \node[num, right=3mm of c] {1\,086 cortes, 414 pacientes};
  \node[box, below=of c] (d) {4. Recorte L1--S1 y detección de discos (ROI 128$\times$128)};
  \node[num, right=3mm of d] {5\,385 ROI};
  \node[box, below=of d] (e) {5. Preetiquetado, revisión humana (A/B) y exclusiones};
  \node[num, right=3mm of e] {5\,123 ROI, 399 pacientes};
  \node[box, below=of e] (f) {6. División por paciente 70/15/15};
  \node[num, right=3mm of f] {279 / 60 / 60 pacientes};
  \node[box, below=of f] (g) {7. Aumento dentro de cada partición (balance por conteo)};
  \node[num, right=3mm of g] {35\,007 imágenes};
  \foreach \i/\j in {a/b,b/c,c/d,d/e,e/f,f/g} \draw[arrow] (\i) -- (\j);
\end{tikzpicture}
\caption{Flujo de preparación del dataset y número de elementos en cada etapa.}
\label{fig:flujo}
\end{figure}

\subsection{Origen de los datos}
Se usó el \emph{Lumbar Spine MRI Dataset} (Mendeley Data) \cite{sudirman2019}, el mismo conjunto público empleado en trabajos previos sobre clasificación de hernia discal \cite{sari2024}. Los estudios se adquirieron con un único equipo, un Siemens MAGNETOM Essenza de 1,5~T, y se distribuyen como archivos DICOM con extensión \arch{.IMA}, en una carpeta por paciente. Los cortes tienen un tamaño de píxel de 0,729~mm, un grosor de 4~mm y una separación entre centros de 4,8~mm. El campo \texttt{PatientID} está vacío en todo el corpus y el repositorio solo conserva identificadores de estudio y de serie (los campos de nombre, nacimiento y fecha están vacíos en \arch{meta_dicom_por_slice.csv}). El conjunto descargado contenía 575 carpetas de pacientes. Se eliminaron las que, en lugar de resonancia lumbar, contenían estudios del segmento cervical o de rodillas, y las lumbares con aberraciones o manchas que habrían introducido ruido en los modelos. Tras esa depuración quedaron 414 pacientes con cortes útiles y, tras las exclusiones de calidad de los ROI (Sección~\ref{sec:roi}), los \textbf{399 pacientes} del dataset final. Cada una de las 414 carpetas restantes corresponde a un \texttt{StudyInstanceUID} distinto y no se repite ningún \texttt{StudyTime}; no hay evidencia de pacientes duplicados, aunque con el identificador vacío no puede descartarse del todo.
% VERIFICAR: el sitio de Mendeley describe 515 pacientes; el equipo reporta 575 carpetas descargadas.

\subsection{Depuración de la fuente}
La descarga inicial contenía cerca de 48\,000 registros en formato \arch{.IMA}. Cada paciente trae varias series; solo una parte sirve para este problema. La Tabla~\ref{tab:descartes} lista lo que se descartó y por qué.

\begin{table}[t]
\centering\small
\caption{Series descartadas de cada estudio y motivo.}
\label{tab:descartes}
\begin{tabular}{p{0.30\linewidth}p{0.60\linewidth}}
\toprule
Serie & Motivo del descarte \\
\midrule
\arch{LOCALIZER_0001} & Localizadores de muy baja resolución usados por el técnico para calibrar la mesa y posicionar al paciente; no aportan información anatómica ni de textura. \\
\arch{POSDISP_[4]_T2_TSE_TRA_...} & Vistas sagitales con líneas blancas superpuestas (\emph{overlays}) que marcan los planos transversales; ese artefacto distorsiona descriptores de textura (GLCM/LBP) y confunde a las redes. \\
\arch{T1_TSE_SAG_...}, \arch{T1_TSE_TRA_...} & En T1 la grasa es brillante y el líquido oscuro, de modo que no permite evaluar la hidratación del disco (base de la escala de Pfirrmann); además, mezclar T1 y T2 altera los rangos de brillo. \\
\arch{T2_TSE_TRA_...} (axial) & Cortes axiales fragmentados en lotes de 3 a 5 por nivel; exigirían un emparejamiento espacial adicional por disco, mientras que el plano sagital cubre L1--S1 en una sola imagen. Se deja para una fase futura. \\
\bottomrule
\end{tabular}
\end{table}

El proceso se hizo en cuatro pasos, cada uno con su script (carpeta \arch{scripts/preprocesamiento/}):
\begin{enumerate}
  \item \arch{ingestion.py} recorre las carpetas de pacientes y copia solo las imágenes T2 sagitales, identificadas por el nombre de la serie.
  \item \arch{limpiar_slices.py} conserva el 40\,\% central del volumen sagital (se descarta el 30\,\% inicial y el 30\,\% final), donde el canal medular y la región L1--S1 son visibles.
  \item \arch{convertidor_dcm_a_png.py} convierte a PNG de 8 bits con una sola escala por serie: percentiles 1 y 99,5 calculados sobre los píxeles de tejido. Se decidió no aplicar CLAHE ni otro realce, para no alterar la señal relativa del T2 que sustenta el criterio diagnóstico.
  \item Depuración manual de los PNG: algunos pacientes tenían estudios del segmento cervical o de rodillas, y otros cortes T2 sagitales quedaban oscuros o con aberraciones; se eliminaron.
\end{enumerate}
Tras esta depuración quedaron \textbf{1\,086 cortes sagitales de 414 pacientes} (2,6 cortes por paciente en promedio; entre 1 y 6).

\subsection{Recorte de la región lumbar y de los ROI por disco}
\label{sec:roi}
\paragraph{Recorte L1--S1.} Con \arch{aplicar_recortes.py} se recortó cada corte en altura para abarcar principalmente L1--S1, replicando por \emph{template matching} unos recortes manuales previos (mediana de la puntuación de coincidencia 0,9996; solo el corte 0547\_008 quedó por debajo de 0,99). El recorte se extendió 60~px hacia arriba y hacia abajo hasta el borde de la imagen; con ello el relleno necesario en L1-L2 bajó de 779 de 1\,080 recortes (72\,\%) a cero.

\paragraph{Detector de discos.} Para localizar los discos se entrenó un detector con un conjunto propio en Roboflow (proyecto \emph{Intervertebral Disks}, versión~1): 100 cortes (70 de entrenamiento, 20 de validación y 10 de prueba), con orientación automática (\emph{auto-orient}) y reescalado a 640$\times$640 con bordes negros, sin aumento de datos en Roboflow. Las cinco clases del detector son los niveles L1-L2 a L5-S1. Se ajustó un YOLOv8n de detección de cajas (\arch{yolov8n.pt}) durante 50 épocas, con paciencia 30, lote 16, imagen de 640~px y semilla 42; el aumento en línea de Ultralytics (mosaico, escala, traslación, brillo) es el predeterminado, con reflexión horizontal. La mejor época en validación fue la 49. El detector solo ubica el nivel del disco; no clasifica patología. La Tabla~\ref{tab:detector} resume su desempeño con el umbral de uso (confianza 0,35).

\begin{table}[t]
\centering\small
\caption{Desempeño del detector de discos (corrida \arch{detector_discos_v1_s42}).}
\label{tab:detector}
\begin{tabular}{lrrrrrrr}
\toprule
Partición & Cortes & Precisión & Recall & mAP50 & mAP50-95 & Cortes con 5 niveles & IoU medio \\
\midrule
Validación & 20 & 0,922 & 0,964 & 0,962 & 0,663 & 95\,\% & 0,819 \\
Prueba     & 10 & 0,983 & 0,994 & 0,995 & 0,691 & 100\,\% & 0,843 \\
\bottomrule
\end{tabular}
\end{table}

El nivel más difícil es L1-L2 (en validación: precisión 0,839, recall 0,820 y IoU 0,74); los demás tienen recall de 1,0. Con solo 20 y 10 cortes, un corte mueve varios puntos porcentuales; es una sola corrida con una sola semilla, y la época se eligió con la validación, por lo que sus cifras están algo infladas.

\paragraph{Extracción del ROI.} Con \arch{autocrop_rois.py} se detecta sobre el recorte y se extrae el parche sobre el PNG completo, sumando el desplazamiento. La caja se amplía (margen posterior 0,40, anterior 0,10 y vertical 0,50), se hace cuadrada y se guarda a 128$\times$128~px, una por nivel (L1-L2, L2-L3, L3-L4, L4-L5 y L5-S1). Cada recorte queda registrado en \arch{roi_log.csv} con su caja, su ordenada central, su confianza y el relleno por lado. El resultado fueron \textbf{5\,385 ROI}. El margen anterior de L5-S1 se probó en 0,20 frente a 0,10 y la ganancia fue marginal, por lo que se mantuvo 0,10. La Figura~\ref{fig:rois} muestra los cinco ROI de un corte.

\begin{figure}[t]
\centering
\includegraphics[width=0.9\linewidth]{guia_visual_5rois.png}
\caption{Los cinco ROI (L1-L2 a L5-S1) de un corte sagital T2. \pendiente{confirmar que es la figura que se quiere en el informe o sustituirla por un montaje propio con las cajas del detector}}
\label{fig:rois}
\end{figure}

\paragraph{Exclusiones por calidad.} Se aplicaron reglas objetivas sobre \arch{roi_log.csv} (función \arch{elegibles()} de \arch{scripts/build_dataset.py}): se descarta el corte completo si falta algún nivel, si dos niveles adyacentes se solapan (IoU $>0{,}5$) o si el orden vertical de los centros no crece de L1-L2 a L5-S1; además se excluyen el paciente 0456, el corte 0547\_008 y los ROI con confianza del detector menor que 0,5. En total se excluyeron \textbf{262 ROI}, que afectan a 48 pacientes; 15 pacientes quedaron fuera por completo. Así, los 5\,385 ROI etiquetados quedan en \textbf{5\,123 ROI de 399 pacientes}; los 15 pacientes restantes (de los 414) no conservaron ningún ROI elegible. De los ROI excluidos, 207 eran Protrusión, 38 Normal y 17 Degeneración.

\subsection{Anotación y retiquetado}
\paragraph{Clases.} La guía de anotación (\arch{docs/guias/guia_anotacion_4_clases.md}) define cuatro clases nominales: Normal, Degeneración, Abombamiento y Hernia. Cuando coinciden varias, se aplica la prioridad Hernia $>$ Abombamiento $>$ Degeneración $>$ Normal, que solo sirve para desempatar el etiquetado y no es un orden de gravedad. La guía recoge criterios de uso común y no está validada por un radiólogo. Las etiquetas las asignaron estudiantes, sin radiólogo.

\paragraph{Fusión a tres clases.} En la primera corrida del clasificador semilla, con las cuatro clases (60 imágenes de validación) la exactitud fue 0,65 (IC 95\,\%: 0,52--0,76) y 10 de los 20 errores fueron confusiones Hernia--Abombamiento. Sin un corte axial, la extensión del abultamiento sobre la circunferencia del disco, que es lo que separa ambas entidades en la nomenclatura de consenso \cite{fardon2014}, no puede medirse. Por eso Hernia y Abombamiento se fusionaron en \textbf{Protrusión}, quedando tres clases: Normal, Degeneración discal y Protrusión. Otros autores también reportan confusión entre abultamiento y protrusión \cite{gao2025}. La fusión es una decisión posterior a ver los datos (\emph{post hoc}) y se declara así.

\paragraph{Preetiquetado y revisión.} Etiquetar 5\,385 ROI a mano desde cero era inviable en el plazo, así que se siguió un esquema de aprendizaje activo. Primero se etiquetó a mano una semilla en Roboflow (proyecto \emph{deepspine-seed-350-cls}); su versión~3, la usada para el preetiquetado final, tiene 719 imágenes de entrenamiento y 84 de validación (24 Degeneración, 30 Normal y 30 Protrusión). Con ella se entrenó un clasificador pequeño, YOLOv8s-cls (entrada 224$\times$224, hasta 60 épocas, paciencia 15, lote 32, semilla 42; cuaderno \arch{entrenamiento_clasificador_semilla_v3.ipynb}). En validación alcanzó una exactitud de 0,833 (IC 95\,\%: 0,739--0,898) y un F1 de 0,784 en Degeneración, 0,852 en Normal y 0,857 en Protrusión; con solo 84 imágenes de validación estas cifras son una guía y no una garantía. El modelo preetiquetó los 5\,385 ROI con su clase y su confianza (media 0,946; mediana 1,0) y marcó 2\,265 (42,1\,\%) para revisión prioritaria: 1\,881 por predecir una clase rara (Protrusión) y el resto por confianza menor que 0,7, por margen menor que 0,2 entre las dos primeras clases, o por combinación de motivos.

Después, todos los ROI pasaron por revisión humana en dos grupos, A y B, de 2\,693 y 2\,692 ROI respectivamente. La clase final es la decisión humana. De los 5\,385 ROI, 3\,366 (62,5\,\%) conservaron la clase del modelo y \textbf{2\,019 (37,5\,\%) fueron corregidos}. La Tabla~\ref{tab:reetiquetado} muestra cómo se reorganizó el dataset.

\begin{table}[t]
\centering\small
\caption{Clase propuesta por el preetiquetador frente a la clase final tras la revisión humana (5\,385 ROI).}
\label{tab:reetiquetado}
\begin{tabular}{lrrrrr}
\toprule
 & \multicolumn{3}{c}{Clase final} & & \\
\cmidrule(lr){2-4}
Predicha & Normal & Degeneración & Protrusión & Total & Corregidas \\
\midrule
Normal       & 1\,086 & 32  & 1\,387 & 2\,505 & 56,6\,\% \\
Degeneración & 51     & 256 & 534    & 841    & 69,6\,\% \\
Protrusión   & 6      & 9   & 2\,024 & 2\,039 & 0,7\,\% \\
\midrule
Total final  & 1\,143 & 297 & 3\,945 & 5\,385 & 37,5\,\% \\
\bottomrule
\end{tabular}
\end{table}

La lectura de la Tabla~\ref{tab:reetiquetado} es clara: el preetiquetador, entrenado con una semilla balanceada, \emph{subestimó} la Protrusión. De los 3\,945 ROI de Protrusión finales, 1\,921 (48,7\,\%) llegaron como Normal o Degeneración, mientras que casi nadie de lo que predijo como Protrusión cambió (15 de 2\,039). Dicho de otro modo, la revisión humana fue lo que fijó la distribución de clases del dataset, no el modelo, y la distribución real está muy desbalanceada (Sección~\ref{sec:balance}).

\paragraph{Fiabilidad de las etiquetas.} Cada ROI fue revisado por una sola persona (A o B); 408 de los 414 pacientes tienen ROI revisados por ambos, pero no hay ROI doble. En consecuencia \textbf{no puede calcularse el coeficiente kappa entre anotadores}. Sí se observan diferencias de criterio: A etiquetó Degeneración en el 6,9\,\% de sus ROI y B en el 4,1\,\%, y corrigió el 35,3\,\% frente al 39,7\,\% de B. No es posible separar cuánto de esa diferencia es criterio y cuánto es composición de la muestra.

\subsection{División por paciente}
\label{sec:split}
Cada paciente aporta unos 13 ROI (5\,123/399) de cortes vecinos casi idénticos. Si se dividiera por imagen, un paciente quedaría repartido entre entrenamiento y prueba y el modelo podría reconocer al paciente en lugar de la patología, inflando las métricas (fuga de datos) \cite{brunekreef2024,rumala2023}. Por eso \textbf{cada paciente va entero a una sola partición}.

El reparto es 70/15/15 por pacientes, con semilla 42. El script prueba 20\,000 repartos aleatorios y conserva el que mejor mantiene, en las tres particiones, la proporción de cada clase (con peso triple) y de cada nivel. Se eligieron tres particiones porque la validación sirve para decidir (época, hiperparámetros, umbrales) y la prueba para medir una sola vez al final, con todo ya decidido. La Tabla~\ref{tab:split} muestra el resultado.

\begin{table}[t]
\centering\small
\caption{División por paciente y ROI originales por clase.}
\label{tab:split}
\begin{tabular}{lrrrrrr}
\toprule
Partición & Pacientes & Normal & Degeneración & Protrusión & ROI & \% ROI \\
\midrule
Entrenamiento & 279 & 771   & 194 & 2\,612 & 3\,577 & 69,8 \\
Validación    & 60  & 170   & 44  & 564    & 778    & 15,2 \\
Prueba        & 60  & 164   & 42  & 562    & 768    & 15,0 \\
\midrule
Total         & 399 & 1\,105 & 280 & 3\,738 & 5\,123 & 100 \\
\bottomrule
\end{tabular}
\end{table}

La intersección de pacientes entre particiones es vacía y se comprueba con aserciones en el script. Los pacientes con al menos un ROI de Degeneración son 81, 17 y 22 en entrenamiento, validación y prueba. La docente aceptó esta división, y los 35\,007 registros que exige la rúbrica son las 35\,007 imágenes del dataset.

\subsection{Construcción del dataset de 35\,007 imágenes}
\label{sec:aumento}
La rúbrica exige un dataset balanceado con más de 35\,000 registros. Como los ROI originales son 5\,123, se recurrió a aumento de datos, siempre \emph{después} de dividir y \emph{dentro} de cada partición, para que ninguna copia cruce de partición. La cantidad de copias por ROI se calcula para que cada clase alcance la misma meta por partición: 8\,167 imágenes por clase en entrenamiento y 1\,751 en validación y en prueba, es decir 35\,007 en total.

\paragraph{Transformaciones.} Rotación de $\pm 8^{\circ}$, desplazamiento de hasta $\pm 8\,\%$ del tamaño, zoom entre 0,90 y 1,10, y brillo y contraste de $\pm 10\,\%$, con reflejo en los bordes. No se usan reflexiones (el eje anterior--posterior tiene significado anatómico), deformación elástica ni realce. Cada copia se genera con una semilla derivada de un hash (SHA-256) del nombre del ROI, la semilla global y el índice de la copia, por lo que el proceso es determinista y reproducible.

\begin{table}[t]
\centering\small
\caption{Imágenes finales por partición y clase (mismo número por clase en cada partición), ROI originales y copias por original.}
\label{tab:35k}
\begin{tabular}{llrrrr}
\toprule
Partición & Clase & Originales & Aumentadas & Total & Copias por original \\
\midrule
Entrenamiento & Normal       & 771    & 7\,396 & 8\,167 & 10,6 \\
              & Degeneración & 194    & 7\,973 & 8\,167 & 42,1 \\
              & Protrusión   & 2\,612 & 5\,555 & 8\,167 & 3,1  \\
Validación    & Normal       & 170    & 1\,581 & 1\,751 & 10,3 \\
              & Degeneración & 44     & 1\,707 & 1\,751 & 39,8 \\
              & Protrusión   & 564    & 1\,187 & 1\,751 & 3,1  \\
Prueba        & Normal       & 164    & 1\,587 & 1\,751 & 10,7 \\
              & Degeneración & 42     & 1\,709 & 1\,751 & 41,7 \\
              & Protrusión   & 562    & 1\,189 & 1\,751 & 3,1  \\
\midrule
Total         &              & 5\,123 & 29\,884 & 35\,007 & \\
\bottomrule
\end{tabular}
\end{table}

\begin{figure}[t]
\centering
\begin{tikzpicture}
\begin{axis}[
  ybar, bar width=14pt, width=0.85\linewidth, height=5.6cm,
  symbolic x coords={Normal,Degeneración,Protrusión}, xtick=data,
  ymin=0, ymax=9800, ylabel={Imágenes (entrenamiento)},
  enlarge x limits=0.25, legend style={at={(0.5,1.04)},anchor=south,legend columns=2,font=\footnotesize},
  nodes near coords, every node near coord/.append style={font=\tiny},
  tick label style={font=\footnotesize}, ylabel style={font=\footnotesize}]
\addplot coordinates {(Normal,771) (Degeneración,194) (Protrusión,2612)};
\addplot coordinates {(Normal,8167) (Degeneración,8167) (Protrusión,8167)};
\legend{Originales, Tras el aumento}
\end{axis}
\end{tikzpicture}
\caption{Balance por conteo en entrenamiento: ROI originales y total tras el aumento.}
\label{fig:balance}
\end{figure}

\paragraph{Trazabilidad.} Cada archivo se llama \arch{T2_TSE_SAG__0188_008_L4-L5__orig.png} (original) o \arch{..._aug012.png} (copia), con paciente, corte y nivel. El manifiesto \arch{manifest_35k.csv} registra, para cada una de las 35\,007 imágenes, su partición, clase, paciente, nivel, ROI original, revisor y, en las copias, los parámetros exactos (semilla, ángulo, zoom, desplazamientos, contraste y brillo); \arch{split_pacientes.csv} guarda la partición de cada paciente. Las copias respetan los rangos declarados (ángulo entre $-8^{\circ}$ y $8^{\circ}$, zoom entre 0,90 y 1,10, contraste entre 0,90 y 1,10).

\paragraph{Verificaciones.} El script comprueba que (i) los conjuntos de pacientes de las tres particiones son disjuntos, (ii) toda copia está en la misma partición que su original, (iii) los nombres de archivo son únicos, (iv) todo PNG está íntegro y (v) el total supera los 35\,000.

\subsection{Balance de clases y sesgo por nivel}
\label{sec:balance}
El balance de 1/3 por clase es \emph{por conteo}, no por originales. Entre los 5\,123 ROI originales, la Protrusión es el 73,0\,\% (3\,738), Normal el 21,6\,\% (1\,105) y Degeneración el 5,5\,\% (280); un clasificador que siempre dijera Protrusión acertaría cerca del 73\,\%. Las copias igualan la contribución de cada clase al entrenamiento, pero no añaden anatomía nueva: el tamaño efectivo del dataset es el número de pacientes (399), y la Degeneración, repetida unas 42 veces, solo tiene 42 originales en la partición de prueba, lo que deja un margen de error amplio en su recall. Para que las métricas no se inflen con copias casi idénticas, los resultados principales se calculan sobre los \textbf{originales} de la partición de prueba, y el conjunto completo se reporta aparte. Además, un modelo entrenado con prior uniforme tiende a sobrepredecir las clases raras en la población real.

\begin{table}[t]
\centering\small
\caption{ROI originales elegibles por nivel y clase.}
\label{tab:nivel}
\begin{tabular}{lrrrrr}
\toprule
Nivel & Normal & Degeneración & Protrusión & Total & \% Protrusión \\
\midrule
L1-L2 & 505 & 121 & 401    & 1\,027 & 39,0 \\
L2-L3 & 383 & 94  & 550    & 1\,027 & 53,6 \\
L3-L4 & 172 & 44  & 811    & 1\,027 & 79,0 \\
L4-L5 & 15  & 5   & 1\,005 & 1\,025 & 98,0 \\
L5-S1 & 30  & 16  & 971    & 1\,017 & 95,5 \\
\bottomrule
\end{tabular}
\end{table}

La Tabla~\ref{tab:nivel} muestra un sesgo importante: casi toda la patología se concentra en L4-L5 y L5-S1. Un clasificador que solo adivinara la clase más frecuente de cada nivel acertaría 3\,842 de 5\,123 ROI (75,0\,\%), de modo que el nivel podría actuar como atajo. Por ello los resultados se reportan también por nivel. Esto relativiza cualquier umbral de exactitud: un 95\,\% medido sobre estas etiquetas mide acuerdo con etiquetas de estudiantes y puede alcanzarse en parte por el atajo, y cifras internas altas suelen bajar al evaluar con datos externos \cite{yi2023}. Por eso se acompaña la exactitud con F1 macro y recall por clase, que además es lo que justifica el uso de F1 con clases desbalanceadas \cite{alkubaisi2022}.

\subsection{Viabilidad del dataset}
\begin{itemize}
  \item \textbf{Disponibilidad y licencia.} Los datos son públicos y los scripts reproducen el procesamiento a partir de ellos. Las imágenes del dataset construido no se suben al repositorio: se distribuyen como un archivo comprimido de 405~MB y el repositorio guarda sus manifiestos.
  \item \textbf{Tamaño.} Cumple el requisito de más de 35\,000 registros. El contenido informativo real es menor: 5\,123 ROI de 399 pacientes, con 280 originales de Degeneración.
  \item \textbf{Balance.} Equilibrado por conteo entre las tres clases y en cada partición; desbalanceado en originales y por nivel (Secciones~\ref{sec:aumento} y~\ref{sec:balance}).
  \item \textbf{Reproducibilidad.} Semilla única (42), manifiesto por imagen y verificaciones automáticas.
  \item \textbf{Cómputo.} Los cuatro modelos se entrenaron con este mismo dataset en una GPU de consumo (RTX 3060~Ti), lo que indica que el tamaño es manejable con recursos modestos.
\end{itemize}

\subsection{Limitaciones}
(1) Las etiquetas son de estudiantes, sin radiólogo, con apoyo de un modelo, y no se midió el acuerdo entre anotadores. (2) Solo hay plano sagital; Hernia y Abombamiento no son separables y se fusionaron después de ver los datos. (3) Un único equipo de resonancia de 1,5~T, sin diversidad de escáner. (4) El tamaño efectivo son 399 pacientes; las copias no aportan anatomía nueva. (5) Degeneración cuenta con pocos pacientes (81, 17 y 22 en las tres particiones). (6) Sesgo por nivel y prior uniforme (Sección~\ref{sec:balance}). (7) Con el identificador de paciente vacío no puede descartarse la presencia de un mismo paciente en dos carpetas. (8) El conjunto es una herramienta académica, no diagnóstica.

% =====================================================================
% Seccion: Seleccion de los cuatro modelos (Hito 2)
% Fragmento SIN preambulo: % =====================================================================
% Seccion: Seleccion de los cuatro modelos (Hito 2)
% Fragmento SIN preambulo: % =====================================================================
% Seccion: Seleccion de los cuatro modelos (Hito 2)
% Fragmento SIN preambulo: \input{secciones/seleccion_modelos}
% Paquetes: booktabs, xcolor. Usa \cite con las claves de referencias_dataset.bib
% Cifras de rendimiento: README.md y docs/resultados/ del repositorio.
% =====================================================================
\providecommand{\pendiente}[1]{\textcolor{red}{\textbf{[PENDIENTE: #1]}}}
\providecommand{\arch}[1]{\texttt{\detokenize{#1}}}

\section{Selección de los modelos}
\label{sec:modelos}

\subsection{Criterios de selección}
La rúbrica pide cuatro modelos de aprendizaje automático acordes con la solución y respaldados por el estado del arte y el marco teórico. Se aplicaron cuatro criterios: (i) que sirvan para la tarea, clasificar un ROI de 128$\times$128~px en tres clases; (ii) que tengan antecedentes en resonancia lumbar o en tareas médicas comparables; (iii) que representen \emph{familias distintas}, para que la comparación diga algo; y (iv) que puedan entrenarse y evaluarse con la misma división por paciente y las mismas métricas, que es la condición para comparar con justicia.

\subsection{Los cuatro modelos}
\begin{table}[t]
\centering\small
\caption{Modelos seleccionados, su familia y su respaldo.}
\label{tab:modelos}
\begin{tabular}{p{0.04\linewidth}p{0.24\linewidth}p{0.26\linewidth}p{0.34\linewidth}}
\toprule
\# & Modelo & Familia & Respaldo principal \\
\midrule
1 & YOLOv8s-cls con ajuste fino & CNN de la familia YOLO & YOLO en resonancia lumbar \cite{diagnostics2023,frontiers2023,sayed2025,wang2025} \\
2 & ResNet-50 con ajuste fino & CNN residual & ResNet-50 en columna lumbar \cite{sayed2025,alkubaisi2022} \\
3 & SVM (núcleo RBF) sobre características de la ResNet-50 & Característica profunda + clasificador de margen & ResNet-50 + SVM sobre este mismo dataset \cite{sari2024}; embeddings + SVM \cite{khoiwal2024}; SVM en tabulares \cite{peerj2023} \\
4 & Random Forest (300 árboles) sobre las mismas características & Característica profunda + \emph{ensemble} de árboles & Solo datos tabulares \cite{peerj2023}; es una extensión propia \\
\bottomrule
\end{tabular}
\end{table}

\paragraph{Modelo 1: YOLOv8s-cls.} La familia YOLO es la más usada en los trabajos revisados sobre resonancia lumbar: YOLOv5 sobre este mismo conjunto de datos \cite{diagnostics2023}, YOLOv5x, YOLOv6 y YOLOv7 en cortes sagitales \cite{frontiers2023}, una variante modificada para grado de degeneración y hernia con validación externa \cite{wang2025}, y un flujo en cascada con YOLOv11 para recortar los discos \cite{sayed2025}. Aquí se usa su variante de clasificación, preentrenada en ImageNet y ajustada con AdamW, capas iniciales congeladas y \emph{dropout}. Además es coherente con el resto del pipeline, que ya emplea YOLO para detectar los discos y para preetiquetar.

\paragraph{Modelo 2: ResNet-50.} Es una red convolucional de otra familia (residual), de modo que la comparación con YOLO no es entre dos variantes de lo mismo. Se ajusta con congelado de las primeras capas, \emph{dropout} de 0,3 en la capa final, suavizado de etiquetas y parada temprana por F1 en validación. Tiene respaldo directo en clasificación de discos lumbares: ResNet-50 sobre cortes sagitales \cite{sayed2025} y transferencia desde ImageNet con ResNet-50 y otras redes \cite{alkubaisi2022}.

\paragraph{Modelos 3 y 4: SVM y Random Forest sobre características profundas.} Un clasificador clásico no funciona bien sobre los píxeles crudos, porque trata cada valor por separado y no reconoce la forma. Por eso la ResNet-50 se usa como extractor: elimina su capa final y entrega 2\,048 valores por imagen, que se estandarizan, se reducen a 256 componentes con PCA y se clasifican con SVM de núcleo RBF o con un Random Forest de 300 árboles. El escalador y el PCA se ajustan solo con el entrenamiento, y los hiperparámetros se eligen con la validación. Esta técnica (características profundas más clasificador clásico) es de uso habitual; con este mismo conjunto de datos, ResNet-50 como extractor más SVM reportó 99,42\,\% en una tarea binaria con validación cruzada \cite{sari2024}, y los \emph{embeddings} congelados con SVM o regresión logística igualan a modelos entrenados desde cero en varias tareas médicas \cite{khoiwal2024}. Ambos modelos, además, corresponden a los temas de la asignatura (máquinas de vectores de soporte y árboles).

\subsection{Alcance del respaldo}
El respaldo no es igual de fuerte para los cuatro modelos, y conviene decirlo:
\begin{itemize}
  \item Los modelos 1 y 2 tienen antecedentes directos en resonancia lumbar.
  \item El modelo 3 tiene antecedentes en el mismo conjunto de datos, pero en una tarea binaria y sin que se declare una división por paciente \cite{sari2024}; sirve como respaldo del método, no como meta de rendimiento.
  \item El modelo 4 solo cuenta con la comparación de clasificadores sobre datos tabulares \cite{peerj2023}, donde el mejor resultado fue de un modelo SMOTE-RFE-XGBoost y SVM quedó entre los mejores. No se encontró un trabajo que aplique Random Forest sobre \emph{embeddings} de resonancia lumbar, así que se presenta como una extensión propia, justificada por la asignatura.
  \item Los resultados altos de la literatura (por ejemplo 97--99\,\%) corresponden a conjuntos de prueba pequeños o a tareas binarias, y cifras internas altas bajan en validación externa \cite{yi2023}.
\end{itemize}

\subsection{Decisión de diseño sobre las características}
El diseño inicial usaba la ResNet-50 con los pesos de ImageNet congelados, para que el SVM y el Random Forest no vieran características ajustadas con las etiquetas de entrenamiento. Al evaluarlo en validación, los modelos rindieron poco (exactitud de 77,6\,\% y 75,3\,\% en prueba de originales con \emph{embeddings} de ImageNet sin ajustar). Se adoptó entonces la ResNet-50 ajustada como extractor: el F1 en validación subió de 0,65 a 0,80 en el SVM y de 0,60 a 0,79 en el Random Forest. El costo es que los modelos 3 y 4 \emph{dependen} del modelo 2 y no son independientes de él; las características de entrenamiento están casi perfectamente separadas porque la red las vio ajustándose a sus etiquetas. Para limitar el riesgo, el escalador, el PCA y los clasificadores se ajustan solo con entrenamiento, y la validación y la prueba nunca intervienen en el ajuste. También se entrenaron regresión logística y KNN sobre las mismas características, como línea base del curso, que no cuentan entre los cuatro.

\subsection{Métricas de comparación}
Los cuatro modelos usan la misma división por paciente, la misma semilla (42) y las mismas métricas, calculadas sobre los originales de la partición de prueba, que se evalúa una sola vez: exactitud, precisión, \emph{recall} y F1 por clase y macro, exactitud balanceada, ROC-AUC macro (uno contra el resto) y pérdida logarítmica. La exactitud sola es engañosa con clases desbalanceadas (un clasificador que siempre dijera Protrusión acertaría cerca del 73\,\% de los originales), por eso se acompaña de F1 macro y \emph{recall} por clase \cite{alkubaisi2022}.

\subsection{Aporte frente a lo existente}
Según la revisión realizada, ninguno de los trabajos leídos combina a la vez (i) tres clases con una fusión declarada de abultamiento y protrusión, (ii) división por paciente con evaluación sobre originales, y (iii) la comparación de cuatro modelos (YOLO, ResNet, SVM y Random Forest) sobre el mismo conjunto de prueba y las mismas métricas. Es una afirmación sobre lo leído, no el resultado de una búsqueda exhaustiva.

% Paquetes: booktabs, xcolor. Usa \cite con las claves de referencias_dataset.bib
% Cifras de rendimiento: README.md y docs/resultados/ del repositorio.
% =====================================================================
\providecommand{\pendiente}[1]{\textcolor{red}{\textbf{[PENDIENTE: #1]}}}
\providecommand{\arch}[1]{\texttt{\detokenize{#1}}}

\section{Selección de los modelos}
\label{sec:modelos}

\subsection{Criterios de selección}
La rúbrica pide cuatro modelos de aprendizaje automático acordes con la solución y respaldados por el estado del arte y el marco teórico. Se aplicaron cuatro criterios: (i) que sirvan para la tarea, clasificar un ROI de 128$\times$128~px en tres clases; (ii) que tengan antecedentes en resonancia lumbar o en tareas médicas comparables; (iii) que representen \emph{familias distintas}, para que la comparación diga algo; y (iv) que puedan entrenarse y evaluarse con la misma división por paciente y las mismas métricas, que es la condición para comparar con justicia.

\subsection{Los cuatro modelos}
\begin{table}[t]
\centering\small
\caption{Modelos seleccionados, su familia y su respaldo.}
\label{tab:modelos}
\begin{tabular}{p{0.04\linewidth}p{0.24\linewidth}p{0.26\linewidth}p{0.34\linewidth}}
\toprule
\# & Modelo & Familia & Respaldo principal \\
\midrule
1 & YOLOv8s-cls con ajuste fino & CNN de la familia YOLO & YOLO en resonancia lumbar \cite{diagnostics2023,frontiers2023,sayed2025,wang2025} \\
2 & ResNet-50 con ajuste fino & CNN residual & ResNet-50 en columna lumbar \cite{sayed2025,alkubaisi2022} \\
3 & SVM (núcleo RBF) sobre características de la ResNet-50 & Característica profunda + clasificador de margen & ResNet-50 + SVM sobre este mismo dataset \cite{sari2024}; embeddings + SVM \cite{khoiwal2024}; SVM en tabulares \cite{peerj2023} \\
4 & Random Forest (300 árboles) sobre las mismas características & Característica profunda + \emph{ensemble} de árboles & Solo datos tabulares \cite{peerj2023}; es una extensión propia \\
\bottomrule
\end{tabular}
\end{table}

\paragraph{Modelo 1: YOLOv8s-cls.} La familia YOLO es la más usada en los trabajos revisados sobre resonancia lumbar: YOLOv5 sobre este mismo conjunto de datos \cite{diagnostics2023}, YOLOv5x, YOLOv6 y YOLOv7 en cortes sagitales \cite{frontiers2023}, una variante modificada para grado de degeneración y hernia con validación externa \cite{wang2025}, y un flujo en cascada con YOLOv11 para recortar los discos \cite{sayed2025}. Aquí se usa su variante de clasificación, preentrenada en ImageNet y ajustada con AdamW, capas iniciales congeladas y \emph{dropout}. Además es coherente con el resto del pipeline, que ya emplea YOLO para detectar los discos y para preetiquetar.

\paragraph{Modelo 2: ResNet-50.} Es una red convolucional de otra familia (residual), de modo que la comparación con YOLO no es entre dos variantes de lo mismo. Se ajusta con congelado de las primeras capas, \emph{dropout} de 0,3 en la capa final, suavizado de etiquetas y parada temprana por F1 en validación. Tiene respaldo directo en clasificación de discos lumbares: ResNet-50 sobre cortes sagitales \cite{sayed2025} y transferencia desde ImageNet con ResNet-50 y otras redes \cite{alkubaisi2022}.

\paragraph{Modelos 3 y 4: SVM y Random Forest sobre características profundas.} Un clasificador clásico no funciona bien sobre los píxeles crudos, porque trata cada valor por separado y no reconoce la forma. Por eso la ResNet-50 se usa como extractor: elimina su capa final y entrega 2\,048 valores por imagen, que se estandarizan, se reducen a 256 componentes con PCA y se clasifican con SVM de núcleo RBF o con un Random Forest de 300 árboles. El escalador y el PCA se ajustan solo con el entrenamiento, y los hiperparámetros se eligen con la validación. Esta técnica (características profundas más clasificador clásico) es de uso habitual; con este mismo conjunto de datos, ResNet-50 como extractor más SVM reportó 99,42\,\% en una tarea binaria con validación cruzada \cite{sari2024}, y los \emph{embeddings} congelados con SVM o regresión logística igualan a modelos entrenados desde cero en varias tareas médicas \cite{khoiwal2024}. Ambos modelos, además, corresponden a los temas de la asignatura (máquinas de vectores de soporte y árboles).

\subsection{Alcance del respaldo}
El respaldo no es igual de fuerte para los cuatro modelos, y conviene decirlo:
\begin{itemize}
  \item Los modelos 1 y 2 tienen antecedentes directos en resonancia lumbar.
  \item El modelo 3 tiene antecedentes en el mismo conjunto de datos, pero en una tarea binaria y sin que se declare una división por paciente \cite{sari2024}; sirve como respaldo del método, no como meta de rendimiento.
  \item El modelo 4 solo cuenta con la comparación de clasificadores sobre datos tabulares \cite{peerj2023}, donde el mejor resultado fue de un modelo SMOTE-RFE-XGBoost y SVM quedó entre los mejores. No se encontró un trabajo que aplique Random Forest sobre \emph{embeddings} de resonancia lumbar, así que se presenta como una extensión propia, justificada por la asignatura.
  \item Los resultados altos de la literatura (por ejemplo 97--99\,\%) corresponden a conjuntos de prueba pequeños o a tareas binarias, y cifras internas altas bajan en validación externa \cite{yi2023}.
\end{itemize}

\subsection{Decisión de diseño sobre las características}
El diseño inicial usaba la ResNet-50 con los pesos de ImageNet congelados, para que el SVM y el Random Forest no vieran características ajustadas con las etiquetas de entrenamiento. Al evaluarlo en validación, los modelos rindieron poco (exactitud de 77,6\,\% y 75,3\,\% en prueba de originales con \emph{embeddings} de ImageNet sin ajustar). Se adoptó entonces la ResNet-50 ajustada como extractor: el F1 en validación subió de 0,65 a 0,80 en el SVM y de 0,60 a 0,79 en el Random Forest. El costo es que los modelos 3 y 4 \emph{dependen} del modelo 2 y no son independientes de él; las características de entrenamiento están casi perfectamente separadas porque la red las vio ajustándose a sus etiquetas. Para limitar el riesgo, el escalador, el PCA y los clasificadores se ajustan solo con entrenamiento, y la validación y la prueba nunca intervienen en el ajuste. También se entrenaron regresión logística y KNN sobre las mismas características, como línea base del curso, que no cuentan entre los cuatro.

\subsection{Métricas de comparación}
Los cuatro modelos usan la misma división por paciente, la misma semilla (42) y las mismas métricas, calculadas sobre los originales de la partición de prueba, que se evalúa una sola vez: exactitud, precisión, \emph{recall} y F1 por clase y macro, exactitud balanceada, ROC-AUC macro (uno contra el resto) y pérdida logarítmica. La exactitud sola es engañosa con clases desbalanceadas (un clasificador que siempre dijera Protrusión acertaría cerca del 73\,\% de los originales), por eso se acompaña de F1 macro y \emph{recall} por clase \cite{alkubaisi2022}.

\subsection{Aporte frente a lo existente}
Según la revisión realizada, ninguno de los trabajos leídos combina a la vez (i) tres clases con una fusión declarada de abultamiento y protrusión, (ii) división por paciente con evaluación sobre originales, y (iii) la comparación de cuatro modelos (YOLO, ResNet, SVM y Random Forest) sobre el mismo conjunto de prueba y las mismas métricas. Es una afirmación sobre lo leído, no el resultado de una búsqueda exhaustiva.

% Paquetes: booktabs, xcolor. Usa \cite con las claves de referencias_dataset.bib
% Cifras de rendimiento: README.md y docs/resultados/ del repositorio.
% =====================================================================
\providecommand{\pendiente}[1]{\textcolor{red}{\textbf{[PENDIENTE: #1]}}}
\providecommand{\arch}[1]{\texttt{\detokenize{#1}}}

\section{Selección de los modelos}
\label{sec:modelos}

\subsection{Criterios de selección}
La rúbrica pide cuatro modelos de aprendizaje automático acordes con la solución y respaldados por el estado del arte y el marco teórico. Se aplicaron cuatro criterios: (i) que sirvan para la tarea, clasificar un ROI de 128$\times$128~px en tres clases; (ii) que tengan antecedentes en resonancia lumbar o en tareas médicas comparables; (iii) que representen \emph{familias distintas}, para que la comparación diga algo; y (iv) que puedan entrenarse y evaluarse con la misma división por paciente y las mismas métricas, que es la condición para comparar con justicia.

\subsection{Los cuatro modelos}
\begin{table}[t]
\centering\small
\caption{Modelos seleccionados, su familia y su respaldo.}
\label{tab:modelos}
\begin{tabular}{p{0.04\linewidth}p{0.24\linewidth}p{0.26\linewidth}p{0.34\linewidth}}
\toprule
\# & Modelo & Familia & Respaldo principal \\
\midrule
1 & YOLOv8s-cls con ajuste fino & CNN de la familia YOLO & YOLO en resonancia lumbar \cite{diagnostics2023,frontiers2023,sayed2025,wang2025} \\
2 & ResNet-50 con ajuste fino & CNN residual & ResNet-50 en columna lumbar \cite{sayed2025,alkubaisi2022} \\
3 & SVM (núcleo RBF) sobre características de la ResNet-50 & Característica profunda + clasificador de margen & ResNet-50 + SVM sobre este mismo dataset \cite{sari2024}; embeddings + SVM \cite{khoiwal2024}; SVM en tabulares \cite{peerj2023} \\
4 & Random Forest (300 árboles) sobre las mismas características & Característica profunda + \emph{ensemble} de árboles & Solo datos tabulares \cite{peerj2023}; es una extensión propia \\
\bottomrule
\end{tabular}
\end{table}

\paragraph{Modelo 1: YOLOv8s-cls.} La familia YOLO es la más usada en los trabajos revisados sobre resonancia lumbar: YOLOv5 sobre este mismo conjunto de datos \cite{diagnostics2023}, YOLOv5x, YOLOv6 y YOLOv7 en cortes sagitales \cite{frontiers2023}, una variante modificada para grado de degeneración y hernia con validación externa \cite{wang2025}, y un flujo en cascada con YOLOv11 para recortar los discos \cite{sayed2025}. Aquí se usa su variante de clasificación, preentrenada en ImageNet y ajustada con AdamW, capas iniciales congeladas y \emph{dropout}. Además es coherente con el resto del pipeline, que ya emplea YOLO para detectar los discos y para preetiquetar.

\paragraph{Modelo 2: ResNet-50.} Es una red convolucional de otra familia (residual), de modo que la comparación con YOLO no es entre dos variantes de lo mismo. Se ajusta con congelado de las primeras capas, \emph{dropout} de 0,3 en la capa final, suavizado de etiquetas y parada temprana por F1 en validación. Tiene respaldo directo en clasificación de discos lumbares: ResNet-50 sobre cortes sagitales \cite{sayed2025} y transferencia desde ImageNet con ResNet-50 y otras redes \cite{alkubaisi2022}.

\paragraph{Modelos 3 y 4: SVM y Random Forest sobre características profundas.} Un clasificador clásico no funciona bien sobre los píxeles crudos, porque trata cada valor por separado y no reconoce la forma. Por eso la ResNet-50 se usa como extractor: elimina su capa final y entrega 2\,048 valores por imagen, que se estandarizan, se reducen a 256 componentes con PCA y se clasifican con SVM de núcleo RBF o con un Random Forest de 300 árboles. El escalador y el PCA se ajustan solo con el entrenamiento, y los hiperparámetros se eligen con la validación. Esta técnica (características profundas más clasificador clásico) es de uso habitual; con este mismo conjunto de datos, ResNet-50 como extractor más SVM reportó 99,42\,\% en una tarea binaria con validación cruzada \cite{sari2024}, y los \emph{embeddings} congelados con SVM o regresión logística igualan a modelos entrenados desde cero en varias tareas médicas \cite{khoiwal2024}. Ambos modelos, además, corresponden a los temas de la asignatura (máquinas de vectores de soporte y árboles).

\subsection{Alcance del respaldo}
El respaldo no es igual de fuerte para los cuatro modelos, y conviene decirlo:
\begin{itemize}
  \item Los modelos 1 y 2 tienen antecedentes directos en resonancia lumbar.
  \item El modelo 3 tiene antecedentes en el mismo conjunto de datos, pero en una tarea binaria y sin que se declare una división por paciente \cite{sari2024}; sirve como respaldo del método, no como meta de rendimiento.
  \item El modelo 4 solo cuenta con la comparación de clasificadores sobre datos tabulares \cite{peerj2023}, donde el mejor resultado fue de un modelo SMOTE-RFE-XGBoost y SVM quedó entre los mejores. No se encontró un trabajo que aplique Random Forest sobre \emph{embeddings} de resonancia lumbar, así que se presenta como una extensión propia, justificada por la asignatura.
  \item Los resultados altos de la literatura (por ejemplo 97--99\,\%) corresponden a conjuntos de prueba pequeños o a tareas binarias, y cifras internas altas bajan en validación externa \cite{yi2023}.
\end{itemize}

\subsection{Decisión de diseño sobre las características}
El diseño inicial usaba la ResNet-50 con los pesos de ImageNet congelados, para que el SVM y el Random Forest no vieran características ajustadas con las etiquetas de entrenamiento. Al evaluarlo en validación, los modelos rindieron poco (exactitud de 77,6\,\% y 75,3\,\% en prueba de originales con \emph{embeddings} de ImageNet sin ajustar). Se adoptó entonces la ResNet-50 ajustada como extractor: el F1 en validación subió de 0,65 a 0,80 en el SVM y de 0,60 a 0,79 en el Random Forest. El costo es que los modelos 3 y 4 \emph{dependen} del modelo 2 y no son independientes de él; las características de entrenamiento están casi perfectamente separadas porque la red las vio ajustándose a sus etiquetas. Para limitar el riesgo, el escalador, el PCA y los clasificadores se ajustan solo con entrenamiento, y la validación y la prueba nunca intervienen en el ajuste. También se entrenaron regresión logística y KNN sobre las mismas características, como línea base del curso, que no cuentan entre los cuatro.

\subsection{Métricas de comparación}
Los cuatro modelos usan la misma división por paciente, la misma semilla (42) y las mismas métricas, calculadas sobre los originales de la partición de prueba, que se evalúa una sola vez: exactitud, precisión, \emph{recall} y F1 por clase y macro, exactitud balanceada, ROC-AUC macro (uno contra el resto) y pérdida logarítmica. La exactitud sola es engañosa con clases desbalanceadas (un clasificador que siempre dijera Protrusión acertaría cerca del 73\,\% de los originales), por eso se acompaña de F1 macro y \emph{recall} por clase \cite{alkubaisi2022}.

\subsection{Aporte frente a lo existente}
Según la revisión realizada, ninguno de los trabajos leídos combina a la vez (i) tres clases con una fusión declarada de abultamiento y protrusión, (ii) división por paciente con evaluación sobre originales, y (iii) la comparación de cuatro modelos (YOLO, ResNet, SVM y Random Forest) sobre el mismo conjunto de prueba y las mismas métricas. Es una afirmación sobre lo leído, no el resultado de una búsqueda exhaustiva.

% =====================================================================
% Seccion: Modelo matematico de los clasificadores y su ensamble
% Fragmento SIN preambulo: % =====================================================================
% Seccion: Modelo matematico de los clasificadores y su ensamble
% Fragmento SIN preambulo: % =====================================================================
% Seccion: Modelo matematico de los clasificadores y su ensamble
% Fragmento SIN preambulo: \input{secciones/modelo_matematico}
% Generada a partir de docs/modelo_matematico/README.md
% Paquetes: amsmath, amssymb, booktabs, graphicx
% =====================================================================
\section{Modelo matemático}
\label{sec:matematico}

Esta sección explica, con sus ecuaciones, cómo cada uno de los cuatro modelos convierte una imagen en una predicción, cómo aprende sus parámetros y cómo se evalúa. Sigue la secuencia vista en clase: datos, variables, modelo, predicción, error, optimización y evaluación. Todas las cifras salen de los modelos entrenados.


\subsection{Planteamiento: X e y}
\label{mm:1}

\textbf{Datos.} $\mathcal{D} = \lbrace (x_i, y_i)\rbrace _{i=1}^{N}$, con $N = 35\,007$ imágenes de 399 pacientes, divididas \textbf{por paciente} en train (24\,501 imágenes, 279 pacientes), validación (5\,253, 60) y test (5\,253, 60).

\textbf{Variables de entrada (X).} Cada imagen es una matriz de intensidades. El recorte (escala de grises) se redimensiona a $224 \times 224$ píxeles y se repite en los 3 canales que esperan las redes preentrenadas. Los valores se llevan a $[0, 1]$. La ResNet-50 (y, a través de ella, los modelos 3 y 4) normaliza además cada canal $c$ con la media y la desviación de ImageNet, igual que \arch{StandardScaler}; YOLOv8 usa los valores en $[0, 1]$ sin normalizar:

\begin{center}
\ajusta{x \in [0,1]^{3 \times 224 \times 224}, \qquad \tilde{x}_{c,h,w} = \frac{x_{c,h,w} - \mu_c}{\sigma_c}, \quad \mu = (0{,}485;\ 0{,}456;\ 0{,}406),\ \sigma = (0{,}229;\ 0{,}224;\ 0{,}225)}
\end{center}

En un modelo tabular, X son columnas como "horas de estudio" o "asistencia". Aquí X son \textbf{150\,528 valores de píxel}. Por eso los modelos 1 y 2 aprenden primero a extraer características, y los modelos 3 y 4 reciben esas características ya extraídas.

\textbf{Variable objetivo (y).} Clasificación \textbf{multiclase} con $K = 3$ clases nominales (sin orden):

\begin{center}
\ajusta{y \in \lbrace 0, 1, 2\rbrace  = \lbrace \text{Degeneración},\ \text{Normal},\ \text{Protrusión}\rbrace , \qquad \mathbf{y} = \text{one-hot}(y) \in \lbrace 0,1\rbrace ^3}
\end{center}

\textbf{Qué se busca.} Una función con parámetros $\theta$ que asigne a cada imagen una probabilidad para cada clase, y luego la clase más probable:

\begin{center}
\ajusta{f_\theta : \mathbb{R}^{3 \times 224 \times 224} \to \Delta^{2} = \Big\lbrace  p \in [0,1]^3 : \textstyle\sum_k p_k = 1 \Big\rbrace , \qquad \hat{y} = \arg\max_k\ p_k}
\end{center}

\begin{table}[htbp]
\centering\small
\begin{tabular}{p{0.313\linewidth}p{0.313\linewidth}p{0.313\linewidth}}
\toprule
Modelo & Qué hace con X & Parámetros aprendidos \\
\midrule
1. YOLOv8s-cls & Red convolucional completa, ajustada de extremo a extremo & 5\,084\,579 pesos \\
2. ResNet-50 & Red convolucional residual, ajustada de extremo a extremo & 23\,514\,179 pesos (225\,344 congelados) \\
3. SVM (RBF) & Clasificador de máximo margen sobre 256 características & 1\,684 vectores de soporte y sus coeficientes $\alpha_i$ \\
4. Random Forest & 300 árboles de decisión sobre las mismas 256 características & Las divisiones (variable, umbral) de cada nodo \\
\bottomrule
\end{tabular}
\end{table}



\subsection{Modelos 1 y 2: redes neuronales convolucionales}
\label{mm:2}

\subsubsection{Capas}
\label{mm:2.1}
\textbf{Convolución.} Un filtro $W$ de tamaño $k \times k$ recorre la imagen y calcula, en cada posición, una suma ponderada de los píxeles vecinos. Detecta bordes, texturas y, en capas profundas, formas como el contorno posterior del disco:

\[
(W * x)_{o,i,j} = \sum_{c} \sum_{u=0}^{k-1} \sum_{v=0}^{k-1} W_{o,c,u,v}\ x_{c,\ i+u,\ j+v} + b_o
\]

\textbf{Normalización por lotes (BatchNorm).} Estandariza cada canal con la media $\mu_B$ y la varianza $\sigma_B^2$ del lote, y luego reescala con dos parámetros aprendidos $\gamma$ y $\beta$:

\[
\hat{h} = \gamma\thinspace  \frac{h - \mu_B}{\sqrt{\sigma_B^2 + \epsilon}} + \beta
\]

\textbf{Activación} (no linealidad, Figura~\ref{mm:fig1}): la ResNet-50 usa $\text{ReLU}(z) = \max(0, z)$ y YOLOv8 usa $\text{SiLU}(z) = z \cdot \sigma(z)$.

\textbf{Bloque residual (ResNet-50).} La entrada se suma a la salida de tres convoluciones ($1{\times}1$, $3{\times}3$, $1{\times}1$). Así la red solo aprende la \emph{corrección} $\mathcal{F}$, y el gradiente fluye bien por 50 capas. La ResNet-50 apila 16 bloques en 4 etapas de 3, 4, 6 y 3 bloques:

\[
h^{(l+1)} = \text{ReLU}\big(\mathcal{F}(h^{(l)};\ W^{(l)}) + h^{(l)}\big)
\]

\textbf{YOLOv8s-cls.} Columna de convoluciones con paso 2 y bloques C2f (convoluciones con conexiones parciales entre etapas, la misma idea de reutilizar la entrada). Termina en una convolución $1{\times}1$ a 1\,280 canales.

\textbf{Promedio global y capa final.} El mapa de características final se promedia por canal. Queda un vector $\phi(x)$ (2\,048 valores en la ResNet-50, 1\,280 en YOLOv8), que pasa por \emph{dropout} ($p = 0{,}3$: durante el entrenamiento se apaga al azar el 30\,\% de los valores) y por una capa lineal:

\[
\phi(x) = \frac{1}{HW}\sum_{i,j} h^{(L)}_{:,i,j}, \qquad z = W\thinspace \phi(x) + b \in \mathbb{R}^3
\]

En la ResNet-50 esta capa tiene $2\,048 \times 3 + 3 = 6\,147$ parámetros.

\subsubsection{De puntajes a probabilidades: softmax}
\label{mm:2.2}
La regresión logística de la semana 9 usa la sigmoide para 2 clases. Con 3 clases se usa su generalización, la \textbf{softmax}, que convierte los 3 puntajes $z_k$ (\emph{logits}) en probabilidades positivas que suman 1:

\[
p_k = P(y = k \mid x) = \frac{e^{z_k}}{\sum_{j=1}^{3} e^{z_j}}
\]

Ejemplo (Figura~\ref{mm:fig1}): los logits $(2{,}0;\ 0{,}5;\ -1{,}0)$ dan las probabilidades $(0{,}786;\ 0{,}175;\ 0{,}039)$, y la predicción es la clase 0. Con $K = 2$ la softmax se reduce exactamente a la sigmoide.

\begin{figure}[htbp]
\centering
\includegraphics[width=0.9\linewidth]{01_activaciones_softmax.png}
\caption{Activaciones y softmax.}
\label{mm:fig1}
\end{figure}


\subsubsection{Función de costo: entropía cruzada}
\label{mm:2.3}
Como en la regresión logística, el error de una imagen se mide con la \textbf{entropía cruzada} (log loss): cuánto se aleja la distribución predicha $p$ de la distribución objetivo $q$. El costo $J(\theta)$ promedia esa pérdida sobre las imágenes:

\begin{center}
\ajusta{\ell(x_i, y_i) = -\sum_{k=1}^{3} q_{ik} \log p_{ik}, \qquad J(\theta) = \frac{1}{N}\sum_{i=1}^{N} \ell(x_i, y_i) + \text{regularización}}
\end{center}
\begin{itemize}
  \item \textbf{YOLOv8s-cls:} $q = \mathbf{y}$ (one-hot), así que $\ell = -\log p_{y}$. Si la red da $p_y = 0{,}01$ a la clase correcta, la pérdida es 4,6; si da 0,99, es 0,01. Penaliza fuerte la \textbf{seguridad equivocada} (Figura~\ref{mm:fig2}).
  \item \textbf{ResNet-50, con suavizado de etiquetas} ($\varepsilon = 0{,}1$): en lugar de exigir probabilidad 1 a la clase real, el objetivo reparte un poco de masa entre las demás:
\end{itemize}

\begin{center}
\ajusta{q_k = (1-\varepsilon)\thinspace \mathbf{y}_k + \frac{\varepsilon}{K} \quad\Rightarrow\quad q = (0{,}933;\ 0{,}033;\ 0{,}033)}
\end{center}

Así la red no puede volverse arbitrariamente segura, lo que reduce el sobreajuste. El costo mínimo posible ya no es 0 sino la entropía de $q$, $H(q) = 0{,}291$. Ese es el piso al que se acerca la pérdida de train (0,305 en la época 6, Figura~\ref{mm:fig2}).

\textbf{¿Por qué no MSE para entrenar?} El MSE de la semana 9 es la función de costo natural cuando $y$ es un número real (regresión). Para probabilidades de clase, la entropía cruzada da gradientes más grandes cuando el error es grande. Además, es la función que se obtiene al maximizar la verosimilitud del modelo. El MSE sí se usa para \textbf{evaluar}: se reporta como puntaje de Brier (Sección~\ref{mm:4.2}) en el registro de desempeño (Sección~\ref{mm:4.4}).

\begin{figure}[htbp]
\centering
\includegraphics[width=0.9\linewidth]{02_funcion_costo.png}
\caption{Función de costo.}
\label{mm:fig2}
\end{figure}


\subsubsection{Optimización: descenso del gradiente con AdamW}
\label{mm:2.4}
\textbf{Descenso del gradiente por lotes.} En cada paso se toma un lote de $B = 64$ imágenes y se mueve cada parámetro en la dirección opuesta al gradiente del costo, como la regla $\theta_j := \theta_j - \alpha\thinspace  \partial J / \partial \theta_j$ de la semana 9:

\[
g_t = \nabla_\theta \frac{1}{B}\sum_{i \in \text{lote}} \ell(x_i, y_i)
\]

El gradiente de cada peso se calcula con la regla de la cadena desde la salida hacia la entrada (retropropagación).

\textbf{AdamW.} Mantiene un promedio móvil del gradiente ($m_t$, el "impulso") y de su cuadrado ($v_t$, la escala). Con eso cada parámetro avanza con un paso adaptado a su propia escala. El término $\lambda\thinspace \theta$ (\emph{weight decay}) encoge los pesos hacia 0, y es la regularización L2:

\begin{center}
\ajusta{m_t = \beta_1 m_{t-1} + (1-\beta_1)\thinspace  g_t, \qquad v_t = \beta_2 v_{t-1} + (1-\beta_2)\thinspace  g_t^2}
\end{center}

\begin{center}
\ajusta{\hat{m}_t = \frac{m_t}{1-\beta_1^t}, \quad \hat{v}_t = \frac{v_t}{1-\beta_2^t}, \qquad \theta_t = \theta_{t-1} - \alpha_t \left( \frac{\hat{m}_t}{\sqrt{\hat{v}_t} + \epsilon} + \lambda\thinspace  \theta_{t-1} \right)}
\end{center}

Valores: $\beta_1 = 0{,}9$ en la ResNet-50 y $0{,}937$ en YOLOv8 (valor por defecto de ultralytics), $\beta_2 = 0{,}999$, $\epsilon = 10^{-8}$, $\lambda = 0{,}01$. Ultralytics no aplica $\lambda$ a los sesgos ni a las capas de normalización.

\textbf{Tasa de aprendizaje en coseno} (Figura~\ref{mm:fig3}). El tamaño del paso $\alpha_t$ empieza en $\alpha_0$ y baja suavemente hasta $\alpha_{\min}$ en $T$ épocas: pasos grandes al principio, finos al final.

\[
\alpha_t = \alpha_{\min} + \tfrac{1}{2}(\alpha_0 - \alpha_{\min})\left(1 + \cos\frac{\pi t}{T}\right)
\]

\begin{table}[htbp]
\centering\small
\begin{tabular}{p{0.313\linewidth}p{0.313\linewidth}p{0.313\linewidth}}
\toprule
 & YOLOv8s-cls & ResNet-50 \\
\midrule
$\alpha_0$ / $\alpha_{\min}$ & $5\cdot10^{-4}$ / $5\cdot10^{-5}$ (1 época de calentamiento) & $10^{-4}$ / 0 \\
Épocas máximas $T$ / paciencia & 15 / 4 & 10 / 4 \\
Capas congeladas (no se actualizan) & primeros 4 módulos & \arch{conv1}, \arch{bn1}, \arch{layer1} (225\,344 pesos) \\
Aumento en línea & recorte aleatorio con el 100\,\% del área (solo varía la proporción, 3:4 a 4:3) y brillo $\pm$5\,\% & rotación $\pm$6$^{\circ}$, desplazamiento 4\,\%, zoom 0,95–1,05, brillo y contraste $\pm$10\,\% \\
Suavizado de etiquetas & no & $\varepsilon = 0{,}1$ \\
Criterio para elegir la época & exactitud top-1 en validación & F1 macro en validación \\
Época elegida / épocas corridas & 2 / 6 & 2 / 6 \\
\bottomrule
\end{tabular}
\end{table}


\textbf{Ajuste fino (\emph{transfer learning}).} Ninguna red parte de cero: $\theta$ se inicializa con los pesos aprendidos en ImageNet (1,28 millones de fotografías), que ya detectan bordes y texturas. Solo la capa final empieza al azar.

\begin{figure}[htbp]
\centering
\includegraphics[width=0.9\linewidth]{03_tasa_aprendizaje.png}
\caption{Tasa de aprendizaje.}
\label{mm:fig3}
\end{figure}


\subsubsection{Diagnóstico: sobreajuste y parada temprana}
\label{mm:2.5}
La Figura~\ref{mm:fig4} muestra el patrón de las dos redes: la pérdida de train sigue bajando, mientras la de validación sube desde la época 2. El modelo empieza a memorizar las imágenes de train, que además incluyen copias aumentadas de las mismas imágenes. Por eso se guarda la época con mejor desempeño en validación y el entrenamiento se detiene tras 4 épocas sin mejora (\textbf{parada temprana}):

\[
e^{*} = \arg\max_{e}\ \text{F1}_{\text{macro}}\big(\text{validación};\ \theta^{(e)}\big)
\]

\begin{figure}[htbp]
\centering
\includegraphics[width=0.9\linewidth]{04_curvas_aprendizaje.png}
\caption{Curvas de aprendizaje.}
\label{mm:fig4}
\end{figure}



\subsection{Modelos 3 y 4: características profundas + clasificador clásico}
\label{mm:3}

Los modelos 3 y 4 no trabajan sobre píxeles. Cada imagen se representa con el vector $\phi(x) \in \mathbb{R}^{2048}$ que produce la ResNet-50 del modelo 2 justo antes de su capa final (Sección~\ref{mm:2.1}). Después se reduce ese vector a 256 dimensiones y se clasifica con un SVM o con un Random Forest.

\textbf{Dependencia declarada:} los modelos 3 y 4 dependen del modelo 2. Con un extractor de ImageNet sin ajustar rinden mucho menos (Sección~\ref{mm:3.4}).

\subsubsection{Estandarización y PCA}
\label{mm:3.1}
\textbf{Estandarización.} Cada una de las 2\,048 características se centra y se escala con la media y la desviación \textbf{calculadas solo con train}. Calcularlas con todo el dataset filtraría información del test (\emph{data leakage}, semanas 6 y 9):

\[
s_j = \frac{\phi_j(x) - \mu_j^{\text{train}}}{\sigma_j^{\text{train}}}
\]

\textbf{Análisis de componentes principales (PCA).} Se busca la base de direcciones ortogonales en la que los datos de train varían más. Son los vectores propios de la matriz de covarianza, y cada imagen se proyecta sobre los $k = 256$ primeros:

\begin{center}
\ajusta{\Sigma = \frac{1}{n-1}\thinspace  S^\top S = V \Lambda V^\top, \qquad u = V_k^\top s \in \mathbb{R}^{256}, \qquad \text{varianza explicada} = \frac{\sum_{j=1}^{k}\lambda_j}{\sum_{j=1}^{2048}\lambda_j}}
\end{center}

Con las características ajustadas, \textbf{256 componentes explican el 80,0\,\% de la varianza} (las 2 primeras ya el 38\,\%). Con las de ImageNet sin ajustar, solo el 58,7\,\% (Figura~\ref{mm:fig5}). Las características ajustadas están más concentradas en pocas direcciones, las que separan las clases.

\begin{figure}[htbp]
\centering
\includegraphics[width=0.9\linewidth]{05_pca_varianza.png}
\caption{Varianza explicada por el PCA.}
\label{mm:fig5}
\end{figure}


\subsubsection{Modelo 3: máquina de vectores de soporte (SVM) con kernel RBF}
\label{mm:3.2}
\textbf{Idea.} Entre todas las fronteras que separan dos clases, el SVM elige la de \textbf{máximo margen}: la que queda más lejos de los ejemplos más cercanos de cada lado. Esos ejemplos se llaman vectores de soporte. Para dos clases, $y_i \in \lbrace -1, +1\rbrace $, con margen blando:

\begin{center}
\ajusta{\min_{w,\thinspace b,\thinspace \xi}\ \ \frac{1}{2}\lVert w \rVert^2 + C \sum_{i=1}^{n} \xi_i \quad \text{sujeto a} \quad y_i\big(w^\top \varphi(u_i) + b\big) \ge 1 - \xi_i,\quad \xi_i \ge 0}
\end{center}
\begin{itemize}
  \item $\frac{1}{2}\lVert w\rVert^2$: minimizarlo maximiza el ancho del margen, $2 / \lVert w \rVert$.
  \item $\xi_i$: holgura que permite que algunas imágenes queden dentro del margen o mal clasificadas.
  \item $C$: equilibrio entre margen ancho y errores en train. Un $C$ pequeño tolera más errores y da una frontera más suave, menos sobreajustada.
\end{itemize}

Es equivalente a minimizar la \textbf{pérdida bisagra} (\emph{hinge}): $\frac{1}{2}\lVert w\rVert^2 + C\sum_i \max\big(0,\ 1 - y_i f(u_i)\big)$.

\textbf{Kernel RBF.} Las clases no se separan con un hiperplano en el espacio original. El SVM resuelve el problema dual, en el que los datos solo aparecen como similitudes $K(u_i, u_j)$; así trabaja implícitamente en un espacio de dimensión infinita (el "truco del kernel"):

\begin{center}
\ajusta{K(u, v) = \exp\big(-\gamma\thinspace \lVert u - v \rVert^2\big), \qquad \gamma = \frac{1}{d \cdot \mathrm{Var}(U)} = 0{,}00061\ \ (\texttt{gamma="scale"},\ d = 256)}
\end{center}

\begin{center}
\ajusta{\max_{\alpha}\ \sum_i \alpha_i - \frac{1}{2}\sum_i\sum_j \alpha_i\alpha_j\thinspace  y_i y_j\thinspace  K(u_i, u_j) \quad \text{s.a.}\ \ 0 \le \alpha_i \le C,\ \ \sum_i \alpha_i y_i = 0}
\end{center}

\[
f(u) = \sum_{i \in \text{SV}} \alpha_i\thinspace  y_i\thinspace  K(u_i, u) + b
\]

Solo los vectores de soporte ($\alpha_i > 0$) intervienen en la predicción. Con $C = 0{,}1$, elegido en validación (Figura~\ref{mm:fig8}), el modelo usa \textbf{1\,684 vectores de soporte de 24\,501 imágenes de train (6,9\,\%)}: 320 de Degeneración, 660 de Normal y 704 de Protrusión.

La Figura~\ref{mm:fig6}, a la derecha, muestra cómo decae la similitud con la distancia. Con el $\gamma$ usado, dos imágenes a la distancia mediana (59) tienen similitud $\approx$ 0,12. Con $\gamma \times 10$ todas las imágenes serían "lejanas" (sobreajuste); con $\gamma / 10$, todas "parecidas" (subajuste).

\textbf{Tres clases.} scikit-learn entrena un SVM por cada par de clases (\emph{uno contra uno}: 3 clasificadores) y combina sus votos en un puntaje por clase; la predicción es la clase con mayor puntaje. Estos puntajes \textbf{no son probabilidades}; por eso el SVM no tiene log loss y la app los muestra normalizados con softmax.

La Figura~\ref{mm:fig6}, a la izquierda, es una \textbf{ilustración en 2 dimensiones} (las 2 primeras componentes, 1\,200 originales de train). En train las clases aparecen muy separadas porque la ResNet ya vio esas imágenes; en validación y test se mezclan más, y por eso la exactitud allí baja a ~88\,\%.

\begin{figure}[htbp]
\centering
\includegraphics[width=0.9\linewidth]{06_svm_rbf.png}
\caption{SVM con kernel RBF.}
\label{mm:fig6}
\end{figure}


\subsubsection{Modelo 4: Random Forest}
\label{mm:3.3}
\textbf{Árbol de decisión.} Hace preguntas del tipo "¿$u_j \le t$?" y divide las imágenes en dos grupos. En cada nodo elige la variable $j$ y el umbral $t$ que más reducen la \textbf{impureza de Gini} (semana 11):

\begin{center}
\ajusta{G(\text{nodo}) = 1 - \sum_{k=1}^{3} p_k^2, \qquad \Delta G = G(\text{padre}) - \frac{n_L}{n}\thinspace G(L) - \frac{n_R}{n}\thinspace G(R)}
\end{center}

$G = 0$ si el nodo tiene una sola clase (puro). El máximo es $1 - 1/K$ si las clases están mezcladas en partes iguales: 0,5 con 2 clases, 0,667 con 3. Ejemplo de la clase: 6 compran y 4 no ($G = 1 - 0{,}6^2 - 0{,}4^2 = 0{,}48$). Tras dividir en un grupo (5, 0) con $G = 0$ y otro (1, 4) con $G = 0{,}32$, la impureza ponderada baja a $0{,}5 \cdot 0 + 0{,}5 \cdot 0{,}32 = 0{,}16$.

\textbf{Bosque.} Un árbol solo sobreajusta. Se entrenan $T = 300$ árboles distintos entre sí por dos fuentes de azar:
\begin{enumerate}
  \item \textbf{Bootstrap:} cada árbol se entrena con una muestra de train tomada con reemplazo, del mismo tamaño.
  \item \textbf{Variables al azar:} en cada división solo se consideran $m = \sqrt{256} = 16$ de las 256 variables.
\end{enumerate}

La predicción promedia las probabilidades de los 300 árboles (votación suave). La probabilidad de cada árbol, $p_t(k \mid u)$, es la proporción de la clase $k$ en la hoja donde cae $u$:

\[
P(y = k \mid u) = \frac{1}{T}\sum_{t=1}^{T} p_t(k \mid u), \qquad \hat{y} = \arg\max_k P(y = k \mid u)
\]

Al promediar árboles diversos baja la varianza (el sobreajuste) sin subir mucho el sesgo. Hiperparámetro elegido en validación (Figura~\ref{mm:fig8}): profundidad máxima 20. Los 300 árboles llegan a esa profundidad y tienen 779 hojas en promedio.

La \textbf{importancia} de cada variable es la disminución total de Gini que aporta, promediada sobre los árboles. Las 10 primeras componentes del PCA concentran el 70\,\% de la importancia (Figura~\ref{mm:fig7}).

\begin{figure}[htbp]
\centering
\includegraphics[width=0.9\linewidth]{07_random_forest.png}
\caption{Random Forest.}
\label{mm:fig7}
\end{figure}


\subsubsection{Extras del curso y comparación con ImageNet}
\label{mm:3.4}
Sobre los mismos 256 componentes se entrenaron, como línea base del curso:
\begin{itemize}
  \item \textbf{Regresión logística multinomial} (semana 9): $p = \text{softmax}(W u + b)$, con regularización L2 de fuerza $1/C$ y $C = 0{,}01$.
  \item \textbf{KNN} (semana 10): predice por mayoría entre los $K = 31$ vecinos más cercanos según la distancia euclídea $d(u, v) = \sqrt{\sum_j (u_j - v_j)^2}$. La escala importa, y por eso se estandariza antes.
\end{itemize}

Con características de ImageNet sin ajustar, el SVM baja de 88,4\,\% a 77,6\,\% y el Random Forest de 88,4\,\% a 75,3\,\% de exactitud en el test de originales. Lo que hace funcionar a los modelos 3 y 4 es la representación aprendida por el modelo 2.

\begin{figure}[htbp]
\centering
\includegraphics[width=0.9\linewidth]{08_busqueda_hiperparametros.png}
\caption{Búsqueda de hiperparámetros.}
\label{mm:fig8}
\end{figure}


\subsubsection{Estructura del sistema: arquitectura híbrida}
\label{mm:3.5}

El sistema tiene \textbf{dos caminos}. YOLOv8s-cls va por el suyo, independiente. La ResNet-50 es la \textbf{base compartida} de los otros tres modelos: sus características $\phi(x)$ alimentan su propia capa final (modelo 2) y, después de estandarizarlas y reducirlas, al SVM (modelo 3) y al Random Forest (modelo 4). Unir una red profunda que extrae características con un clasificador clásico que decide es lo que se llama \textbf{arquitectura híbrida} (Figura~\ref{fig:estructura}).


\textbf{Las 4 funciones.} Sea $g_Y$ la red YOLOv8s-cls hasta sus logits, $g_R$ la ResNet-50 sin su capa final y $\phi(x) = g_R(x)$. Cada modelo es una composición de funciones:

\[
p_1(x) = \text{softmax}\big(g_Y(x)\big)
\]

\[
p_2(x) = \text{softmax}\big(W\thinspace \phi(x) + b\big)
\]

\begin{center}
\ajusta{u(x) = V_k^\top\thinspace D^{-1}\big(\phi(x) - \mu\big), \qquad D = \text{diag}(\sigma_1, \dots, \sigma_{2048})}
\end{center}

\begin{center}
\ajusta{p_3(x) = \text{softmax}\big(\text{SVM}(u(x))\big), \qquad p_4(x) = \frac{1}{T}\sum_{t=1}^{T} p_t\big(u(x)\big)}
\end{center}

Las tres últimas dependen de la misma $\phi(x)$. Por eso, en la aplicación la ResNet-50 se ejecuta una sola vez por imagen. La softmax de $p_3$ solo lleva los puntajes del SVM a la misma escala de los demás: no son probabilidades calibradas.

\textbf{Orden de entrenamiento.} Cada etapa usa solo lo aprendido en las anteriores, y nada se ajusta con el test:

\begin{table}[htbp]
\centering\small
\begin{tabular}{p{0.235\linewidth}p{0.235\linewidth}p{0.235\linewidth}p{0.235\linewidth}}
\toprule
Paso & Qué se aprende & Con qué datos & Qué queda fijo \\
\midrule
1 & Pesos de YOLOv8s-cls & train; época elegida con validación & — \\
2 & Pesos de la ResNet-50, incluidos $W$ y $b$ & train; época elegida con validación & — \\
3 & $\mu$, $D$ y $V_k$ (escalador y PCA) & características de train & La ResNet-50 \\
4 & SVM y Random Forest & $u$ de train; hiperparámetros elegidos con validación & ResNet-50, escalador y PCA \\
\bottomrule
\end{tabular}
\end{table}


\textbf{Por qué es híbrida y no un ensamble.} En esta estructura un modelo le pasa sus \textbf{características} a otro, y cada uno sigue dando su propia predicción. Un ensamble, en cambio, combina las \textbf{respuestas finales} de varios modelos. Esa combinación es la sección siguiente.

\subsubsection{Ensamble por promedio}
\label{mm:3.6}

\textbf{Idea.} Es la misma de la semana 11: consultar a varios modelos reduce la dependencia de una sola decisión. Cada modelo entrega 3 puntajes para la imagen; se promedian por clase y gana la clase con el mayor promedio (\textbf{votación suave}). Con $M = 4$ modelos:

\begin{center}
\ajusta{P_{\text{ens}}(k \mid x) = \frac{1}{M}\sum_{m=1}^{M} p_m(k \mid x), \qquad \hat{y}_{\text{ens}} = \arg\max_k\ P_{\text{ens}}(k \mid x)}
\end{center}

El ensamble \textbf{no se entrena}: todos los modelos pesan igual ($1/4$) y no tiene parámetros que aprender. Se calcula después, con las salidas de los 4 modelos ya entrenados. Por eso es un \textbf{extra} y no un quinto modelo.

\textbf{Ejemplo real} (imagen de test: paciente 0319, L1-L2, etiqueta real Degeneración):

\begin{table}[htbp]
\centering\small
\begin{tabular}{lllll}
\toprule
Modelo & Degeneración & Normal & Protrusión & Predicción \\
\midrule
1. YOLOv8s-cls & 0,109 & 0,871 & 0,020 & Normal (falla) \\
2. ResNet-50 & 0,936 & 0,025 & 0,039 & Degeneración \\
3. SVM (con softmax) & 0,746 & 0,064 & 0,190 & Degeneración \\
4. Random Forest & 0,866 & 0,050 & 0,083 & Degeneración \\
\textbf{Ensamble (promedio)} & \textbf{0,664} & 0,253 & 0,083 & \textbf{Degeneración} \\
\bottomrule
\end{tabular}
\end{table}


Por ejemplo, $P_{\text{ens}}(\text{Degeneración}) = (0{,}109 + 0{,}936 + 0{,}746 + 0{,}866) / 4 = 0{,}664$. Los otros tres modelos corrigen el error de YOLO.

\textbf{Por qué puede ayudar, y su límite.} Si cada modelo tiene una varianza de error $\sigma^2$ y la correlación entre modelos es $\rho$, la varianza del promedio de $M$ modelos es:

\begin{center}
\ajusta{\text{Var}\left(\frac{1}{M}\sum_{m=1}^{M} p_m\right) = \rho\thinspace \sigma^2 + \frac{1-\rho}{M}\thinspace \sigma^2}
\end{center}

Si los modelos fueran independientes ($\rho = 0$), la varianza se dividiría entre 4. Si fueran idénticos ($\rho = 1$), promediar no serviría de nada. Aquí los modelos 2, 3 y 4 comparten $\phi(x)$, así que se espera una correlación alta y una mejora limitada.

\textbf{Resultados} (test de originales; detalle en las Secciones~\ref{mm:4.3} y~\ref{mm:4.4}):

\begin{table}[htbp]
\centering\small
\begin{tabular}{lll}
\toprule
Métrica & Mejor modelo individual & Ensamble \\
\midrule
Exactitud & 0,889 (ResNet-50) & 0,889 \\
F1 macro & 0,769 (ResNet-50) & 0,745 \\
ROC-AUC macro & 0,942 (Random Forest) & \textbf{0,947} \\
MSE (Brier) & 0,173 (ResNet-50) & \textbf{0,164} \\
\bottomrule
\end{tabular}
\end{table}


El ensamble empata en exactitud y da las probabilidades más confiables (mejor ROC-AUC y MSE), pero detecta menos casos de Degeneración: 17 de 42, frente a 21 de la ResNet-50. En validación, el criterio acordado, su F1 macro es 0,822 frente a 0,824 de la ResNet-50. Por eso \textbf{la ResNet-50 sigue siendo el modelo principal} y el ensamble se reporta como resultado adicional.


\subsection{Evaluación}
\label{mm:4}

\subsubsection{Protocolo}
\label{mm:4.1}
\begin{itemize}
  \item Los parámetros $\theta$ se aprenden solo con \textbf{train}.
  \item Los hiperparámetros ($C$, profundidad, $K$, época) se eligen con \textbf{validación}, por F1 macro.
  \item El \textbf{test} se evalúa una sola vez, al final. Ningún paciente aparece en dos particiones, y cada copia aumentada queda en la partición de su original.
  \item Se reporta el test de \textbf{originales} (768 imágenes en proporción natural: 42 Degeneración, 164 Normal, 562 Protrusión) y el \textbf{test completo} (5\,253, balanceado con copias).
\end{itemize}

\subsubsection{Matriz de confusión y métricas}
\label{mm:4.2}
En clase se vieron las fórmulas para \textbf{2 clases} (positiva y negativa):

\begin{center}
\ajusta{\text{Accuracy} = \frac{TP + TN}{TP + TN + FP + FN}, \qquad \text{Precision} = \frac{TP}{TP + FP}, \qquad \text{Recall} = \frac{TP}{TP + FN}, \qquad F1 = 2\cdot\frac{\text{Precision}\cdot\text{Recall}}{\text{Precision} + \text{Recall}}}
\end{center}

Con \textbf{3 clases} se aplican las mismas fórmulas tres veces. Cada clase $k$ se trata como la "positiva" y las otras dos como la "negativa" (uno contra el resto): TP = predijo $k$ y era $k$; FP = predijo $k$ y no era $k$; FN = era $k$ y predijo otra; TN = el resto. Luego se promedian las tres (promedio \emph{macro}). La exactitud no necesita promedio: es el total de aciertos (la diagonal de la matriz) sobre el total de imágenes.

\begin{center}
\ajusta{\text{Exactitud} = \frac{\sum_k \text{TP}_k}{N}, \qquad \text{Precisión}_k = \frac{\text{TP}_k}{\text{TP}_k + \text{FP}_k}, \qquad \text{Recall}_k = \frac{\text{TP}_k}{\text{TP}_k + \text{FN}_k}}
\end{center}

\begin{center}
\ajusta{F1_k = 2\thinspace \frac{\text{Precisión}_k \cdot \text{Recall}_k}{\text{Precisión}_k + \text{Recall}_k}, \qquad F1_{\text{macro}} = \frac{1}{K}\sum_k F1_k, \qquad \text{Exactitud balanceada} = \frac{1}{K}\sum_k \text{Recall}_k}
\end{center}

\textbf{Ejemplo con el modelo 2 en el test de originales} (filas = real, columnas = predicha):

\begin{table}[htbp]
\centering\small
\begin{tabular}{llll}
\toprule
Real \textbackslash{} Predicha & Degeneración & Normal & Protrusión \\
\midrule
\textbf{Degeneración} (42) & \textbf{21} & 6 & 15 \\
\textbf{Normal} (164) & 5 & \textbf{133} & 26 \\
\textbf{Protrusión} (562) & 6 & 27 & \textbf{529} \\
\bottomrule
\end{tabular}
\end{table}


\begin{table}[htbp]
\centering\small
\begin{tabular}{lllllll}
\toprule
Clase & TP & FP & FN & Precisión & Recall & F1 \\
\midrule
Degeneración & 21 & 11 & 21 & 0,656 & 0,500 & 0,568 \\
Normal & 133 & 33 & 31 & 0,801 & 0,811 & 0,806 \\
Protrusión & 529 & 41 & 33 & 0,928 & 0,941 & 0,935 \\
\bottomrule
\end{tabular}
\end{table}

\begin{itemize}
  \item \textbf{Exactitud} = (21 + 133 + 529) / 768 = 683 / 768 = \textbf{0,889}.
  \item \textbf{F1 macro} = (0,568 + 0,806 + 0,935) / 3 = \textbf{0,769}.
\end{itemize}

\textbf{Por qué la exactitud sola engaña.} Un modelo que siempre dijera "Protrusión" acertaría 562 / 768 = \textbf{73,2\,\%} sin detectar ni un caso de Degeneración. Es el mismo ejemplo del fraude visto en clase. El F1 macro pesa igual las tres clases y deja ver que Degeneración se detecta solo en la mitad de los casos.

\textbf{Intervalo de confianza.} Con 768 imágenes, la exactitud de 88,9\,\% tiene un intervalo de Wilson al 95\,\% de \textbf{[86,5\,\%; 91,0\,\%]}:

\begin{center}
\ajusta{\frac{\hat{p} + \frac{z^2}{2n} \pm z\sqrt{\frac{\hat{p}(1-\hat{p})}{n} + \frac{z^2}{4n^2}}}{1 + \frac{z^2}{n}}, \qquad z = 1{,}96}
\end{center}

Las diferencias entre los 4 modelos (87,8\,\% a 88,9\,\%) están dentro de ese margen.

\textbf{ROC-AUC.} Para cada clase, la curva ROC recorre todos los umbrales $\tau$ sobre $p_k$ y grafica la tasa de verdaderos positivos contra la de falsos positivos. El área bajo la curva es la probabilidad de que una imagen de la clase $k$ reciba más puntaje que una que no lo es. Se promedia sobre las 3 clases (uno contra el resto).

\textbf{Log loss.} Es la entropía cruzada de la Sección~\ref{mm:2.3}, evaluada con las probabilidades sobre el conjunto de prueba.

\textbf{MSE (puntaje de Brier).} Es el MSE de la semana 9 aplicado a la clasificación. El "valor real" de cada imagen es su etiqueta en one-hot (por ejemplo, Protrusión = $(0, 0, 1)$) y la "predicción" son las 3 probabilidades del modelo. Se suman los errores al cuadrado de las 3 clases y se promedia sobre las $N$ imágenes:

\begin{center}
\ajusta{\text{MSE} = \frac{1}{N}\sum_{i=1}^{N}\sum_{k=1}^{3}\left(p_{ik} - y_{ik}\right)^2, \qquad \text{RMSE} = \sqrt{\text{MSE}}}
\end{center}

Vale 0 si el modelo da probabilidad 1 a la clase correcta y como máximo 2. Igual que el MSE de clase, castiga más los errores grandes. Ejemplo con una imagen real de Protrusión:
\begin{itemize}
  \item Si el modelo da $p = (0{,}1;\ 0{,}2;\ 0{,}7)$: $0{,}1^2 + 0{,}2^2 + 0{,}3^2 = 0{,}14$ (acierta, con algo de duda).
  \item Si da $p = (0{,}9;\ 0{,}05;\ 0{,}05)$: $0{,}9^2 + 0{,}05^2 + 0{,}95^2 = 1{,}715$ (falla con mucha seguridad).
\end{itemize}

A diferencia de la exactitud, el MSE también mide si la \textbf{confianza} del modelo es la adecuada. No aplica al SVM, porque sus puntajes no son probabilidades.

\subsubsection{Resultados (test de originales)}
\label{mm:4.3}
Los resultados en los originales de prueba están en la Tabla~\ref{tab:resultados}, con sus matrices de confusión y curvas ROC en las Figuras~\ref{fig:matrices} y~\ref{fig:roc}.

\subsubsection{Registro de desempeño (formato de la clase)}
\label{mm:4.4}

Es la tabla "Registro de desempeño" de la semana 9, con columnas Train y Test.
\begin{itemize}
  \item \textbf{Train:} las 3\,577 imágenes originales de entrenamiento.
  \item \textbf{Test:} las 768 originales de test.
  \item \textbf{Test completo:} las 5\,253 imágenes de test con copias aumentadas, como referencia adicional.
\end{itemize}

Precision, recall y F1 son promedios \emph{macro} de las 3 clases. \textbf{R$^2$ no aplica}: es una métrica de regresión y este problema es de clasificación (como advierte la diapositiva, no todas las métricas aplican a todos los modelos).

\begin{table}[htbp]
\centering\small
\caption{Registro de desempeño. Entrenamiento: 3\,577 originales. Prueba: 768 originales. Prueba completa: 5\,253 imágenes con copias, balanceada. Precision, recall y F1 son promedios macro.}
\label{tab:registro}
\begin{tabular}{lllll}
\toprule
Modelo & Métrica & Train & Test & Test completo \\
\midrule
\textbf{1. YOLOv8s-cls} & MSE / RMSE (Brier) & 0,024 / 0,154 & 0,196 / 0,442 & 0,403 / 0,635 \\
 & Accuracy & 0,984 & 0,878 & 0,759 \\
 & Precision / Recall & 0,959 / 0,989 & 0,771 / 0,732 & 0,779 / 0,759 \\
 & F1 / ROC-AUC & 0,973 / 0,999 & 0,750 / 0,931 & 0,752 / 0,923 \\
\textbf{2. ResNet-50} & MSE / RMSE (Brier) & 0,047 / 0,217 & 0,173 / 0,416 & 0,349 / 0,591 \\
 & Accuracy & 0,971 & \textbf{0,889} & 0,775 \\
 & Precision / Recall & 0,948 / 0,981 & 0,795 / 0,751 & 0,799 / 0,775 \\
 & F1 / ROC-AUC & 0,964 / 0,998 & 0,769 / 0,919 & 0,768 / 0,915 \\
\textbf{3. SVM (RBF)} & MSE / RMSE (Brier) & no aplica & no aplica & no aplica \\
 & Accuracy & 0,979 & 0,884 & 0,754 \\
 & Precision / Recall & 0,964 / 0,984 & 0,797 / 0,718 & 0,789 / 0,754 \\
 & F1 / ROC-AUC & 0,974 / 0,998 & 0,748 / 0,899 & 0,746 / 0,905 \\
\textbf{4. Random Forest} & MSE / RMSE (Brier) & 0,010 / 0,098 & 0,174 / 0,416 & 0,370 / 0,608 \\
 & Accuracy & 0,998 & 0,884 & 0,743 \\
 & Precision / Recall & 0,994 / 0,999 & 0,810 / 0,710 & 0,784 / 0,743 \\
 & F1 / ROC-AUC & 0,996 / 1,000 & 0,746 / 0,942 & 0,731 / 0,920 \\
\textbf{Extra: ensamble} & MSE / RMSE (Brier) & 0,031 / 0,175 & 0,164 / 0,405 & 0,334 / 0,578 \\
 & Accuracy & 0,989 & 0,889 & 0,762 \\
 & Precision / Recall & 0,978 / 0,992 & 0,782 / 0,721 & 0,794 / 0,762 \\
 & F1 / ROC-AUC & 0,985 / 1,000 & 0,745 / 0,947 & 0,752 / 0,937 \\
\bottomrule
\end{tabular}
\end{table}


\textbf{Lectura:}
\begin{itemize}
  \item La distancia entre Train y Test es la que se vio en las curvas de aprendizaje (Sección~\ref{mm:2.5}): \textbf{sobreajuste}. En train los modelos casi no se equivocan (MSE de 0,01 a 0,05); en test el MSE sube a ~0,17–0,20.
  \item Por MSE, el ensamble es el mejor en test (0,164), seguido de la ResNet-50 y el Random Forest (0,173 y 0,174).
  \item Todas las cifras salen de \arch{docs/resultados/comparacion/tabla_metricas_larga.csv}, que \arch{scripts/entrenamiento/05_comparar_modelos.py} calcula a partir de los puntajes guardados de cada modelo.
\end{itemize}


\subsection{Supuestos y limitaciones del modelo}
\label{mm:5}
\begin{itemize}
  \item \textbf{Independencia:} se supone que los pacientes son independientes entre sí. Las imágenes de un mismo paciente no lo son, por eso se divide por paciente. Las copias aumentadas no son datos nuevos: el tamaño efectivo es de 399 pacientes, no de 35\,007 imágenes.
  \item \textbf{Etiquetas con ruido:} $y$ la asignaron estudiantes, sin radiólogo y sin corte axial. El modelo aprende a reproducir esas etiquetas, no un diagnóstico.
  \item \textbf{Desbalance:} el entrenamiento ve las clases balanceadas (por copias), pero en la proporción natural el 73\,\% es Protrusión. Se probó corregir por esa proporción multiplicando $p_k$ por la frecuencia de cada clase; no mejoró la exactitud y hundió el recall de Degeneración, así que se descartó (\arch{docs/resultados/experimentos_validacion.md}).
  \item \textbf{Sobreajuste:} las redes llegan a ~97–99\,\% en train y ~89\,\% en test de originales. La regularización (dropout, weight decay, suavizado, aumento) y la parada temprana lo reducen, pero no lo eliminan.
  \item \textbf{Dependencia de los modelos 3 y 4} respecto del modelo 2: comparten la representación $\phi(x)$, así que no son cuatro opiniones independientes. Eso limita lo que puede ganar el ensamble (Sección~\ref{mm:3.6}).
  \item \textbf{Herramienta académica, no diagnóstica.}
\end{itemize}

% Generada a partir de docs/modelo_matematico/README.md
% Paquetes: amsmath, amssymb, booktabs, graphicx
% =====================================================================
\section{Modelo matemático}
\label{sec:matematico}

Esta sección explica, con sus ecuaciones, cómo cada uno de los cuatro modelos convierte una imagen en una predicción, cómo aprende sus parámetros y cómo se evalúa. Sigue la secuencia vista en clase: datos, variables, modelo, predicción, error, optimización y evaluación. Todas las cifras salen de los modelos entrenados.


\subsection{Planteamiento: X e y}
\label{mm:1}

\textbf{Datos.} $\mathcal{D} = \lbrace (x_i, y_i)\rbrace _{i=1}^{N}$, con $N = 35\,007$ imágenes de 399 pacientes, divididas \textbf{por paciente} en train (24\,501 imágenes, 279 pacientes), validación (5\,253, 60) y test (5\,253, 60).

\textbf{Variables de entrada (X).} Cada imagen es una matriz de intensidades. El recorte (escala de grises) se redimensiona a $224 \times 224$ píxeles y se repite en los 3 canales que esperan las redes preentrenadas. Los valores se llevan a $[0, 1]$. La ResNet-50 (y, a través de ella, los modelos 3 y 4) normaliza además cada canal $c$ con la media y la desviación de ImageNet, igual que \arch{StandardScaler}; YOLOv8 usa los valores en $[0, 1]$ sin normalizar:

\begin{center}
\ajusta{x \in [0,1]^{3 \times 224 \times 224}, \qquad \tilde{x}_{c,h,w} = \frac{x_{c,h,w} - \mu_c}{\sigma_c}, \quad \mu = (0{,}485;\ 0{,}456;\ 0{,}406),\ \sigma = (0{,}229;\ 0{,}224;\ 0{,}225)}
\end{center}

En un modelo tabular, X son columnas como "horas de estudio" o "asistencia". Aquí X son \textbf{150\,528 valores de píxel}. Por eso los modelos 1 y 2 aprenden primero a extraer características, y los modelos 3 y 4 reciben esas características ya extraídas.

\textbf{Variable objetivo (y).} Clasificación \textbf{multiclase} con $K = 3$ clases nominales (sin orden):

\begin{center}
\ajusta{y \in \lbrace 0, 1, 2\rbrace  = \lbrace \text{Degeneración},\ \text{Normal},\ \text{Protrusión}\rbrace , \qquad \mathbf{y} = \text{one-hot}(y) \in \lbrace 0,1\rbrace ^3}
\end{center}

\textbf{Qué se busca.} Una función con parámetros $\theta$ que asigne a cada imagen una probabilidad para cada clase, y luego la clase más probable:

\begin{center}
\ajusta{f_\theta : \mathbb{R}^{3 \times 224 \times 224} \to \Delta^{2} = \Big\lbrace  p \in [0,1]^3 : \textstyle\sum_k p_k = 1 \Big\rbrace , \qquad \hat{y} = \arg\max_k\ p_k}
\end{center}

\begin{table}[htbp]
\centering\small
\begin{tabular}{p{0.313\linewidth}p{0.313\linewidth}p{0.313\linewidth}}
\toprule
Modelo & Qué hace con X & Parámetros aprendidos \\
\midrule
1. YOLOv8s-cls & Red convolucional completa, ajustada de extremo a extremo & 5\,084\,579 pesos \\
2. ResNet-50 & Red convolucional residual, ajustada de extremo a extremo & 23\,514\,179 pesos (225\,344 congelados) \\
3. SVM (RBF) & Clasificador de máximo margen sobre 256 características & 1\,684 vectores de soporte y sus coeficientes $\alpha_i$ \\
4. Random Forest & 300 árboles de decisión sobre las mismas 256 características & Las divisiones (variable, umbral) de cada nodo \\
\bottomrule
\end{tabular}
\end{table}



\subsection{Modelos 1 y 2: redes neuronales convolucionales}
\label{mm:2}

\subsubsection{Capas}
\label{mm:2.1}
\textbf{Convolución.} Un filtro $W$ de tamaño $k \times k$ recorre la imagen y calcula, en cada posición, una suma ponderada de los píxeles vecinos. Detecta bordes, texturas y, en capas profundas, formas como el contorno posterior del disco:

\[
(W * x)_{o,i,j} = \sum_{c} \sum_{u=0}^{k-1} \sum_{v=0}^{k-1} W_{o,c,u,v}\ x_{c,\ i+u,\ j+v} + b_o
\]

\textbf{Normalización por lotes (BatchNorm).} Estandariza cada canal con la media $\mu_B$ y la varianza $\sigma_B^2$ del lote, y luego reescala con dos parámetros aprendidos $\gamma$ y $\beta$:

\[
\hat{h} = \gamma\thinspace  \frac{h - \mu_B}{\sqrt{\sigma_B^2 + \epsilon}} + \beta
\]

\textbf{Activación} (no linealidad, Figura~\ref{mm:fig1}): la ResNet-50 usa $\text{ReLU}(z) = \max(0, z)$ y YOLOv8 usa $\text{SiLU}(z) = z \cdot \sigma(z)$.

\textbf{Bloque residual (ResNet-50).} La entrada se suma a la salida de tres convoluciones ($1{\times}1$, $3{\times}3$, $1{\times}1$). Así la red solo aprende la \emph{corrección} $\mathcal{F}$, y el gradiente fluye bien por 50 capas. La ResNet-50 apila 16 bloques en 4 etapas de 3, 4, 6 y 3 bloques:

\[
h^{(l+1)} = \text{ReLU}\big(\mathcal{F}(h^{(l)};\ W^{(l)}) + h^{(l)}\big)
\]

\textbf{YOLOv8s-cls.} Columna de convoluciones con paso 2 y bloques C2f (convoluciones con conexiones parciales entre etapas, la misma idea de reutilizar la entrada). Termina en una convolución $1{\times}1$ a 1\,280 canales.

\textbf{Promedio global y capa final.} El mapa de características final se promedia por canal. Queda un vector $\phi(x)$ (2\,048 valores en la ResNet-50, 1\,280 en YOLOv8), que pasa por \emph{dropout} ($p = 0{,}3$: durante el entrenamiento se apaga al azar el 30\,\% de los valores) y por una capa lineal:

\[
\phi(x) = \frac{1}{HW}\sum_{i,j} h^{(L)}_{:,i,j}, \qquad z = W\thinspace \phi(x) + b \in \mathbb{R}^3
\]

En la ResNet-50 esta capa tiene $2\,048 \times 3 + 3 = 6\,147$ parámetros.

\subsubsection{De puntajes a probabilidades: softmax}
\label{mm:2.2}
La regresión logística de la semana 9 usa la sigmoide para 2 clases. Con 3 clases se usa su generalización, la \textbf{softmax}, que convierte los 3 puntajes $z_k$ (\emph{logits}) en probabilidades positivas que suman 1:

\[
p_k = P(y = k \mid x) = \frac{e^{z_k}}{\sum_{j=1}^{3} e^{z_j}}
\]

Ejemplo (Figura~\ref{mm:fig1}): los logits $(2{,}0;\ 0{,}5;\ -1{,}0)$ dan las probabilidades $(0{,}786;\ 0{,}175;\ 0{,}039)$, y la predicción es la clase 0. Con $K = 2$ la softmax se reduce exactamente a la sigmoide.

\begin{figure}[htbp]
\centering
\includegraphics[width=0.9\linewidth]{01_activaciones_softmax.png}
\caption{Activaciones y softmax.}
\label{mm:fig1}
\end{figure}


\subsubsection{Función de costo: entropía cruzada}
\label{mm:2.3}
Como en la regresión logística, el error de una imagen se mide con la \textbf{entropía cruzada} (log loss): cuánto se aleja la distribución predicha $p$ de la distribución objetivo $q$. El costo $J(\theta)$ promedia esa pérdida sobre las imágenes:

\begin{center}
\ajusta{\ell(x_i, y_i) = -\sum_{k=1}^{3} q_{ik} \log p_{ik}, \qquad J(\theta) = \frac{1}{N}\sum_{i=1}^{N} \ell(x_i, y_i) + \text{regularización}}
\end{center}
\begin{itemize}
  \item \textbf{YOLOv8s-cls:} $q = \mathbf{y}$ (one-hot), así que $\ell = -\log p_{y}$. Si la red da $p_y = 0{,}01$ a la clase correcta, la pérdida es 4,6; si da 0,99, es 0,01. Penaliza fuerte la \textbf{seguridad equivocada} (Figura~\ref{mm:fig2}).
  \item \textbf{ResNet-50, con suavizado de etiquetas} ($\varepsilon = 0{,}1$): en lugar de exigir probabilidad 1 a la clase real, el objetivo reparte un poco de masa entre las demás:
\end{itemize}

\begin{center}
\ajusta{q_k = (1-\varepsilon)\thinspace \mathbf{y}_k + \frac{\varepsilon}{K} \quad\Rightarrow\quad q = (0{,}933;\ 0{,}033;\ 0{,}033)}
\end{center}

Así la red no puede volverse arbitrariamente segura, lo que reduce el sobreajuste. El costo mínimo posible ya no es 0 sino la entropía de $q$, $H(q) = 0{,}291$. Ese es el piso al que se acerca la pérdida de train (0,305 en la época 6, Figura~\ref{mm:fig2}).

\textbf{¿Por qué no MSE para entrenar?} El MSE de la semana 9 es la función de costo natural cuando $y$ es un número real (regresión). Para probabilidades de clase, la entropía cruzada da gradientes más grandes cuando el error es grande. Además, es la función que se obtiene al maximizar la verosimilitud del modelo. El MSE sí se usa para \textbf{evaluar}: se reporta como puntaje de Brier (Sección~\ref{mm:4.2}) en el registro de desempeño (Sección~\ref{mm:4.4}).

\begin{figure}[htbp]
\centering
\includegraphics[width=0.9\linewidth]{02_funcion_costo.png}
\caption{Función de costo.}
\label{mm:fig2}
\end{figure}


\subsubsection{Optimización: descenso del gradiente con AdamW}
\label{mm:2.4}
\textbf{Descenso del gradiente por lotes.} En cada paso se toma un lote de $B = 64$ imágenes y se mueve cada parámetro en la dirección opuesta al gradiente del costo, como la regla $\theta_j := \theta_j - \alpha\thinspace  \partial J / \partial \theta_j$ de la semana 9:

\[
g_t = \nabla_\theta \frac{1}{B}\sum_{i \in \text{lote}} \ell(x_i, y_i)
\]

El gradiente de cada peso se calcula con la regla de la cadena desde la salida hacia la entrada (retropropagación).

\textbf{AdamW.} Mantiene un promedio móvil del gradiente ($m_t$, el "impulso") y de su cuadrado ($v_t$, la escala). Con eso cada parámetro avanza con un paso adaptado a su propia escala. El término $\lambda\thinspace \theta$ (\emph{weight decay}) encoge los pesos hacia 0, y es la regularización L2:

\begin{center}
\ajusta{m_t = \beta_1 m_{t-1} + (1-\beta_1)\thinspace  g_t, \qquad v_t = \beta_2 v_{t-1} + (1-\beta_2)\thinspace  g_t^2}
\end{center}

\begin{center}
\ajusta{\hat{m}_t = \frac{m_t}{1-\beta_1^t}, \quad \hat{v}_t = \frac{v_t}{1-\beta_2^t}, \qquad \theta_t = \theta_{t-1} - \alpha_t \left( \frac{\hat{m}_t}{\sqrt{\hat{v}_t} + \epsilon} + \lambda\thinspace  \theta_{t-1} \right)}
\end{center}

Valores: $\beta_1 = 0{,}9$ en la ResNet-50 y $0{,}937$ en YOLOv8 (valor por defecto de ultralytics), $\beta_2 = 0{,}999$, $\epsilon = 10^{-8}$, $\lambda = 0{,}01$. Ultralytics no aplica $\lambda$ a los sesgos ni a las capas de normalización.

\textbf{Tasa de aprendizaje en coseno} (Figura~\ref{mm:fig3}). El tamaño del paso $\alpha_t$ empieza en $\alpha_0$ y baja suavemente hasta $\alpha_{\min}$ en $T$ épocas: pasos grandes al principio, finos al final.

\[
\alpha_t = \alpha_{\min} + \tfrac{1}{2}(\alpha_0 - \alpha_{\min})\left(1 + \cos\frac{\pi t}{T}\right)
\]

\begin{table}[htbp]
\centering\small
\begin{tabular}{p{0.313\linewidth}p{0.313\linewidth}p{0.313\linewidth}}
\toprule
 & YOLOv8s-cls & ResNet-50 \\
\midrule
$\alpha_0$ / $\alpha_{\min}$ & $5\cdot10^{-4}$ / $5\cdot10^{-5}$ (1 época de calentamiento) & $10^{-4}$ / 0 \\
Épocas máximas $T$ / paciencia & 15 / 4 & 10 / 4 \\
Capas congeladas (no se actualizan) & primeros 4 módulos & \arch{conv1}, \arch{bn1}, \arch{layer1} (225\,344 pesos) \\
Aumento en línea & recorte aleatorio con el 100\,\% del área (solo varía la proporción, 3:4 a 4:3) y brillo $\pm$5\,\% & rotación $\pm$6$^{\circ}$, desplazamiento 4\,\%, zoom 0,95–1,05, brillo y contraste $\pm$10\,\% \\
Suavizado de etiquetas & no & $\varepsilon = 0{,}1$ \\
Criterio para elegir la época & exactitud top-1 en validación & F1 macro en validación \\
Época elegida / épocas corridas & 2 / 6 & 2 / 6 \\
\bottomrule
\end{tabular}
\end{table}


\textbf{Ajuste fino (\emph{transfer learning}).} Ninguna red parte de cero: $\theta$ se inicializa con los pesos aprendidos en ImageNet (1,28 millones de fotografías), que ya detectan bordes y texturas. Solo la capa final empieza al azar.

\begin{figure}[htbp]
\centering
\includegraphics[width=0.9\linewidth]{03_tasa_aprendizaje.png}
\caption{Tasa de aprendizaje.}
\label{mm:fig3}
\end{figure}


\subsubsection{Diagnóstico: sobreajuste y parada temprana}
\label{mm:2.5}
La Figura~\ref{mm:fig4} muestra el patrón de las dos redes: la pérdida de train sigue bajando, mientras la de validación sube desde la época 2. El modelo empieza a memorizar las imágenes de train, que además incluyen copias aumentadas de las mismas imágenes. Por eso se guarda la época con mejor desempeño en validación y el entrenamiento se detiene tras 4 épocas sin mejora (\textbf{parada temprana}):

\[
e^{*} = \arg\max_{e}\ \text{F1}_{\text{macro}}\big(\text{validación};\ \theta^{(e)}\big)
\]

\begin{figure}[htbp]
\centering
\includegraphics[width=0.9\linewidth]{04_curvas_aprendizaje.png}
\caption{Curvas de aprendizaje.}
\label{mm:fig4}
\end{figure}



\subsection{Modelos 3 y 4: características profundas + clasificador clásico}
\label{mm:3}

Los modelos 3 y 4 no trabajan sobre píxeles. Cada imagen se representa con el vector $\phi(x) \in \mathbb{R}^{2048}$ que produce la ResNet-50 del modelo 2 justo antes de su capa final (Sección~\ref{mm:2.1}). Después se reduce ese vector a 256 dimensiones y se clasifica con un SVM o con un Random Forest.

\textbf{Dependencia declarada:} los modelos 3 y 4 dependen del modelo 2. Con un extractor de ImageNet sin ajustar rinden mucho menos (Sección~\ref{mm:3.4}).

\subsubsection{Estandarización y PCA}
\label{mm:3.1}
\textbf{Estandarización.} Cada una de las 2\,048 características se centra y se escala con la media y la desviación \textbf{calculadas solo con train}. Calcularlas con todo el dataset filtraría información del test (\emph{data leakage}, semanas 6 y 9):

\[
s_j = \frac{\phi_j(x) - \mu_j^{\text{train}}}{\sigma_j^{\text{train}}}
\]

\textbf{Análisis de componentes principales (PCA).} Se busca la base de direcciones ortogonales en la que los datos de train varían más. Son los vectores propios de la matriz de covarianza, y cada imagen se proyecta sobre los $k = 256$ primeros:

\begin{center}
\ajusta{\Sigma = \frac{1}{n-1}\thinspace  S^\top S = V \Lambda V^\top, \qquad u = V_k^\top s \in \mathbb{R}^{256}, \qquad \text{varianza explicada} = \frac{\sum_{j=1}^{k}\lambda_j}{\sum_{j=1}^{2048}\lambda_j}}
\end{center}

Con las características ajustadas, \textbf{256 componentes explican el 80,0\,\% de la varianza} (las 2 primeras ya el 38\,\%). Con las de ImageNet sin ajustar, solo el 58,7\,\% (Figura~\ref{mm:fig5}). Las características ajustadas están más concentradas en pocas direcciones, las que separan las clases.

\begin{figure}[htbp]
\centering
\includegraphics[width=0.9\linewidth]{05_pca_varianza.png}
\caption{Varianza explicada por el PCA.}
\label{mm:fig5}
\end{figure}


\subsubsection{Modelo 3: máquina de vectores de soporte (SVM) con kernel RBF}
\label{mm:3.2}
\textbf{Idea.} Entre todas las fronteras que separan dos clases, el SVM elige la de \textbf{máximo margen}: la que queda más lejos de los ejemplos más cercanos de cada lado. Esos ejemplos se llaman vectores de soporte. Para dos clases, $y_i \in \lbrace -1, +1\rbrace $, con margen blando:

\begin{center}
\ajusta{\min_{w,\thinspace b,\thinspace \xi}\ \ \frac{1}{2}\lVert w \rVert^2 + C \sum_{i=1}^{n} \xi_i \quad \text{sujeto a} \quad y_i\big(w^\top \varphi(u_i) + b\big) \ge 1 - \xi_i,\quad \xi_i \ge 0}
\end{center}
\begin{itemize}
  \item $\frac{1}{2}\lVert w\rVert^2$: minimizarlo maximiza el ancho del margen, $2 / \lVert w \rVert$.
  \item $\xi_i$: holgura que permite que algunas imágenes queden dentro del margen o mal clasificadas.
  \item $C$: equilibrio entre margen ancho y errores en train. Un $C$ pequeño tolera más errores y da una frontera más suave, menos sobreajustada.
\end{itemize}

Es equivalente a minimizar la \textbf{pérdida bisagra} (\emph{hinge}): $\frac{1}{2}\lVert w\rVert^2 + C\sum_i \max\big(0,\ 1 - y_i f(u_i)\big)$.

\textbf{Kernel RBF.} Las clases no se separan con un hiperplano en el espacio original. El SVM resuelve el problema dual, en el que los datos solo aparecen como similitudes $K(u_i, u_j)$; así trabaja implícitamente en un espacio de dimensión infinita (el "truco del kernel"):

\begin{center}
\ajusta{K(u, v) = \exp\big(-\gamma\thinspace \lVert u - v \rVert^2\big), \qquad \gamma = \frac{1}{d \cdot \mathrm{Var}(U)} = 0{,}00061\ \ (\texttt{gamma="scale"},\ d = 256)}
\end{center}

\begin{center}
\ajusta{\max_{\alpha}\ \sum_i \alpha_i - \frac{1}{2}\sum_i\sum_j \alpha_i\alpha_j\thinspace  y_i y_j\thinspace  K(u_i, u_j) \quad \text{s.a.}\ \ 0 \le \alpha_i \le C,\ \ \sum_i \alpha_i y_i = 0}
\end{center}

\[
f(u) = \sum_{i \in \text{SV}} \alpha_i\thinspace  y_i\thinspace  K(u_i, u) + b
\]

Solo los vectores de soporte ($\alpha_i > 0$) intervienen en la predicción. Con $C = 0{,}1$, elegido en validación (Figura~\ref{mm:fig8}), el modelo usa \textbf{1\,684 vectores de soporte de 24\,501 imágenes de train (6,9\,\%)}: 320 de Degeneración, 660 de Normal y 704 de Protrusión.

La Figura~\ref{mm:fig6}, a la derecha, muestra cómo decae la similitud con la distancia. Con el $\gamma$ usado, dos imágenes a la distancia mediana (59) tienen similitud $\approx$ 0,12. Con $\gamma \times 10$ todas las imágenes serían "lejanas" (sobreajuste); con $\gamma / 10$, todas "parecidas" (subajuste).

\textbf{Tres clases.} scikit-learn entrena un SVM por cada par de clases (\emph{uno contra uno}: 3 clasificadores) y combina sus votos en un puntaje por clase; la predicción es la clase con mayor puntaje. Estos puntajes \textbf{no son probabilidades}; por eso el SVM no tiene log loss y la app los muestra normalizados con softmax.

La Figura~\ref{mm:fig6}, a la izquierda, es una \textbf{ilustración en 2 dimensiones} (las 2 primeras componentes, 1\,200 originales de train). En train las clases aparecen muy separadas porque la ResNet ya vio esas imágenes; en validación y test se mezclan más, y por eso la exactitud allí baja a ~88\,\%.

\begin{figure}[htbp]
\centering
\includegraphics[width=0.9\linewidth]{06_svm_rbf.png}
\caption{SVM con kernel RBF.}
\label{mm:fig6}
\end{figure}


\subsubsection{Modelo 4: Random Forest}
\label{mm:3.3}
\textbf{Árbol de decisión.} Hace preguntas del tipo "¿$u_j \le t$?" y divide las imágenes en dos grupos. En cada nodo elige la variable $j$ y el umbral $t$ que más reducen la \textbf{impureza de Gini} (semana 11):

\begin{center}
\ajusta{G(\text{nodo}) = 1 - \sum_{k=1}^{3} p_k^2, \qquad \Delta G = G(\text{padre}) - \frac{n_L}{n}\thinspace G(L) - \frac{n_R}{n}\thinspace G(R)}
\end{center}

$G = 0$ si el nodo tiene una sola clase (puro). El máximo es $1 - 1/K$ si las clases están mezcladas en partes iguales: 0,5 con 2 clases, 0,667 con 3. Ejemplo de la clase: 6 compran y 4 no ($G = 1 - 0{,}6^2 - 0{,}4^2 = 0{,}48$). Tras dividir en un grupo (5, 0) con $G = 0$ y otro (1, 4) con $G = 0{,}32$, la impureza ponderada baja a $0{,}5 \cdot 0 + 0{,}5 \cdot 0{,}32 = 0{,}16$.

\textbf{Bosque.} Un árbol solo sobreajusta. Se entrenan $T = 300$ árboles distintos entre sí por dos fuentes de azar:
\begin{enumerate}
  \item \textbf{Bootstrap:} cada árbol se entrena con una muestra de train tomada con reemplazo, del mismo tamaño.
  \item \textbf{Variables al azar:} en cada división solo se consideran $m = \sqrt{256} = 16$ de las 256 variables.
\end{enumerate}

La predicción promedia las probabilidades de los 300 árboles (votación suave). La probabilidad de cada árbol, $p_t(k \mid u)$, es la proporción de la clase $k$ en la hoja donde cae $u$:

\[
P(y = k \mid u) = \frac{1}{T}\sum_{t=1}^{T} p_t(k \mid u), \qquad \hat{y} = \arg\max_k P(y = k \mid u)
\]

Al promediar árboles diversos baja la varianza (el sobreajuste) sin subir mucho el sesgo. Hiperparámetro elegido en validación (Figura~\ref{mm:fig8}): profundidad máxima 20. Los 300 árboles llegan a esa profundidad y tienen 779 hojas en promedio.

La \textbf{importancia} de cada variable es la disminución total de Gini que aporta, promediada sobre los árboles. Las 10 primeras componentes del PCA concentran el 70\,\% de la importancia (Figura~\ref{mm:fig7}).

\begin{figure}[htbp]
\centering
\includegraphics[width=0.9\linewidth]{07_random_forest.png}
\caption{Random Forest.}
\label{mm:fig7}
\end{figure}


\subsubsection{Extras del curso y comparación con ImageNet}
\label{mm:3.4}
Sobre los mismos 256 componentes se entrenaron, como línea base del curso:
\begin{itemize}
  \item \textbf{Regresión logística multinomial} (semana 9): $p = \text{softmax}(W u + b)$, con regularización L2 de fuerza $1/C$ y $C = 0{,}01$.
  \item \textbf{KNN} (semana 10): predice por mayoría entre los $K = 31$ vecinos más cercanos según la distancia euclídea $d(u, v) = \sqrt{\sum_j (u_j - v_j)^2}$. La escala importa, y por eso se estandariza antes.
\end{itemize}

Con características de ImageNet sin ajustar, el SVM baja de 88,4\,\% a 77,6\,\% y el Random Forest de 88,4\,\% a 75,3\,\% de exactitud en el test de originales. Lo que hace funcionar a los modelos 3 y 4 es la representación aprendida por el modelo 2.

\begin{figure}[htbp]
\centering
\includegraphics[width=0.9\linewidth]{08_busqueda_hiperparametros.png}
\caption{Búsqueda de hiperparámetros.}
\label{mm:fig8}
\end{figure}


\subsubsection{Estructura del sistema: arquitectura híbrida}
\label{mm:3.5}

El sistema tiene \textbf{dos caminos}. YOLOv8s-cls va por el suyo, independiente. La ResNet-50 es la \textbf{base compartida} de los otros tres modelos: sus características $\phi(x)$ alimentan su propia capa final (modelo 2) y, después de estandarizarlas y reducirlas, al SVM (modelo 3) y al Random Forest (modelo 4). Unir una red profunda que extrae características con un clasificador clásico que decide es lo que se llama \textbf{arquitectura híbrida} (Figura~\ref{fig:estructura}).


\textbf{Las 4 funciones.} Sea $g_Y$ la red YOLOv8s-cls hasta sus logits, $g_R$ la ResNet-50 sin su capa final y $\phi(x) = g_R(x)$. Cada modelo es una composición de funciones:

\[
p_1(x) = \text{softmax}\big(g_Y(x)\big)
\]

\[
p_2(x) = \text{softmax}\big(W\thinspace \phi(x) + b\big)
\]

\begin{center}
\ajusta{u(x) = V_k^\top\thinspace D^{-1}\big(\phi(x) - \mu\big), \qquad D = \text{diag}(\sigma_1, \dots, \sigma_{2048})}
\end{center}

\begin{center}
\ajusta{p_3(x) = \text{softmax}\big(\text{SVM}(u(x))\big), \qquad p_4(x) = \frac{1}{T}\sum_{t=1}^{T} p_t\big(u(x)\big)}
\end{center}

Las tres últimas dependen de la misma $\phi(x)$. Por eso, en la aplicación la ResNet-50 se ejecuta una sola vez por imagen. La softmax de $p_3$ solo lleva los puntajes del SVM a la misma escala de los demás: no son probabilidades calibradas.

\textbf{Orden de entrenamiento.} Cada etapa usa solo lo aprendido en las anteriores, y nada se ajusta con el test:

\begin{table}[htbp]
\centering\small
\begin{tabular}{p{0.235\linewidth}p{0.235\linewidth}p{0.235\linewidth}p{0.235\linewidth}}
\toprule
Paso & Qué se aprende & Con qué datos & Qué queda fijo \\
\midrule
1 & Pesos de YOLOv8s-cls & train; época elegida con validación & — \\
2 & Pesos de la ResNet-50, incluidos $W$ y $b$ & train; época elegida con validación & — \\
3 & $\mu$, $D$ y $V_k$ (escalador y PCA) & características de train & La ResNet-50 \\
4 & SVM y Random Forest & $u$ de train; hiperparámetros elegidos con validación & ResNet-50, escalador y PCA \\
\bottomrule
\end{tabular}
\end{table}


\textbf{Por qué es híbrida y no un ensamble.} En esta estructura un modelo le pasa sus \textbf{características} a otro, y cada uno sigue dando su propia predicción. Un ensamble, en cambio, combina las \textbf{respuestas finales} de varios modelos. Esa combinación es la sección siguiente.

\subsubsection{Ensamble por promedio}
\label{mm:3.6}

\textbf{Idea.} Es la misma de la semana 11: consultar a varios modelos reduce la dependencia de una sola decisión. Cada modelo entrega 3 puntajes para la imagen; se promedian por clase y gana la clase con el mayor promedio (\textbf{votación suave}). Con $M = 4$ modelos:

\begin{center}
\ajusta{P_{\text{ens}}(k \mid x) = \frac{1}{M}\sum_{m=1}^{M} p_m(k \mid x), \qquad \hat{y}_{\text{ens}} = \arg\max_k\ P_{\text{ens}}(k \mid x)}
\end{center}

El ensamble \textbf{no se entrena}: todos los modelos pesan igual ($1/4$) y no tiene parámetros que aprender. Se calcula después, con las salidas de los 4 modelos ya entrenados. Por eso es un \textbf{extra} y no un quinto modelo.

\textbf{Ejemplo real} (imagen de test: paciente 0319, L1-L2, etiqueta real Degeneración):

\begin{table}[htbp]
\centering\small
\begin{tabular}{lllll}
\toprule
Modelo & Degeneración & Normal & Protrusión & Predicción \\
\midrule
1. YOLOv8s-cls & 0,109 & 0,871 & 0,020 & Normal (falla) \\
2. ResNet-50 & 0,936 & 0,025 & 0,039 & Degeneración \\
3. SVM (con softmax) & 0,746 & 0,064 & 0,190 & Degeneración \\
4. Random Forest & 0,866 & 0,050 & 0,083 & Degeneración \\
\textbf{Ensamble (promedio)} & \textbf{0,664} & 0,253 & 0,083 & \textbf{Degeneración} \\
\bottomrule
\end{tabular}
\end{table}


Por ejemplo, $P_{\text{ens}}(\text{Degeneración}) = (0{,}109 + 0{,}936 + 0{,}746 + 0{,}866) / 4 = 0{,}664$. Los otros tres modelos corrigen el error de YOLO.

\textbf{Por qué puede ayudar, y su límite.} Si cada modelo tiene una varianza de error $\sigma^2$ y la correlación entre modelos es $\rho$, la varianza del promedio de $M$ modelos es:

\begin{center}
\ajusta{\text{Var}\left(\frac{1}{M}\sum_{m=1}^{M} p_m\right) = \rho\thinspace \sigma^2 + \frac{1-\rho}{M}\thinspace \sigma^2}
\end{center}

Si los modelos fueran independientes ($\rho = 0$), la varianza se dividiría entre 4. Si fueran idénticos ($\rho = 1$), promediar no serviría de nada. Aquí los modelos 2, 3 y 4 comparten $\phi(x)$, así que se espera una correlación alta y una mejora limitada.

\textbf{Resultados} (test de originales; detalle en las Secciones~\ref{mm:4.3} y~\ref{mm:4.4}):

\begin{table}[htbp]
\centering\small
\begin{tabular}{lll}
\toprule
Métrica & Mejor modelo individual & Ensamble \\
\midrule
Exactitud & 0,889 (ResNet-50) & 0,889 \\
F1 macro & 0,769 (ResNet-50) & 0,745 \\
ROC-AUC macro & 0,942 (Random Forest) & \textbf{0,947} \\
MSE (Brier) & 0,173 (ResNet-50) & \textbf{0,164} \\
\bottomrule
\end{tabular}
\end{table}


El ensamble empata en exactitud y da las probabilidades más confiables (mejor ROC-AUC y MSE), pero detecta menos casos de Degeneración: 17 de 42, frente a 21 de la ResNet-50. En validación, el criterio acordado, su F1 macro es 0,822 frente a 0,824 de la ResNet-50. Por eso \textbf{la ResNet-50 sigue siendo el modelo principal} y el ensamble se reporta como resultado adicional.


\subsection{Evaluación}
\label{mm:4}

\subsubsection{Protocolo}
\label{mm:4.1}
\begin{itemize}
  \item Los parámetros $\theta$ se aprenden solo con \textbf{train}.
  \item Los hiperparámetros ($C$, profundidad, $K$, época) se eligen con \textbf{validación}, por F1 macro.
  \item El \textbf{test} se evalúa una sola vez, al final. Ningún paciente aparece en dos particiones, y cada copia aumentada queda en la partición de su original.
  \item Se reporta el test de \textbf{originales} (768 imágenes en proporción natural: 42 Degeneración, 164 Normal, 562 Protrusión) y el \textbf{test completo} (5\,253, balanceado con copias).
\end{itemize}

\subsubsection{Matriz de confusión y métricas}
\label{mm:4.2}
En clase se vieron las fórmulas para \textbf{2 clases} (positiva y negativa):

\begin{center}
\ajusta{\text{Accuracy} = \frac{TP + TN}{TP + TN + FP + FN}, \qquad \text{Precision} = \frac{TP}{TP + FP}, \qquad \text{Recall} = \frac{TP}{TP + FN}, \qquad F1 = 2\cdot\frac{\text{Precision}\cdot\text{Recall}}{\text{Precision} + \text{Recall}}}
\end{center}

Con \textbf{3 clases} se aplican las mismas fórmulas tres veces. Cada clase $k$ se trata como la "positiva" y las otras dos como la "negativa" (uno contra el resto): TP = predijo $k$ y era $k$; FP = predijo $k$ y no era $k$; FN = era $k$ y predijo otra; TN = el resto. Luego se promedian las tres (promedio \emph{macro}). La exactitud no necesita promedio: es el total de aciertos (la diagonal de la matriz) sobre el total de imágenes.

\begin{center}
\ajusta{\text{Exactitud} = \frac{\sum_k \text{TP}_k}{N}, \qquad \text{Precisión}_k = \frac{\text{TP}_k}{\text{TP}_k + \text{FP}_k}, \qquad \text{Recall}_k = \frac{\text{TP}_k}{\text{TP}_k + \text{FN}_k}}
\end{center}

\begin{center}
\ajusta{F1_k = 2\thinspace \frac{\text{Precisión}_k \cdot \text{Recall}_k}{\text{Precisión}_k + \text{Recall}_k}, \qquad F1_{\text{macro}} = \frac{1}{K}\sum_k F1_k, \qquad \text{Exactitud balanceada} = \frac{1}{K}\sum_k \text{Recall}_k}
\end{center}

\textbf{Ejemplo con el modelo 2 en el test de originales} (filas = real, columnas = predicha):

\begin{table}[htbp]
\centering\small
\begin{tabular}{llll}
\toprule
Real \textbackslash{} Predicha & Degeneración & Normal & Protrusión \\
\midrule
\textbf{Degeneración} (42) & \textbf{21} & 6 & 15 \\
\textbf{Normal} (164) & 5 & \textbf{133} & 26 \\
\textbf{Protrusión} (562) & 6 & 27 & \textbf{529} \\
\bottomrule
\end{tabular}
\end{table}


\begin{table}[htbp]
\centering\small
\begin{tabular}{lllllll}
\toprule
Clase & TP & FP & FN & Precisión & Recall & F1 \\
\midrule
Degeneración & 21 & 11 & 21 & 0,656 & 0,500 & 0,568 \\
Normal & 133 & 33 & 31 & 0,801 & 0,811 & 0,806 \\
Protrusión & 529 & 41 & 33 & 0,928 & 0,941 & 0,935 \\
\bottomrule
\end{tabular}
\end{table}

\begin{itemize}
  \item \textbf{Exactitud} = (21 + 133 + 529) / 768 = 683 / 768 = \textbf{0,889}.
  \item \textbf{F1 macro} = (0,568 + 0,806 + 0,935) / 3 = \textbf{0,769}.
\end{itemize}

\textbf{Por qué la exactitud sola engaña.} Un modelo que siempre dijera "Protrusión" acertaría 562 / 768 = \textbf{73,2\,\%} sin detectar ni un caso de Degeneración. Es el mismo ejemplo del fraude visto en clase. El F1 macro pesa igual las tres clases y deja ver que Degeneración se detecta solo en la mitad de los casos.

\textbf{Intervalo de confianza.} Con 768 imágenes, la exactitud de 88,9\,\% tiene un intervalo de Wilson al 95\,\% de \textbf{[86,5\,\%; 91,0\,\%]}:

\begin{center}
\ajusta{\frac{\hat{p} + \frac{z^2}{2n} \pm z\sqrt{\frac{\hat{p}(1-\hat{p})}{n} + \frac{z^2}{4n^2}}}{1 + \frac{z^2}{n}}, \qquad z = 1{,}96}
\end{center}

Las diferencias entre los 4 modelos (87,8\,\% a 88,9\,\%) están dentro de ese margen.

\textbf{ROC-AUC.} Para cada clase, la curva ROC recorre todos los umbrales $\tau$ sobre $p_k$ y grafica la tasa de verdaderos positivos contra la de falsos positivos. El área bajo la curva es la probabilidad de que una imagen de la clase $k$ reciba más puntaje que una que no lo es. Se promedia sobre las 3 clases (uno contra el resto).

\textbf{Log loss.} Es la entropía cruzada de la Sección~\ref{mm:2.3}, evaluada con las probabilidades sobre el conjunto de prueba.

\textbf{MSE (puntaje de Brier).} Es el MSE de la semana 9 aplicado a la clasificación. El "valor real" de cada imagen es su etiqueta en one-hot (por ejemplo, Protrusión = $(0, 0, 1)$) y la "predicción" son las 3 probabilidades del modelo. Se suman los errores al cuadrado de las 3 clases y se promedia sobre las $N$ imágenes:

\begin{center}
\ajusta{\text{MSE} = \frac{1}{N}\sum_{i=1}^{N}\sum_{k=1}^{3}\left(p_{ik} - y_{ik}\right)^2, \qquad \text{RMSE} = \sqrt{\text{MSE}}}
\end{center}

Vale 0 si el modelo da probabilidad 1 a la clase correcta y como máximo 2. Igual que el MSE de clase, castiga más los errores grandes. Ejemplo con una imagen real de Protrusión:
\begin{itemize}
  \item Si el modelo da $p = (0{,}1;\ 0{,}2;\ 0{,}7)$: $0{,}1^2 + 0{,}2^2 + 0{,}3^2 = 0{,}14$ (acierta, con algo de duda).
  \item Si da $p = (0{,}9;\ 0{,}05;\ 0{,}05)$: $0{,}9^2 + 0{,}05^2 + 0{,}95^2 = 1{,}715$ (falla con mucha seguridad).
\end{itemize}

A diferencia de la exactitud, el MSE también mide si la \textbf{confianza} del modelo es la adecuada. No aplica al SVM, porque sus puntajes no son probabilidades.

\subsubsection{Resultados (test de originales)}
\label{mm:4.3}
Los resultados en los originales de prueba están en la Tabla~\ref{tab:resultados}, con sus matrices de confusión y curvas ROC en las Figuras~\ref{fig:matrices} y~\ref{fig:roc}.

\subsubsection{Registro de desempeño (formato de la clase)}
\label{mm:4.4}

Es la tabla "Registro de desempeño" de la semana 9, con columnas Train y Test.
\begin{itemize}
  \item \textbf{Train:} las 3\,577 imágenes originales de entrenamiento.
  \item \textbf{Test:} las 768 originales de test.
  \item \textbf{Test completo:} las 5\,253 imágenes de test con copias aumentadas, como referencia adicional.
\end{itemize}

Precision, recall y F1 son promedios \emph{macro} de las 3 clases. \textbf{R$^2$ no aplica}: es una métrica de regresión y este problema es de clasificación (como advierte la diapositiva, no todas las métricas aplican a todos los modelos).

\begin{table}[htbp]
\centering\small
\caption{Registro de desempeño. Entrenamiento: 3\,577 originales. Prueba: 768 originales. Prueba completa: 5\,253 imágenes con copias, balanceada. Precision, recall y F1 son promedios macro.}
\label{tab:registro}
\begin{tabular}{lllll}
\toprule
Modelo & Métrica & Train & Test & Test completo \\
\midrule
\textbf{1. YOLOv8s-cls} & MSE / RMSE (Brier) & 0,024 / 0,154 & 0,196 / 0,442 & 0,403 / 0,635 \\
 & Accuracy & 0,984 & 0,878 & 0,759 \\
 & Precision / Recall & 0,959 / 0,989 & 0,771 / 0,732 & 0,779 / 0,759 \\
 & F1 / ROC-AUC & 0,973 / 0,999 & 0,750 / 0,931 & 0,752 / 0,923 \\
\textbf{2. ResNet-50} & MSE / RMSE (Brier) & 0,047 / 0,217 & 0,173 / 0,416 & 0,349 / 0,591 \\
 & Accuracy & 0,971 & \textbf{0,889} & 0,775 \\
 & Precision / Recall & 0,948 / 0,981 & 0,795 / 0,751 & 0,799 / 0,775 \\
 & F1 / ROC-AUC & 0,964 / 0,998 & 0,769 / 0,919 & 0,768 / 0,915 \\
\textbf{3. SVM (RBF)} & MSE / RMSE (Brier) & no aplica & no aplica & no aplica \\
 & Accuracy & 0,979 & 0,884 & 0,754 \\
 & Precision / Recall & 0,964 / 0,984 & 0,797 / 0,718 & 0,789 / 0,754 \\
 & F1 / ROC-AUC & 0,974 / 0,998 & 0,748 / 0,899 & 0,746 / 0,905 \\
\textbf{4. Random Forest} & MSE / RMSE (Brier) & 0,010 / 0,098 & 0,174 / 0,416 & 0,370 / 0,608 \\
 & Accuracy & 0,998 & 0,884 & 0,743 \\
 & Precision / Recall & 0,994 / 0,999 & 0,810 / 0,710 & 0,784 / 0,743 \\
 & F1 / ROC-AUC & 0,996 / 1,000 & 0,746 / 0,942 & 0,731 / 0,920 \\
\textbf{Extra: ensamble} & MSE / RMSE (Brier) & 0,031 / 0,175 & 0,164 / 0,405 & 0,334 / 0,578 \\
 & Accuracy & 0,989 & 0,889 & 0,762 \\
 & Precision / Recall & 0,978 / 0,992 & 0,782 / 0,721 & 0,794 / 0,762 \\
 & F1 / ROC-AUC & 0,985 / 1,000 & 0,745 / 0,947 & 0,752 / 0,937 \\
\bottomrule
\end{tabular}
\end{table}


\textbf{Lectura:}
\begin{itemize}
  \item La distancia entre Train y Test es la que se vio en las curvas de aprendizaje (Sección~\ref{mm:2.5}): \textbf{sobreajuste}. En train los modelos casi no se equivocan (MSE de 0,01 a 0,05); en test el MSE sube a ~0,17–0,20.
  \item Por MSE, el ensamble es el mejor en test (0,164), seguido de la ResNet-50 y el Random Forest (0,173 y 0,174).
  \item Todas las cifras salen de \arch{docs/resultados/comparacion/tabla_metricas_larga.csv}, que \arch{scripts/entrenamiento/05_comparar_modelos.py} calcula a partir de los puntajes guardados de cada modelo.
\end{itemize}


\subsection{Supuestos y limitaciones del modelo}
\label{mm:5}
\begin{itemize}
  \item \textbf{Independencia:} se supone que los pacientes son independientes entre sí. Las imágenes de un mismo paciente no lo son, por eso se divide por paciente. Las copias aumentadas no son datos nuevos: el tamaño efectivo es de 399 pacientes, no de 35\,007 imágenes.
  \item \textbf{Etiquetas con ruido:} $y$ la asignaron estudiantes, sin radiólogo y sin corte axial. El modelo aprende a reproducir esas etiquetas, no un diagnóstico.
  \item \textbf{Desbalance:} el entrenamiento ve las clases balanceadas (por copias), pero en la proporción natural el 73\,\% es Protrusión. Se probó corregir por esa proporción multiplicando $p_k$ por la frecuencia de cada clase; no mejoró la exactitud y hundió el recall de Degeneración, así que se descartó (\arch{docs/resultados/experimentos_validacion.md}).
  \item \textbf{Sobreajuste:} las redes llegan a ~97–99\,\% en train y ~89\,\% en test de originales. La regularización (dropout, weight decay, suavizado, aumento) y la parada temprana lo reducen, pero no lo eliminan.
  \item \textbf{Dependencia de los modelos 3 y 4} respecto del modelo 2: comparten la representación $\phi(x)$, así que no son cuatro opiniones independientes. Eso limita lo que puede ganar el ensamble (Sección~\ref{mm:3.6}).
  \item \textbf{Herramienta académica, no diagnóstica.}
\end{itemize}

% Generada a partir de docs/modelo_matematico/README.md
% Paquetes: amsmath, amssymb, booktabs, graphicx
% =====================================================================
\section{Modelo matemático}
\label{sec:matematico}

Esta sección explica, con sus ecuaciones, cómo cada uno de los cuatro modelos convierte una imagen en una predicción, cómo aprende sus parámetros y cómo se evalúa. Sigue la secuencia vista en clase: datos, variables, modelo, predicción, error, optimización y evaluación. Todas las cifras salen de los modelos entrenados.


\subsection{Planteamiento: X e y}
\label{mm:1}

\textbf{Datos.} $\mathcal{D} = \lbrace (x_i, y_i)\rbrace _{i=1}^{N}$, con $N = 35\,007$ imágenes de 399 pacientes, divididas \textbf{por paciente} en train (24\,501 imágenes, 279 pacientes), validación (5\,253, 60) y test (5\,253, 60).

\textbf{Variables de entrada (X).} Cada imagen es una matriz de intensidades. El recorte (escala de grises) se redimensiona a $224 \times 224$ píxeles y se repite en los 3 canales que esperan las redes preentrenadas. Los valores se llevan a $[0, 1]$. La ResNet-50 (y, a través de ella, los modelos 3 y 4) normaliza además cada canal $c$ con la media y la desviación de ImageNet, igual que \arch{StandardScaler}; YOLOv8 usa los valores en $[0, 1]$ sin normalizar:

\begin{center}
\ajusta{x \in [0,1]^{3 \times 224 \times 224}, \qquad \tilde{x}_{c,h,w} = \frac{x_{c,h,w} - \mu_c}{\sigma_c}, \quad \mu = (0{,}485;\ 0{,}456;\ 0{,}406),\ \sigma = (0{,}229;\ 0{,}224;\ 0{,}225)}
\end{center}

En un modelo tabular, X son columnas como "horas de estudio" o "asistencia". Aquí X son \textbf{150\,528 valores de píxel}. Por eso los modelos 1 y 2 aprenden primero a extraer características, y los modelos 3 y 4 reciben esas características ya extraídas.

\textbf{Variable objetivo (y).} Clasificación \textbf{multiclase} con $K = 3$ clases nominales (sin orden):

\begin{center}
\ajusta{y \in \lbrace 0, 1, 2\rbrace  = \lbrace \text{Degeneración},\ \text{Normal},\ \text{Protrusión}\rbrace , \qquad \mathbf{y} = \text{one-hot}(y) \in \lbrace 0,1\rbrace ^3}
\end{center}

\textbf{Qué se busca.} Una función con parámetros $\theta$ que asigne a cada imagen una probabilidad para cada clase, y luego la clase más probable:

\begin{center}
\ajusta{f_\theta : \mathbb{R}^{3 \times 224 \times 224} \to \Delta^{2} = \Big\lbrace  p \in [0,1]^3 : \textstyle\sum_k p_k = 1 \Big\rbrace , \qquad \hat{y} = \arg\max_k\ p_k}
\end{center}

\begin{table}[htbp]
\centering\small
\begin{tabular}{p{0.313\linewidth}p{0.313\linewidth}p{0.313\linewidth}}
\toprule
Modelo & Qué hace con X & Parámetros aprendidos \\
\midrule
1. YOLOv8s-cls & Red convolucional completa, ajustada de extremo a extremo & 5\,084\,579 pesos \\
2. ResNet-50 & Red convolucional residual, ajustada de extremo a extremo & 23\,514\,179 pesos (225\,344 congelados) \\
3. SVM (RBF) & Clasificador de máximo margen sobre 256 características & 1\,684 vectores de soporte y sus coeficientes $\alpha_i$ \\
4. Random Forest & 300 árboles de decisión sobre las mismas 256 características & Las divisiones (variable, umbral) de cada nodo \\
\bottomrule
\end{tabular}
\end{table}



\subsection{Modelos 1 y 2: redes neuronales convolucionales}
\label{mm:2}

\subsubsection{Capas}
\label{mm:2.1}
\textbf{Convolución.} Un filtro $W$ de tamaño $k \times k$ recorre la imagen y calcula, en cada posición, una suma ponderada de los píxeles vecinos. Detecta bordes, texturas y, en capas profundas, formas como el contorno posterior del disco:

\[
(W * x)_{o,i,j} = \sum_{c} \sum_{u=0}^{k-1} \sum_{v=0}^{k-1} W_{o,c,u,v}\ x_{c,\ i+u,\ j+v} + b_o
\]

\textbf{Normalización por lotes (BatchNorm).} Estandariza cada canal con la media $\mu_B$ y la varianza $\sigma_B^2$ del lote, y luego reescala con dos parámetros aprendidos $\gamma$ y $\beta$:

\[
\hat{h} = \gamma\thinspace  \frac{h - \mu_B}{\sqrt{\sigma_B^2 + \epsilon}} + \beta
\]

\textbf{Activación} (no linealidad, Figura~\ref{mm:fig1}): la ResNet-50 usa $\text{ReLU}(z) = \max(0, z)$ y YOLOv8 usa $\text{SiLU}(z) = z \cdot \sigma(z)$.

\textbf{Bloque residual (ResNet-50).} La entrada se suma a la salida de tres convoluciones ($1{\times}1$, $3{\times}3$, $1{\times}1$). Así la red solo aprende la \emph{corrección} $\mathcal{F}$, y el gradiente fluye bien por 50 capas. La ResNet-50 apila 16 bloques en 4 etapas de 3, 4, 6 y 3 bloques:

\[
h^{(l+1)} = \text{ReLU}\big(\mathcal{F}(h^{(l)};\ W^{(l)}) + h^{(l)}\big)
\]

\textbf{YOLOv8s-cls.} Columna de convoluciones con paso 2 y bloques C2f (convoluciones con conexiones parciales entre etapas, la misma idea de reutilizar la entrada). Termina en una convolución $1{\times}1$ a 1\,280 canales.

\textbf{Promedio global y capa final.} El mapa de características final se promedia por canal. Queda un vector $\phi(x)$ (2\,048 valores en la ResNet-50, 1\,280 en YOLOv8), que pasa por \emph{dropout} ($p = 0{,}3$: durante el entrenamiento se apaga al azar el 30\,\% de los valores) y por una capa lineal:

\[
\phi(x) = \frac{1}{HW}\sum_{i,j} h^{(L)}_{:,i,j}, \qquad z = W\thinspace \phi(x) + b \in \mathbb{R}^3
\]

En la ResNet-50 esta capa tiene $2\,048 \times 3 + 3 = 6\,147$ parámetros.

\subsubsection{De puntajes a probabilidades: softmax}
\label{mm:2.2}
La regresión logística de la semana 9 usa la sigmoide para 2 clases. Con 3 clases se usa su generalización, la \textbf{softmax}, que convierte los 3 puntajes $z_k$ (\emph{logits}) en probabilidades positivas que suman 1:

\[
p_k = P(y = k \mid x) = \frac{e^{z_k}}{\sum_{j=1}^{3} e^{z_j}}
\]

Ejemplo (Figura~\ref{mm:fig1}): los logits $(2{,}0;\ 0{,}5;\ -1{,}0)$ dan las probabilidades $(0{,}786;\ 0{,}175;\ 0{,}039)$, y la predicción es la clase 0. Con $K = 2$ la softmax se reduce exactamente a la sigmoide.

\begin{figure}[htbp]
\centering
\includegraphics[width=0.9\linewidth]{01_activaciones_softmax.png}
\caption{Activaciones y softmax.}
\label{mm:fig1}
\end{figure}


\subsubsection{Función de costo: entropía cruzada}
\label{mm:2.3}
Como en la regresión logística, el error de una imagen se mide con la \textbf{entropía cruzada} (log loss): cuánto se aleja la distribución predicha $p$ de la distribución objetivo $q$. El costo $J(\theta)$ promedia esa pérdida sobre las imágenes:

\begin{center}
\ajusta{\ell(x_i, y_i) = -\sum_{k=1}^{3} q_{ik} \log p_{ik}, \qquad J(\theta) = \frac{1}{N}\sum_{i=1}^{N} \ell(x_i, y_i) + \text{regularización}}
\end{center}
\begin{itemize}
  \item \textbf{YOLOv8s-cls:} $q = \mathbf{y}$ (one-hot), así que $\ell = -\log p_{y}$. Si la red da $p_y = 0{,}01$ a la clase correcta, la pérdida es 4,6; si da 0,99, es 0,01. Penaliza fuerte la \textbf{seguridad equivocada} (Figura~\ref{mm:fig2}).
  \item \textbf{ResNet-50, con suavizado de etiquetas} ($\varepsilon = 0{,}1$): en lugar de exigir probabilidad 1 a la clase real, el objetivo reparte un poco de masa entre las demás:
\end{itemize}

\begin{center}
\ajusta{q_k = (1-\varepsilon)\thinspace \mathbf{y}_k + \frac{\varepsilon}{K} \quad\Rightarrow\quad q = (0{,}933;\ 0{,}033;\ 0{,}033)}
\end{center}

Así la red no puede volverse arbitrariamente segura, lo que reduce el sobreajuste. El costo mínimo posible ya no es 0 sino la entropía de $q$, $H(q) = 0{,}291$. Ese es el piso al que se acerca la pérdida de train (0,305 en la época 6, Figura~\ref{mm:fig2}).

\textbf{¿Por qué no MSE para entrenar?} El MSE de la semana 9 es la función de costo natural cuando $y$ es un número real (regresión). Para probabilidades de clase, la entropía cruzada da gradientes más grandes cuando el error es grande. Además, es la función que se obtiene al maximizar la verosimilitud del modelo. El MSE sí se usa para \textbf{evaluar}: se reporta como puntaje de Brier (Sección~\ref{mm:4.2}) en el registro de desempeño (Sección~\ref{mm:4.4}).

\begin{figure}[htbp]
\centering
\includegraphics[width=0.9\linewidth]{02_funcion_costo.png}
\caption{Función de costo.}
\label{mm:fig2}
\end{figure}


\subsubsection{Optimización: descenso del gradiente con AdamW}
\label{mm:2.4}
\textbf{Descenso del gradiente por lotes.} En cada paso se toma un lote de $B = 64$ imágenes y se mueve cada parámetro en la dirección opuesta al gradiente del costo, como la regla $\theta_j := \theta_j - \alpha\thinspace  \partial J / \partial \theta_j$ de la semana 9:

\[
g_t = \nabla_\theta \frac{1}{B}\sum_{i \in \text{lote}} \ell(x_i, y_i)
\]

El gradiente de cada peso se calcula con la regla de la cadena desde la salida hacia la entrada (retropropagación).

\textbf{AdamW.} Mantiene un promedio móvil del gradiente ($m_t$, el "impulso") y de su cuadrado ($v_t$, la escala). Con eso cada parámetro avanza con un paso adaptado a su propia escala. El término $\lambda\thinspace \theta$ (\emph{weight decay}) encoge los pesos hacia 0, y es la regularización L2:

\begin{center}
\ajusta{m_t = \beta_1 m_{t-1} + (1-\beta_1)\thinspace  g_t, \qquad v_t = \beta_2 v_{t-1} + (1-\beta_2)\thinspace  g_t^2}
\end{center}

\begin{center}
\ajusta{\hat{m}_t = \frac{m_t}{1-\beta_1^t}, \quad \hat{v}_t = \frac{v_t}{1-\beta_2^t}, \qquad \theta_t = \theta_{t-1} - \alpha_t \left( \frac{\hat{m}_t}{\sqrt{\hat{v}_t} + \epsilon} + \lambda\thinspace  \theta_{t-1} \right)}
\end{center}

Valores: $\beta_1 = 0{,}9$ en la ResNet-50 y $0{,}937$ en YOLOv8 (valor por defecto de ultralytics), $\beta_2 = 0{,}999$, $\epsilon = 10^{-8}$, $\lambda = 0{,}01$. Ultralytics no aplica $\lambda$ a los sesgos ni a las capas de normalización.

\textbf{Tasa de aprendizaje en coseno} (Figura~\ref{mm:fig3}). El tamaño del paso $\alpha_t$ empieza en $\alpha_0$ y baja suavemente hasta $\alpha_{\min}$ en $T$ épocas: pasos grandes al principio, finos al final.

\[
\alpha_t = \alpha_{\min} + \tfrac{1}{2}(\alpha_0 - \alpha_{\min})\left(1 + \cos\frac{\pi t}{T}\right)
\]

\begin{table}[htbp]
\centering\small
\begin{tabular}{p{0.313\linewidth}p{0.313\linewidth}p{0.313\linewidth}}
\toprule
 & YOLOv8s-cls & ResNet-50 \\
\midrule
$\alpha_0$ / $\alpha_{\min}$ & $5\cdot10^{-4}$ / $5\cdot10^{-5}$ (1 época de calentamiento) & $10^{-4}$ / 0 \\
Épocas máximas $T$ / paciencia & 15 / 4 & 10 / 4 \\
Capas congeladas (no se actualizan) & primeros 4 módulos & \arch{conv1}, \arch{bn1}, \arch{layer1} (225\,344 pesos) \\
Aumento en línea & recorte aleatorio con el 100\,\% del área (solo varía la proporción, 3:4 a 4:3) y brillo $\pm$5\,\% & rotación $\pm$6$^{\circ}$, desplazamiento 4\,\%, zoom 0,95–1,05, brillo y contraste $\pm$10\,\% \\
Suavizado de etiquetas & no & $\varepsilon = 0{,}1$ \\
Criterio para elegir la época & exactitud top-1 en validación & F1 macro en validación \\
Época elegida / épocas corridas & 2 / 6 & 2 / 6 \\
\bottomrule
\end{tabular}
\end{table}


\textbf{Ajuste fino (\emph{transfer learning}).} Ninguna red parte de cero: $\theta$ se inicializa con los pesos aprendidos en ImageNet (1,28 millones de fotografías), que ya detectan bordes y texturas. Solo la capa final empieza al azar.

\begin{figure}[htbp]
\centering
\includegraphics[width=0.9\linewidth]{03_tasa_aprendizaje.png}
\caption{Tasa de aprendizaje.}
\label{mm:fig3}
\end{figure}


\subsubsection{Diagnóstico: sobreajuste y parada temprana}
\label{mm:2.5}
La Figura~\ref{mm:fig4} muestra el patrón de las dos redes: la pérdida de train sigue bajando, mientras la de validación sube desde la época 2. El modelo empieza a memorizar las imágenes de train, que además incluyen copias aumentadas de las mismas imágenes. Por eso se guarda la época con mejor desempeño en validación y el entrenamiento se detiene tras 4 épocas sin mejora (\textbf{parada temprana}):

\[
e^{*} = \arg\max_{e}\ \text{F1}_{\text{macro}}\big(\text{validación};\ \theta^{(e)}\big)
\]

\begin{figure}[htbp]
\centering
\includegraphics[width=0.9\linewidth]{04_curvas_aprendizaje.png}
\caption{Curvas de aprendizaje.}
\label{mm:fig4}
\end{figure}



\subsection{Modelos 3 y 4: características profundas + clasificador clásico}
\label{mm:3}

Los modelos 3 y 4 no trabajan sobre píxeles. Cada imagen se representa con el vector $\phi(x) \in \mathbb{R}^{2048}$ que produce la ResNet-50 del modelo 2 justo antes de su capa final (Sección~\ref{mm:2.1}). Después se reduce ese vector a 256 dimensiones y se clasifica con un SVM o con un Random Forest.

\textbf{Dependencia declarada:} los modelos 3 y 4 dependen del modelo 2. Con un extractor de ImageNet sin ajustar rinden mucho menos (Sección~\ref{mm:3.4}).

\subsubsection{Estandarización y PCA}
\label{mm:3.1}
\textbf{Estandarización.} Cada una de las 2\,048 características se centra y se escala con la media y la desviación \textbf{calculadas solo con train}. Calcularlas con todo el dataset filtraría información del test (\emph{data leakage}, semanas 6 y 9):

\[
s_j = \frac{\phi_j(x) - \mu_j^{\text{train}}}{\sigma_j^{\text{train}}}
\]

\textbf{Análisis de componentes principales (PCA).} Se busca la base de direcciones ortogonales en la que los datos de train varían más. Son los vectores propios de la matriz de covarianza, y cada imagen se proyecta sobre los $k = 256$ primeros:

\begin{center}
\ajusta{\Sigma = \frac{1}{n-1}\thinspace  S^\top S = V \Lambda V^\top, \qquad u = V_k^\top s \in \mathbb{R}^{256}, \qquad \text{varianza explicada} = \frac{\sum_{j=1}^{k}\lambda_j}{\sum_{j=1}^{2048}\lambda_j}}
\end{center}

Con las características ajustadas, \textbf{256 componentes explican el 80,0\,\% de la varianza} (las 2 primeras ya el 38\,\%). Con las de ImageNet sin ajustar, solo el 58,7\,\% (Figura~\ref{mm:fig5}). Las características ajustadas están más concentradas en pocas direcciones, las que separan las clases.

\begin{figure}[htbp]
\centering
\includegraphics[width=0.9\linewidth]{05_pca_varianza.png}
\caption{Varianza explicada por el PCA.}
\label{mm:fig5}
\end{figure}


\subsubsection{Modelo 3: máquina de vectores de soporte (SVM) con kernel RBF}
\label{mm:3.2}
\textbf{Idea.} Entre todas las fronteras que separan dos clases, el SVM elige la de \textbf{máximo margen}: la que queda más lejos de los ejemplos más cercanos de cada lado. Esos ejemplos se llaman vectores de soporte. Para dos clases, $y_i \in \lbrace -1, +1\rbrace $, con margen blando:

\begin{center}
\ajusta{\min_{w,\thinspace b,\thinspace \xi}\ \ \frac{1}{2}\lVert w \rVert^2 + C \sum_{i=1}^{n} \xi_i \quad \text{sujeto a} \quad y_i\big(w^\top \varphi(u_i) + b\big) \ge 1 - \xi_i,\quad \xi_i \ge 0}
\end{center}
\begin{itemize}
  \item $\frac{1}{2}\lVert w\rVert^2$: minimizarlo maximiza el ancho del margen, $2 / \lVert w \rVert$.
  \item $\xi_i$: holgura que permite que algunas imágenes queden dentro del margen o mal clasificadas.
  \item $C$: equilibrio entre margen ancho y errores en train. Un $C$ pequeño tolera más errores y da una frontera más suave, menos sobreajustada.
\end{itemize}

Es equivalente a minimizar la \textbf{pérdida bisagra} (\emph{hinge}): $\frac{1}{2}\lVert w\rVert^2 + C\sum_i \max\big(0,\ 1 - y_i f(u_i)\big)$.

\textbf{Kernel RBF.} Las clases no se separan con un hiperplano en el espacio original. El SVM resuelve el problema dual, en el que los datos solo aparecen como similitudes $K(u_i, u_j)$; así trabaja implícitamente en un espacio de dimensión infinita (el "truco del kernel"):

\begin{center}
\ajusta{K(u, v) = \exp\big(-\gamma\thinspace \lVert u - v \rVert^2\big), \qquad \gamma = \frac{1}{d \cdot \mathrm{Var}(U)} = 0{,}00061\ \ (\texttt{gamma="scale"},\ d = 256)}
\end{center}

\begin{center}
\ajusta{\max_{\alpha}\ \sum_i \alpha_i - \frac{1}{2}\sum_i\sum_j \alpha_i\alpha_j\thinspace  y_i y_j\thinspace  K(u_i, u_j) \quad \text{s.a.}\ \ 0 \le \alpha_i \le C,\ \ \sum_i \alpha_i y_i = 0}
\end{center}

\[
f(u) = \sum_{i \in \text{SV}} \alpha_i\thinspace  y_i\thinspace  K(u_i, u) + b
\]

Solo los vectores de soporte ($\alpha_i > 0$) intervienen en la predicción. Con $C = 0{,}1$, elegido en validación (Figura~\ref{mm:fig8}), el modelo usa \textbf{1\,684 vectores de soporte de 24\,501 imágenes de train (6,9\,\%)}: 320 de Degeneración, 660 de Normal y 704 de Protrusión.

La Figura~\ref{mm:fig6}, a la derecha, muestra cómo decae la similitud con la distancia. Con el $\gamma$ usado, dos imágenes a la distancia mediana (59) tienen similitud $\approx$ 0,12. Con $\gamma \times 10$ todas las imágenes serían "lejanas" (sobreajuste); con $\gamma / 10$, todas "parecidas" (subajuste).

\textbf{Tres clases.} scikit-learn entrena un SVM por cada par de clases (\emph{uno contra uno}: 3 clasificadores) y combina sus votos en un puntaje por clase; la predicción es la clase con mayor puntaje. Estos puntajes \textbf{no son probabilidades}; por eso el SVM no tiene log loss y la app los muestra normalizados con softmax.

La Figura~\ref{mm:fig6}, a la izquierda, es una \textbf{ilustración en 2 dimensiones} (las 2 primeras componentes, 1\,200 originales de train). En train las clases aparecen muy separadas porque la ResNet ya vio esas imágenes; en validación y test se mezclan más, y por eso la exactitud allí baja a ~88\,\%.

\begin{figure}[htbp]
\centering
\includegraphics[width=0.9\linewidth]{06_svm_rbf.png}
\caption{SVM con kernel RBF.}
\label{mm:fig6}
\end{figure}


\subsubsection{Modelo 4: Random Forest}
\label{mm:3.3}
\textbf{Árbol de decisión.} Hace preguntas del tipo "¿$u_j \le t$?" y divide las imágenes en dos grupos. En cada nodo elige la variable $j$ y el umbral $t$ que más reducen la \textbf{impureza de Gini} (semana 11):

\begin{center}
\ajusta{G(\text{nodo}) = 1 - \sum_{k=1}^{3} p_k^2, \qquad \Delta G = G(\text{padre}) - \frac{n_L}{n}\thinspace G(L) - \frac{n_R}{n}\thinspace G(R)}
\end{center}

$G = 0$ si el nodo tiene una sola clase (puro). El máximo es $1 - 1/K$ si las clases están mezcladas en partes iguales: 0,5 con 2 clases, 0,667 con 3. Ejemplo de la clase: 6 compran y 4 no ($G = 1 - 0{,}6^2 - 0{,}4^2 = 0{,}48$). Tras dividir en un grupo (5, 0) con $G = 0$ y otro (1, 4) con $G = 0{,}32$, la impureza ponderada baja a $0{,}5 \cdot 0 + 0{,}5 \cdot 0{,}32 = 0{,}16$.

\textbf{Bosque.} Un árbol solo sobreajusta. Se entrenan $T = 300$ árboles distintos entre sí por dos fuentes de azar:
\begin{enumerate}
  \item \textbf{Bootstrap:} cada árbol se entrena con una muestra de train tomada con reemplazo, del mismo tamaño.
  \item \textbf{Variables al azar:} en cada división solo se consideran $m = \sqrt{256} = 16$ de las 256 variables.
\end{enumerate}

La predicción promedia las probabilidades de los 300 árboles (votación suave). La probabilidad de cada árbol, $p_t(k \mid u)$, es la proporción de la clase $k$ en la hoja donde cae $u$:

\[
P(y = k \mid u) = \frac{1}{T}\sum_{t=1}^{T} p_t(k \mid u), \qquad \hat{y} = \arg\max_k P(y = k \mid u)
\]

Al promediar árboles diversos baja la varianza (el sobreajuste) sin subir mucho el sesgo. Hiperparámetro elegido en validación (Figura~\ref{mm:fig8}): profundidad máxima 20. Los 300 árboles llegan a esa profundidad y tienen 779 hojas en promedio.

La \textbf{importancia} de cada variable es la disminución total de Gini que aporta, promediada sobre los árboles. Las 10 primeras componentes del PCA concentran el 70\,\% de la importancia (Figura~\ref{mm:fig7}).

\begin{figure}[htbp]
\centering
\includegraphics[width=0.9\linewidth]{07_random_forest.png}
\caption{Random Forest.}
\label{mm:fig7}
\end{figure}


\subsubsection{Extras del curso y comparación con ImageNet}
\label{mm:3.4}
Sobre los mismos 256 componentes se entrenaron, como línea base del curso:
\begin{itemize}
  \item \textbf{Regresión logística multinomial} (semana 9): $p = \text{softmax}(W u + b)$, con regularización L2 de fuerza $1/C$ y $C = 0{,}01$.
  \item \textbf{KNN} (semana 10): predice por mayoría entre los $K = 31$ vecinos más cercanos según la distancia euclídea $d(u, v) = \sqrt{\sum_j (u_j - v_j)^2}$. La escala importa, y por eso se estandariza antes.
\end{itemize}

Con características de ImageNet sin ajustar, el SVM baja de 88,4\,\% a 77,6\,\% y el Random Forest de 88,4\,\% a 75,3\,\% de exactitud en el test de originales. Lo que hace funcionar a los modelos 3 y 4 es la representación aprendida por el modelo 2.

\begin{figure}[htbp]
\centering
\includegraphics[width=0.9\linewidth]{08_busqueda_hiperparametros.png}
\caption{Búsqueda de hiperparámetros.}
\label{mm:fig8}
\end{figure}


\subsubsection{Estructura del sistema: arquitectura híbrida}
\label{mm:3.5}

El sistema tiene \textbf{dos caminos}. YOLOv8s-cls va por el suyo, independiente. La ResNet-50 es la \textbf{base compartida} de los otros tres modelos: sus características $\phi(x)$ alimentan su propia capa final (modelo 2) y, después de estandarizarlas y reducirlas, al SVM (modelo 3) y al Random Forest (modelo 4). Unir una red profunda que extrae características con un clasificador clásico que decide es lo que se llama \textbf{arquitectura híbrida} (Figura~\ref{fig:estructura}).


\textbf{Las 4 funciones.} Sea $g_Y$ la red YOLOv8s-cls hasta sus logits, $g_R$ la ResNet-50 sin su capa final y $\phi(x) = g_R(x)$. Cada modelo es una composición de funciones:

\[
p_1(x) = \text{softmax}\big(g_Y(x)\big)
\]

\[
p_2(x) = \text{softmax}\big(W\thinspace \phi(x) + b\big)
\]

\begin{center}
\ajusta{u(x) = V_k^\top\thinspace D^{-1}\big(\phi(x) - \mu\big), \qquad D = \text{diag}(\sigma_1, \dots, \sigma_{2048})}
\end{center}

\begin{center}
\ajusta{p_3(x) = \text{softmax}\big(\text{SVM}(u(x))\big), \qquad p_4(x) = \frac{1}{T}\sum_{t=1}^{T} p_t\big(u(x)\big)}
\end{center}

Las tres últimas dependen de la misma $\phi(x)$. Por eso, en la aplicación la ResNet-50 se ejecuta una sola vez por imagen. La softmax de $p_3$ solo lleva los puntajes del SVM a la misma escala de los demás: no son probabilidades calibradas.

\textbf{Orden de entrenamiento.} Cada etapa usa solo lo aprendido en las anteriores, y nada se ajusta con el test:

\begin{table}[htbp]
\centering\small
\begin{tabular}{p{0.235\linewidth}p{0.235\linewidth}p{0.235\linewidth}p{0.235\linewidth}}
\toprule
Paso & Qué se aprende & Con qué datos & Qué queda fijo \\
\midrule
1 & Pesos de YOLOv8s-cls & train; época elegida con validación & — \\
2 & Pesos de la ResNet-50, incluidos $W$ y $b$ & train; época elegida con validación & — \\
3 & $\mu$, $D$ y $V_k$ (escalador y PCA) & características de train & La ResNet-50 \\
4 & SVM y Random Forest & $u$ de train; hiperparámetros elegidos con validación & ResNet-50, escalador y PCA \\
\bottomrule
\end{tabular}
\end{table}


\textbf{Por qué es híbrida y no un ensamble.} En esta estructura un modelo le pasa sus \textbf{características} a otro, y cada uno sigue dando su propia predicción. Un ensamble, en cambio, combina las \textbf{respuestas finales} de varios modelos. Esa combinación es la sección siguiente.

\subsubsection{Ensamble por promedio}
\label{mm:3.6}

\textbf{Idea.} Es la misma de la semana 11: consultar a varios modelos reduce la dependencia de una sola decisión. Cada modelo entrega 3 puntajes para la imagen; se promedian por clase y gana la clase con el mayor promedio (\textbf{votación suave}). Con $M = 4$ modelos:

\begin{center}
\ajusta{P_{\text{ens}}(k \mid x) = \frac{1}{M}\sum_{m=1}^{M} p_m(k \mid x), \qquad \hat{y}_{\text{ens}} = \arg\max_k\ P_{\text{ens}}(k \mid x)}
\end{center}

El ensamble \textbf{no se entrena}: todos los modelos pesan igual ($1/4$) y no tiene parámetros que aprender. Se calcula después, con las salidas de los 4 modelos ya entrenados. Por eso es un \textbf{extra} y no un quinto modelo.

\textbf{Ejemplo real} (imagen de test: paciente 0319, L1-L2, etiqueta real Degeneración):

\begin{table}[htbp]
\centering\small
\begin{tabular}{lllll}
\toprule
Modelo & Degeneración & Normal & Protrusión & Predicción \\
\midrule
1. YOLOv8s-cls & 0,109 & 0,871 & 0,020 & Normal (falla) \\
2. ResNet-50 & 0,936 & 0,025 & 0,039 & Degeneración \\
3. SVM (con softmax) & 0,746 & 0,064 & 0,190 & Degeneración \\
4. Random Forest & 0,866 & 0,050 & 0,083 & Degeneración \\
\textbf{Ensamble (promedio)} & \textbf{0,664} & 0,253 & 0,083 & \textbf{Degeneración} \\
\bottomrule
\end{tabular}
\end{table}


Por ejemplo, $P_{\text{ens}}(\text{Degeneración}) = (0{,}109 + 0{,}936 + 0{,}746 + 0{,}866) / 4 = 0{,}664$. Los otros tres modelos corrigen el error de YOLO.

\textbf{Por qué puede ayudar, y su límite.} Si cada modelo tiene una varianza de error $\sigma^2$ y la correlación entre modelos es $\rho$, la varianza del promedio de $M$ modelos es:

\begin{center}
\ajusta{\text{Var}\left(\frac{1}{M}\sum_{m=1}^{M} p_m\right) = \rho\thinspace \sigma^2 + \frac{1-\rho}{M}\thinspace \sigma^2}
\end{center}

Si los modelos fueran independientes ($\rho = 0$), la varianza se dividiría entre 4. Si fueran idénticos ($\rho = 1$), promediar no serviría de nada. Aquí los modelos 2, 3 y 4 comparten $\phi(x)$, así que se espera una correlación alta y una mejora limitada.

\textbf{Resultados} (test de originales; detalle en las Secciones~\ref{mm:4.3} y~\ref{mm:4.4}):

\begin{table}[htbp]
\centering\small
\begin{tabular}{lll}
\toprule
Métrica & Mejor modelo individual & Ensamble \\
\midrule
Exactitud & 0,889 (ResNet-50) & 0,889 \\
F1 macro & 0,769 (ResNet-50) & 0,745 \\
ROC-AUC macro & 0,942 (Random Forest) & \textbf{0,947} \\
MSE (Brier) & 0,173 (ResNet-50) & \textbf{0,164} \\
\bottomrule
\end{tabular}
\end{table}


El ensamble empata en exactitud y da las probabilidades más confiables (mejor ROC-AUC y MSE), pero detecta menos casos de Degeneración: 17 de 42, frente a 21 de la ResNet-50. En validación, el criterio acordado, su F1 macro es 0,822 frente a 0,824 de la ResNet-50. Por eso \textbf{la ResNet-50 sigue siendo el modelo principal} y el ensamble se reporta como resultado adicional.


\subsection{Evaluación}
\label{mm:4}

\subsubsection{Protocolo}
\label{mm:4.1}
\begin{itemize}
  \item Los parámetros $\theta$ se aprenden solo con \textbf{train}.
  \item Los hiperparámetros ($C$, profundidad, $K$, época) se eligen con \textbf{validación}, por F1 macro.
  \item El \textbf{test} se evalúa una sola vez, al final. Ningún paciente aparece en dos particiones, y cada copia aumentada queda en la partición de su original.
  \item Se reporta el test de \textbf{originales} (768 imágenes en proporción natural: 42 Degeneración, 164 Normal, 562 Protrusión) y el \textbf{test completo} (5\,253, balanceado con copias).
\end{itemize}

\subsubsection{Matriz de confusión y métricas}
\label{mm:4.2}
En clase se vieron las fórmulas para \textbf{2 clases} (positiva y negativa):

\begin{center}
\ajusta{\text{Accuracy} = \frac{TP + TN}{TP + TN + FP + FN}, \qquad \text{Precision} = \frac{TP}{TP + FP}, \qquad \text{Recall} = \frac{TP}{TP + FN}, \qquad F1 = 2\cdot\frac{\text{Precision}\cdot\text{Recall}}{\text{Precision} + \text{Recall}}}
\end{center}

Con \textbf{3 clases} se aplican las mismas fórmulas tres veces. Cada clase $k$ se trata como la "positiva" y las otras dos como la "negativa" (uno contra el resto): TP = predijo $k$ y era $k$; FP = predijo $k$ y no era $k$; FN = era $k$ y predijo otra; TN = el resto. Luego se promedian las tres (promedio \emph{macro}). La exactitud no necesita promedio: es el total de aciertos (la diagonal de la matriz) sobre el total de imágenes.

\begin{center}
\ajusta{\text{Exactitud} = \frac{\sum_k \text{TP}_k}{N}, \qquad \text{Precisión}_k = \frac{\text{TP}_k}{\text{TP}_k + \text{FP}_k}, \qquad \text{Recall}_k = \frac{\text{TP}_k}{\text{TP}_k + \text{FN}_k}}
\end{center}

\begin{center}
\ajusta{F1_k = 2\thinspace \frac{\text{Precisión}_k \cdot \text{Recall}_k}{\text{Precisión}_k + \text{Recall}_k}, \qquad F1_{\text{macro}} = \frac{1}{K}\sum_k F1_k, \qquad \text{Exactitud balanceada} = \frac{1}{K}\sum_k \text{Recall}_k}
\end{center}

\textbf{Ejemplo con el modelo 2 en el test de originales} (filas = real, columnas = predicha):

\begin{table}[htbp]
\centering\small
\begin{tabular}{llll}
\toprule
Real \textbackslash{} Predicha & Degeneración & Normal & Protrusión \\
\midrule
\textbf{Degeneración} (42) & \textbf{21} & 6 & 15 \\
\textbf{Normal} (164) & 5 & \textbf{133} & 26 \\
\textbf{Protrusión} (562) & 6 & 27 & \textbf{529} \\
\bottomrule
\end{tabular}
\end{table}


\begin{table}[htbp]
\centering\small
\begin{tabular}{lllllll}
\toprule
Clase & TP & FP & FN & Precisión & Recall & F1 \\
\midrule
Degeneración & 21 & 11 & 21 & 0,656 & 0,500 & 0,568 \\
Normal & 133 & 33 & 31 & 0,801 & 0,811 & 0,806 \\
Protrusión & 529 & 41 & 33 & 0,928 & 0,941 & 0,935 \\
\bottomrule
\end{tabular}
\end{table}

\begin{itemize}
  \item \textbf{Exactitud} = (21 + 133 + 529) / 768 = 683 / 768 = \textbf{0,889}.
  \item \textbf{F1 macro} = (0,568 + 0,806 + 0,935) / 3 = \textbf{0,769}.
\end{itemize}

\textbf{Por qué la exactitud sola engaña.} Un modelo que siempre dijera "Protrusión" acertaría 562 / 768 = \textbf{73,2\,\%} sin detectar ni un caso de Degeneración. Es el mismo ejemplo del fraude visto en clase. El F1 macro pesa igual las tres clases y deja ver que Degeneración se detecta solo en la mitad de los casos.

\textbf{Intervalo de confianza.} Con 768 imágenes, la exactitud de 88,9\,\% tiene un intervalo de Wilson al 95\,\% de \textbf{[86,5\,\%; 91,0\,\%]}:

\begin{center}
\ajusta{\frac{\hat{p} + \frac{z^2}{2n} \pm z\sqrt{\frac{\hat{p}(1-\hat{p})}{n} + \frac{z^2}{4n^2}}}{1 + \frac{z^2}{n}}, \qquad z = 1{,}96}
\end{center}

Las diferencias entre los 4 modelos (87,8\,\% a 88,9\,\%) están dentro de ese margen.

\textbf{ROC-AUC.} Para cada clase, la curva ROC recorre todos los umbrales $\tau$ sobre $p_k$ y grafica la tasa de verdaderos positivos contra la de falsos positivos. El área bajo la curva es la probabilidad de que una imagen de la clase $k$ reciba más puntaje que una que no lo es. Se promedia sobre las 3 clases (uno contra el resto).

\textbf{Log loss.} Es la entropía cruzada de la Sección~\ref{mm:2.3}, evaluada con las probabilidades sobre el conjunto de prueba.

\textbf{MSE (puntaje de Brier).} Es el MSE de la semana 9 aplicado a la clasificación. El "valor real" de cada imagen es su etiqueta en one-hot (por ejemplo, Protrusión = $(0, 0, 1)$) y la "predicción" son las 3 probabilidades del modelo. Se suman los errores al cuadrado de las 3 clases y se promedia sobre las $N$ imágenes:

\begin{center}
\ajusta{\text{MSE} = \frac{1}{N}\sum_{i=1}^{N}\sum_{k=1}^{3}\left(p_{ik} - y_{ik}\right)^2, \qquad \text{RMSE} = \sqrt{\text{MSE}}}
\end{center}

Vale 0 si el modelo da probabilidad 1 a la clase correcta y como máximo 2. Igual que el MSE de clase, castiga más los errores grandes. Ejemplo con una imagen real de Protrusión:
\begin{itemize}
  \item Si el modelo da $p = (0{,}1;\ 0{,}2;\ 0{,}7)$: $0{,}1^2 + 0{,}2^2 + 0{,}3^2 = 0{,}14$ (acierta, con algo de duda).
  \item Si da $p = (0{,}9;\ 0{,}05;\ 0{,}05)$: $0{,}9^2 + 0{,}05^2 + 0{,}95^2 = 1{,}715$ (falla con mucha seguridad).
\end{itemize}

A diferencia de la exactitud, el MSE también mide si la \textbf{confianza} del modelo es la adecuada. No aplica al SVM, porque sus puntajes no son probabilidades.

\subsubsection{Resultados (test de originales)}
\label{mm:4.3}
Los resultados en los originales de prueba están en la Tabla~\ref{tab:resultados}, con sus matrices de confusión y curvas ROC en las Figuras~\ref{fig:matrices} y~\ref{fig:roc}.

\subsubsection{Registro de desempeño (formato de la clase)}
\label{mm:4.4}

Es la tabla "Registro de desempeño" de la semana 9, con columnas Train y Test.
\begin{itemize}
  \item \textbf{Train:} las 3\,577 imágenes originales de entrenamiento.
  \item \textbf{Test:} las 768 originales de test.
  \item \textbf{Test completo:} las 5\,253 imágenes de test con copias aumentadas, como referencia adicional.
\end{itemize}

Precision, recall y F1 son promedios \emph{macro} de las 3 clases. \textbf{R$^2$ no aplica}: es una métrica de regresión y este problema es de clasificación (como advierte la diapositiva, no todas las métricas aplican a todos los modelos).

\begin{table}[htbp]
\centering\small
\caption{Registro de desempeño. Entrenamiento: 3\,577 originales. Prueba: 768 originales. Prueba completa: 5\,253 imágenes con copias, balanceada. Precision, recall y F1 son promedios macro.}
\label{tab:registro}
\begin{tabular}{lllll}
\toprule
Modelo & Métrica & Train & Test & Test completo \\
\midrule
\textbf{1. YOLOv8s-cls} & MSE / RMSE (Brier) & 0,024 / 0,154 & 0,196 / 0,442 & 0,403 / 0,635 \\
 & Accuracy & 0,984 & 0,878 & 0,759 \\
 & Precision / Recall & 0,959 / 0,989 & 0,771 / 0,732 & 0,779 / 0,759 \\
 & F1 / ROC-AUC & 0,973 / 0,999 & 0,750 / 0,931 & 0,752 / 0,923 \\
\textbf{2. ResNet-50} & MSE / RMSE (Brier) & 0,047 / 0,217 & 0,173 / 0,416 & 0,349 / 0,591 \\
 & Accuracy & 0,971 & \textbf{0,889} & 0,775 \\
 & Precision / Recall & 0,948 / 0,981 & 0,795 / 0,751 & 0,799 / 0,775 \\
 & F1 / ROC-AUC & 0,964 / 0,998 & 0,769 / 0,919 & 0,768 / 0,915 \\
\textbf{3. SVM (RBF)} & MSE / RMSE (Brier) & no aplica & no aplica & no aplica \\
 & Accuracy & 0,979 & 0,884 & 0,754 \\
 & Precision / Recall & 0,964 / 0,984 & 0,797 / 0,718 & 0,789 / 0,754 \\
 & F1 / ROC-AUC & 0,974 / 0,998 & 0,748 / 0,899 & 0,746 / 0,905 \\
\textbf{4. Random Forest} & MSE / RMSE (Brier) & 0,010 / 0,098 & 0,174 / 0,416 & 0,370 / 0,608 \\
 & Accuracy & 0,998 & 0,884 & 0,743 \\
 & Precision / Recall & 0,994 / 0,999 & 0,810 / 0,710 & 0,784 / 0,743 \\
 & F1 / ROC-AUC & 0,996 / 1,000 & 0,746 / 0,942 & 0,731 / 0,920 \\
\textbf{Extra: ensamble} & MSE / RMSE (Brier) & 0,031 / 0,175 & 0,164 / 0,405 & 0,334 / 0,578 \\
 & Accuracy & 0,989 & 0,889 & 0,762 \\
 & Precision / Recall & 0,978 / 0,992 & 0,782 / 0,721 & 0,794 / 0,762 \\
 & F1 / ROC-AUC & 0,985 / 1,000 & 0,745 / 0,947 & 0,752 / 0,937 \\
\bottomrule
\end{tabular}
\end{table}


\textbf{Lectura:}
\begin{itemize}
  \item La distancia entre Train y Test es la que se vio en las curvas de aprendizaje (Sección~\ref{mm:2.5}): \textbf{sobreajuste}. En train los modelos casi no se equivocan (MSE de 0,01 a 0,05); en test el MSE sube a ~0,17–0,20.
  \item Por MSE, el ensamble es el mejor en test (0,164), seguido de la ResNet-50 y el Random Forest (0,173 y 0,174).
  \item Todas las cifras salen de \arch{docs/resultados/comparacion/tabla_metricas_larga.csv}, que \arch{scripts/entrenamiento/05_comparar_modelos.py} calcula a partir de los puntajes guardados de cada modelo.
\end{itemize}


\subsection{Supuestos y limitaciones del modelo}
\label{mm:5}
\begin{itemize}
  \item \textbf{Independencia:} se supone que los pacientes son independientes entre sí. Las imágenes de un mismo paciente no lo son, por eso se divide por paciente. Las copias aumentadas no son datos nuevos: el tamaño efectivo es de 399 pacientes, no de 35\,007 imágenes.
  \item \textbf{Etiquetas con ruido:} $y$ la asignaron estudiantes, sin radiólogo y sin corte axial. El modelo aprende a reproducir esas etiquetas, no un diagnóstico.
  \item \textbf{Desbalance:} el entrenamiento ve las clases balanceadas (por copias), pero en la proporción natural el 73\,\% es Protrusión. Se probó corregir por esa proporción multiplicando $p_k$ por la frecuencia de cada clase; no mejoró la exactitud y hundió el recall de Degeneración, así que se descartó (\arch{docs/resultados/experimentos_validacion.md}).
  \item \textbf{Sobreajuste:} las redes llegan a ~97–99\,\% en train y ~89\,\% en test de originales. La regularización (dropout, weight decay, suavizado, aumento) y la parada temprana lo reducen, pero no lo eliminan.
  \item \textbf{Dependencia de los modelos 3 y 4} respecto del modelo 2: comparten la representación $\phi(x)$, así que no son cuatro opiniones independientes. Eso limita lo que puede ganar el ensamble (Sección~\ref{mm:3.6}).
  \item \textbf{Herramienta académica, no diagnóstica.}
\end{itemize}

% =====================================================================
% Seccion: Entrenamiento y evaluacion de los cuatro modelos
% Fragmento SIN preambulo: % =====================================================================
% Seccion: Entrenamiento y evaluacion de los cuatro modelos
% Fragmento SIN preambulo: % =====================================================================
% Seccion: Entrenamiento y evaluacion de los cuatro modelos
% Fragmento SIN preambulo: \input{secciones/entrenamiento_modelos}
% Cifras: docs/resultados/ y docs/modelo_matematico/README.md
% =====================================================================
\section{Entrenamiento y evaluación de los modelos}
\label{sec:entrenamiento}

Los cuatro modelos se entrenaron con el dataset de la Sección~\ref{sec:dataset}, con la misma división por paciente y la semilla 42. Cada script comprueba que ningún paciente esté en dos particiones y que el dataset coincida con \arch{manifest_35k.csv}. Los scripts están en \arch{scripts/entrenamiento/}, los resultados en \arch{docs/resultados/} y las ecuaciones de cada modelo en \arch{docs/modelo_matematico/README.md}.

\subsection{Estructura del sistema}
La Figura~\ref{fig:estructura} muestra la arquitectura. YOLOv8s-cls sigue un camino independiente. La ResNet-50 es la base de los otros tres modelos: sus 2\,048 características alimentan su propia capa final (modelo~2) y, una vez estandarizadas y reducidas a 256 componentes con PCA, al SVM (modelo~3) y al Random Forest (modelo~4). El ensamble promedia las salidas de los cuatro; no se entrena y se reporta como resultado adicional, no como quinto modelo.

\begin{figure}[htbp]
\centering
\includegraphics[width=\linewidth]{09_estructura_sistema.png}
\caption{Estructura del sistema: los modelos 2, 3 y 4 comparten las características de la ResNet-50.}
\label{fig:estructura}
\end{figure}

\subsection{Protocolo de entrenamiento}
El orden fue el siguiente, y cada etapa usó solo lo aprendido en las anteriores:
\begin{enumerate}
  \item Ajuste de YOLOv8s-cls con entrenamiento, eligiendo la época con validación.
  \item Ajuste de la ResNet-50 con entrenamiento, eligiendo la época con validación.
  \item Ajuste del escalador y del PCA con las características de entrenamiento.
  \item Ajuste del SVM y del Random Forest, eligiendo sus hiperparámetros con validación.
  \item Evaluación única sobre la partición de prueba, con la configuración ya cerrada.
\end{enumerate}
La Tabla~\ref{tab:config} resume la configuración de las dos redes. Ambas parten de pesos de ImageNet y usan AdamW con tasa de aprendizaje en coseno.

\begin{table}[htbp]
\centering\small
\caption{Configuración de entrenamiento de las redes.}
\label{tab:config}
\begin{tabular}{p{0.30\linewidth}p{0.29\linewidth}p{0.29\linewidth}}
\toprule
 & YOLOv8s-cls & ResNet-50 \\
\midrule
Parámetros & 5\,084\,579 & 23\,514\,179 (225\,344 congelados) \\
Tasa inicial / mínima & $5\cdot10^{-4}$ / $5\cdot10^{-5}$ & $10^{-4}$ / 0 \\
Épocas máximas / paciencia & 15 / 4 & 10 / 4 \\
Capas congeladas & primeros 4 módulos & \arch{conv1}, \arch{bn1}, \arch{layer1} \\
Suavizado de etiquetas & no & $\varepsilon = 0{,}1$ \\
Criterio de la época & exactitud en validación & F1 macro en validación \\
Época elegida / corridas & 2 / 6 & 2 / 6 \\
\bottomrule
\end{tabular}
\end{table}

\paragraph{SVM y Random Forest.} Con validación se eligió $C = 0{,}1$ para el SVM (núcleo RBF), que quedó con 1\,684 vectores de soporte de 24\,501 imágenes, y una profundidad máxima de 20 para el Random Forest de 300 árboles (Figura~\ref{fig:hiper}). Con PCA, 256 componentes explican el 80,0\,\% de la varianza de las características ajustadas.

\begin{figure}[htbp]
\centering
\includegraphics[width=\linewidth]{08_busqueda_hiperparametros.png}
\caption{Búsqueda de hiperparámetros con la partición de validación.}
\label{fig:hiper}
\end{figure}

\subsection{Evidencia del entrenamiento}
Las Figuras~\ref{fig:curvas1} y~\ref{fig:curvas2} son las curvas registradas. En las dos redes la mejor época es la 2: desde ahí la pérdida de entrenamiento sigue bajando y la de validación empeora. En la ResNet-50, el F1 macro de validación llegó a 0,824 en la época 2 y bajó en las siguientes, y la parada temprana detuvo el proceso en la época 6.

\begin{figure}[htbp]
\centering
\includegraphics[width=0.85\linewidth]{modelo1_yolov8s_cls/curvas_entrenamiento.png}
\caption{Curvas de entrenamiento de YOLOv8s-cls.}
\label{fig:curvas1}
\end{figure}

\begin{figure}[htbp]
\centering
\includegraphics[width=0.85\linewidth]{modelo2_resnet50/curvas_entrenamiento.png}
\caption{Curvas de entrenamiento de la ResNet-50.}
\label{fig:curvas2}
\end{figure}

\subsection{Ajustes probados en validación}
Se probaron variantes para acercarse al 89\,\% de exactitud que pide la rúbrica, todas decididas con validación (registro en \arch{docs/resultados/experimentos_validacion.md}). El cambio de mayor efecto fue sustituir el extractor de ImageNet por la ResNet-50 ajustada: el F1 de validación subió de 0,65 a 0,80 en el SVM y de 0,60 a 0,79 en el Random Forest. A cambio, los modelos 3 y 4 dependen del modelo~2. Las demás variantes se estancaron en una exactitud de 86--87,5\,\% en validación.

\subsection{Resultados}
\label{sec:resultados}
La Tabla~\ref{tab:resultados} recoge los resultados en la prueba de originales. El registro de la Tabla~\ref{tab:registro} compara entrenamiento y prueba. Las Figuras~\ref{fig:matrices}, \ref{fig:roc} y~\ref{fig:barras} muestran las matrices de confusión, las curvas ROC y la comparación con el conjunto de prueba completo.

\begin{table}[htbp]
\centering\small
\caption{Resultados en los originales de prueba (768 imágenes: 42 Degeneración, 164 Normal, 562 Protrusión).}
\label{tab:resultados}
\begin{tabular}{lrrrrrr}
\toprule
 & & & & \multicolumn{3}{c}{Recall} \\
\cmidrule(lr){5-7}
Modelo & Exactitud & F1 macro & ROC-AUC & Deg. & Normal & Prot. \\
\midrule
1. YOLOv8s-cls        & 87,8\,\% & 0,750 & 0,931 & 0,50 & 0,76 & 0,94 \\
\textbf{2. ResNet-50} & \textbf{88,9\,\%} & \textbf{0,769} & 0,919 & 0,50 & 0,81 & 0,94 \\
3. SVM (RBF)          & 88,4\,\% & 0,748 & 0,899 & 0,43 & 0,77 & 0,95 \\
4. Random Forest      & 88,4\,\% & 0,746 & 0,942 & 0,40 & 0,77 & 0,95 \\
\midrule
Extra: ensamble       & 88,9\,\% & 0,745 & 0,947 & 0,40 & 0,81 & 0,95 \\
\bottomrule
\end{tabular}
\end{table}

\begin{figure}[htbp]
\centering
\includegraphics[width=\linewidth]{comparacion/matrices_confusion_test_orig.png}
\caption{Matrices de confusión de los cuatro modelos en los originales de prueba.}
\label{fig:matrices}
\end{figure}

\begin{figure}[htbp]
\centering
\includegraphics[width=0.7\linewidth]{comparacion/roc_test_orig.png}
\caption{Curvas ROC macro en los originales de prueba.}
\label{fig:roc}
\end{figure}

\begin{figure}[htbp]
\centering
\includegraphics[width=\linewidth]{comparacion/barras_test.png}
\caption{Exactitud y F1 macro por modelo en la prueba de originales y en la prueba completa.}
\label{fig:barras}
\end{figure}

\paragraph{Lectura de los resultados.}
\begin{itemize}
  \item \textbf{Ningún modelo alcanza el 89\,\%.} La ResNet-50 acierta 683 de 768 imágenes (88,9\,\%; IC de Wilson al 95\,\%: 86,5--91,0\,\%). Las diferencias entre los cuatro modelos están dentro de ese margen.
  \item \textbf{La exactitud está favorecida por la clase mayoritaria.} Decir siempre Protrusión daría 73,2\,\%. Por eso se usa el F1 macro como criterio.
  \item \textbf{Degeneración es la clase débil:} como mucho se detecta la mitad de los 42 casos, y muchos se confunden con Protrusión.
  \item \textbf{Hay sobreajuste:} los modelos superan el 97\,\% en entrenamiento y quedan cerca del 89\,\% en prueba. En la prueba completa, que pesa igual las tres clases, la exactitud baja a 74--78\,\%.
  \item \textbf{Posible atajo por nivel:} la ResNet-50 acierta 97,4\,\% en L4-L5 y 78,4\,\% en L1-L2, en coherencia con el sesgo de la Tabla~\ref{tab:nivel}.
  \item \textbf{El ensamble no reemplaza a la ResNet-50.} Mejora ROC-AUC y MSE, pero detecta 17 de 42 casos de Degeneración frente a 21, y su F1 macro de validación es 0,822 frente a 0,824.
\end{itemize}

\subsection{Publicación de los modelos}
Los pesos (\arch{yolov8s_cls.pt}, \arch{resnet50.pt}, \arch{svm.joblib} y \arch{random_forest.joblib}, unos 121~MB) están en la \emph{Release} \arch{modelos-v1.0} de GitHub. El script \arch{scripts/descargar_modelos.py} los descarga y verifica su SHA-256. El detector de discos se versiona con Git LFS.

% Cifras: docs/resultados/ y docs/modelo_matematico/README.md
% =====================================================================
\section{Entrenamiento y evaluación de los modelos}
\label{sec:entrenamiento}

Los cuatro modelos se entrenaron con el dataset de la Sección~\ref{sec:dataset}, con la misma división por paciente y la semilla 42. Cada script comprueba que ningún paciente esté en dos particiones y que el dataset coincida con \arch{manifest_35k.csv}. Los scripts están en \arch{scripts/entrenamiento/}, los resultados en \arch{docs/resultados/} y las ecuaciones de cada modelo en \arch{docs/modelo_matematico/README.md}.

\subsection{Estructura del sistema}
La Figura~\ref{fig:estructura} muestra la arquitectura. YOLOv8s-cls sigue un camino independiente. La ResNet-50 es la base de los otros tres modelos: sus 2\,048 características alimentan su propia capa final (modelo~2) y, una vez estandarizadas y reducidas a 256 componentes con PCA, al SVM (modelo~3) y al Random Forest (modelo~4). El ensamble promedia las salidas de los cuatro; no se entrena y se reporta como resultado adicional, no como quinto modelo.

\begin{figure}[htbp]
\centering
\includegraphics[width=\linewidth]{09_estructura_sistema.png}
\caption{Estructura del sistema: los modelos 2, 3 y 4 comparten las características de la ResNet-50.}
\label{fig:estructura}
\end{figure}

\subsection{Protocolo de entrenamiento}
El orden fue el siguiente, y cada etapa usó solo lo aprendido en las anteriores:
\begin{enumerate}
  \item Ajuste de YOLOv8s-cls con entrenamiento, eligiendo la época con validación.
  \item Ajuste de la ResNet-50 con entrenamiento, eligiendo la época con validación.
  \item Ajuste del escalador y del PCA con las características de entrenamiento.
  \item Ajuste del SVM y del Random Forest, eligiendo sus hiperparámetros con validación.
  \item Evaluación única sobre la partición de prueba, con la configuración ya cerrada.
\end{enumerate}
La Tabla~\ref{tab:config} resume la configuración de las dos redes. Ambas parten de pesos de ImageNet y usan AdamW con tasa de aprendizaje en coseno.

\begin{table}[htbp]
\centering\small
\caption{Configuración de entrenamiento de las redes.}
\label{tab:config}
\begin{tabular}{p{0.30\linewidth}p{0.29\linewidth}p{0.29\linewidth}}
\toprule
 & YOLOv8s-cls & ResNet-50 \\
\midrule
Parámetros & 5\,084\,579 & 23\,514\,179 (225\,344 congelados) \\
Tasa inicial / mínima & $5\cdot10^{-4}$ / $5\cdot10^{-5}$ & $10^{-4}$ / 0 \\
Épocas máximas / paciencia & 15 / 4 & 10 / 4 \\
Capas congeladas & primeros 4 módulos & \arch{conv1}, \arch{bn1}, \arch{layer1} \\
Suavizado de etiquetas & no & $\varepsilon = 0{,}1$ \\
Criterio de la época & exactitud en validación & F1 macro en validación \\
Época elegida / corridas & 2 / 6 & 2 / 6 \\
\bottomrule
\end{tabular}
\end{table}

\paragraph{SVM y Random Forest.} Con validación se eligió $C = 0{,}1$ para el SVM (núcleo RBF), que quedó con 1\,684 vectores de soporte de 24\,501 imágenes, y una profundidad máxima de 20 para el Random Forest de 300 árboles (Figura~\ref{fig:hiper}). Con PCA, 256 componentes explican el 80,0\,\% de la varianza de las características ajustadas.

\begin{figure}[htbp]
\centering
\includegraphics[width=\linewidth]{08_busqueda_hiperparametros.png}
\caption{Búsqueda de hiperparámetros con la partición de validación.}
\label{fig:hiper}
\end{figure}

\subsection{Evidencia del entrenamiento}
Las Figuras~\ref{fig:curvas1} y~\ref{fig:curvas2} son las curvas registradas. En las dos redes la mejor época es la 2: desde ahí la pérdida de entrenamiento sigue bajando y la de validación empeora. En la ResNet-50, el F1 macro de validación llegó a 0,824 en la época 2 y bajó en las siguientes, y la parada temprana detuvo el proceso en la época 6.

\begin{figure}[htbp]
\centering
\includegraphics[width=0.85\linewidth]{modelo1_yolov8s_cls/curvas_entrenamiento.png}
\caption{Curvas de entrenamiento de YOLOv8s-cls.}
\label{fig:curvas1}
\end{figure}

\begin{figure}[htbp]
\centering
\includegraphics[width=0.85\linewidth]{modelo2_resnet50/curvas_entrenamiento.png}
\caption{Curvas de entrenamiento de la ResNet-50.}
\label{fig:curvas2}
\end{figure}

\subsection{Ajustes probados en validación}
Se probaron variantes para acercarse al 89\,\% de exactitud que pide la rúbrica, todas decididas con validación (registro en \arch{docs/resultados/experimentos_validacion.md}). El cambio de mayor efecto fue sustituir el extractor de ImageNet por la ResNet-50 ajustada: el F1 de validación subió de 0,65 a 0,80 en el SVM y de 0,60 a 0,79 en el Random Forest. A cambio, los modelos 3 y 4 dependen del modelo~2. Las demás variantes se estancaron en una exactitud de 86--87,5\,\% en validación.

\subsection{Resultados}
\label{sec:resultados}
La Tabla~\ref{tab:resultados} recoge los resultados en la prueba de originales. El registro de la Tabla~\ref{tab:registro} compara entrenamiento y prueba. Las Figuras~\ref{fig:matrices}, \ref{fig:roc} y~\ref{fig:barras} muestran las matrices de confusión, las curvas ROC y la comparación con el conjunto de prueba completo.

\begin{table}[htbp]
\centering\small
\caption{Resultados en los originales de prueba (768 imágenes: 42 Degeneración, 164 Normal, 562 Protrusión).}
\label{tab:resultados}
\begin{tabular}{lrrrrrr}
\toprule
 & & & & \multicolumn{3}{c}{Recall} \\
\cmidrule(lr){5-7}
Modelo & Exactitud & F1 macro & ROC-AUC & Deg. & Normal & Prot. \\
\midrule
1. YOLOv8s-cls        & 87,8\,\% & 0,750 & 0,931 & 0,50 & 0,76 & 0,94 \\
\textbf{2. ResNet-50} & \textbf{88,9\,\%} & \textbf{0,769} & 0,919 & 0,50 & 0,81 & 0,94 \\
3. SVM (RBF)          & 88,4\,\% & 0,748 & 0,899 & 0,43 & 0,77 & 0,95 \\
4. Random Forest      & 88,4\,\% & 0,746 & 0,942 & 0,40 & 0,77 & 0,95 \\
\midrule
Extra: ensamble       & 88,9\,\% & 0,745 & 0,947 & 0,40 & 0,81 & 0,95 \\
\bottomrule
\end{tabular}
\end{table}

\begin{figure}[htbp]
\centering
\includegraphics[width=\linewidth]{comparacion/matrices_confusion_test_orig.png}
\caption{Matrices de confusión de los cuatro modelos en los originales de prueba.}
\label{fig:matrices}
\end{figure}

\begin{figure}[htbp]
\centering
\includegraphics[width=0.7\linewidth]{comparacion/roc_test_orig.png}
\caption{Curvas ROC macro en los originales de prueba.}
\label{fig:roc}
\end{figure}

\begin{figure}[htbp]
\centering
\includegraphics[width=\linewidth]{comparacion/barras_test.png}
\caption{Exactitud y F1 macro por modelo en la prueba de originales y en la prueba completa.}
\label{fig:barras}
\end{figure}

\paragraph{Lectura de los resultados.}
\begin{itemize}
  \item \textbf{Ningún modelo alcanza el 89\,\%.} La ResNet-50 acierta 683 de 768 imágenes (88,9\,\%; IC de Wilson al 95\,\%: 86,5--91,0\,\%). Las diferencias entre los cuatro modelos están dentro de ese margen.
  \item \textbf{La exactitud está favorecida por la clase mayoritaria.} Decir siempre Protrusión daría 73,2\,\%. Por eso se usa el F1 macro como criterio.
  \item \textbf{Degeneración es la clase débil:} como mucho se detecta la mitad de los 42 casos, y muchos se confunden con Protrusión.
  \item \textbf{Hay sobreajuste:} los modelos superan el 97\,\% en entrenamiento y quedan cerca del 89\,\% en prueba. En la prueba completa, que pesa igual las tres clases, la exactitud baja a 74--78\,\%.
  \item \textbf{Posible atajo por nivel:} la ResNet-50 acierta 97,4\,\% en L4-L5 y 78,4\,\% en L1-L2, en coherencia con el sesgo de la Tabla~\ref{tab:nivel}.
  \item \textbf{El ensamble no reemplaza a la ResNet-50.} Mejora ROC-AUC y MSE, pero detecta 17 de 42 casos de Degeneración frente a 21, y su F1 macro de validación es 0,822 frente a 0,824.
\end{itemize}

\subsection{Publicación de los modelos}
Los pesos (\arch{yolov8s_cls.pt}, \arch{resnet50.pt}, \arch{svm.joblib} y \arch{random_forest.joblib}, unos 121~MB) están en la \emph{Release} \arch{modelos-v1.0} de GitHub. El script \arch{scripts/descargar_modelos.py} los descarga y verifica su SHA-256. El detector de discos se versiona con Git LFS.

% Cifras: docs/resultados/ y docs/modelo_matematico/README.md
% =====================================================================
\section{Entrenamiento y evaluación de los modelos}
\label{sec:entrenamiento}

Los cuatro modelos se entrenaron con el dataset de la Sección~\ref{sec:dataset}, con la misma división por paciente y la semilla 42. Cada script comprueba que ningún paciente esté en dos particiones y que el dataset coincida con \arch{manifest_35k.csv}. Los scripts están en \arch{scripts/entrenamiento/}, los resultados en \arch{docs/resultados/} y las ecuaciones de cada modelo en \arch{docs/modelo_matematico/README.md}.

\subsection{Estructura del sistema}
La Figura~\ref{fig:estructura} muestra la arquitectura. YOLOv8s-cls sigue un camino independiente. La ResNet-50 es la base de los otros tres modelos: sus 2\,048 características alimentan su propia capa final (modelo~2) y, una vez estandarizadas y reducidas a 256 componentes con PCA, al SVM (modelo~3) y al Random Forest (modelo~4). El ensamble promedia las salidas de los cuatro; no se entrena y se reporta como resultado adicional, no como quinto modelo.

\begin{figure}[htbp]
\centering
\includegraphics[width=\linewidth]{09_estructura_sistema.png}
\caption{Estructura del sistema: los modelos 2, 3 y 4 comparten las características de la ResNet-50.}
\label{fig:estructura}
\end{figure}

\subsection{Protocolo de entrenamiento}
El orden fue el siguiente, y cada etapa usó solo lo aprendido en las anteriores:
\begin{enumerate}
  \item Ajuste de YOLOv8s-cls con entrenamiento, eligiendo la época con validación.
  \item Ajuste de la ResNet-50 con entrenamiento, eligiendo la época con validación.
  \item Ajuste del escalador y del PCA con las características de entrenamiento.
  \item Ajuste del SVM y del Random Forest, eligiendo sus hiperparámetros con validación.
  \item Evaluación única sobre la partición de prueba, con la configuración ya cerrada.
\end{enumerate}
La Tabla~\ref{tab:config} resume la configuración de las dos redes. Ambas parten de pesos de ImageNet y usan AdamW con tasa de aprendizaje en coseno.

\begin{table}[htbp]
\centering\small
\caption{Configuración de entrenamiento de las redes.}
\label{tab:config}
\begin{tabular}{p{0.30\linewidth}p{0.29\linewidth}p{0.29\linewidth}}
\toprule
 & YOLOv8s-cls & ResNet-50 \\
\midrule
Parámetros & 5\,084\,579 & 23\,514\,179 (225\,344 congelados) \\
Tasa inicial / mínima & $5\cdot10^{-4}$ / $5\cdot10^{-5}$ & $10^{-4}$ / 0 \\
Épocas máximas / paciencia & 15 / 4 & 10 / 4 \\
Capas congeladas & primeros 4 módulos & \arch{conv1}, \arch{bn1}, \arch{layer1} \\
Suavizado de etiquetas & no & $\varepsilon = 0{,}1$ \\
Criterio de la época & exactitud en validación & F1 macro en validación \\
Época elegida / corridas & 2 / 6 & 2 / 6 \\
\bottomrule
\end{tabular}
\end{table}

\paragraph{SVM y Random Forest.} Con validación se eligió $C = 0{,}1$ para el SVM (núcleo RBF), que quedó con 1\,684 vectores de soporte de 24\,501 imágenes, y una profundidad máxima de 20 para el Random Forest de 300 árboles (Figura~\ref{fig:hiper}). Con PCA, 256 componentes explican el 80,0\,\% de la varianza de las características ajustadas.

\begin{figure}[htbp]
\centering
\includegraphics[width=\linewidth]{08_busqueda_hiperparametros.png}
\caption{Búsqueda de hiperparámetros con la partición de validación.}
\label{fig:hiper}
\end{figure}

\subsection{Evidencia del entrenamiento}
Las Figuras~\ref{fig:curvas1} y~\ref{fig:curvas2} son las curvas registradas. En las dos redes la mejor época es la 2: desde ahí la pérdida de entrenamiento sigue bajando y la de validación empeora. En la ResNet-50, el F1 macro de validación llegó a 0,824 en la época 2 y bajó en las siguientes, y la parada temprana detuvo el proceso en la época 6.

\begin{figure}[htbp]
\centering
\includegraphics[width=0.85\linewidth]{modelo1_yolov8s_cls/curvas_entrenamiento.png}
\caption{Curvas de entrenamiento de YOLOv8s-cls.}
\label{fig:curvas1}
\end{figure}

\begin{figure}[htbp]
\centering
\includegraphics[width=0.85\linewidth]{modelo2_resnet50/curvas_entrenamiento.png}
\caption{Curvas de entrenamiento de la ResNet-50.}
\label{fig:curvas2}
\end{figure}

\subsection{Ajustes probados en validación}
Se probaron variantes para acercarse al 89\,\% de exactitud que pide la rúbrica, todas decididas con validación (registro en \arch{docs/resultados/experimentos_validacion.md}). El cambio de mayor efecto fue sustituir el extractor de ImageNet por la ResNet-50 ajustada: el F1 de validación subió de 0,65 a 0,80 en el SVM y de 0,60 a 0,79 en el Random Forest. A cambio, los modelos 3 y 4 dependen del modelo~2. Las demás variantes se estancaron en una exactitud de 86--87,5\,\% en validación.

\subsection{Resultados}
\label{sec:resultados}
La Tabla~\ref{tab:resultados} recoge los resultados en la prueba de originales. El registro de la Tabla~\ref{tab:registro} compara entrenamiento y prueba. Las Figuras~\ref{fig:matrices}, \ref{fig:roc} y~\ref{fig:barras} muestran las matrices de confusión, las curvas ROC y la comparación con el conjunto de prueba completo.

\begin{table}[htbp]
\centering\small
\caption{Resultados en los originales de prueba (768 imágenes: 42 Degeneración, 164 Normal, 562 Protrusión).}
\label{tab:resultados}
\begin{tabular}{lrrrrrr}
\toprule
 & & & & \multicolumn{3}{c}{Recall} \\
\cmidrule(lr){5-7}
Modelo & Exactitud & F1 macro & ROC-AUC & Deg. & Normal & Prot. \\
\midrule
1. YOLOv8s-cls        & 87,8\,\% & 0,750 & 0,931 & 0,50 & 0,76 & 0,94 \\
\textbf{2. ResNet-50} & \textbf{88,9\,\%} & \textbf{0,769} & 0,919 & 0,50 & 0,81 & 0,94 \\
3. SVM (RBF)          & 88,4\,\% & 0,748 & 0,899 & 0,43 & 0,77 & 0,95 \\
4. Random Forest      & 88,4\,\% & 0,746 & 0,942 & 0,40 & 0,77 & 0,95 \\
\midrule
Extra: ensamble       & 88,9\,\% & 0,745 & 0,947 & 0,40 & 0,81 & 0,95 \\
\bottomrule
\end{tabular}
\end{table}

\begin{figure}[htbp]
\centering
\includegraphics[width=\linewidth]{comparacion/matrices_confusion_test_orig.png}
\caption{Matrices de confusión de los cuatro modelos en los originales de prueba.}
\label{fig:matrices}
\end{figure}

\begin{figure}[htbp]
\centering
\includegraphics[width=0.7\linewidth]{comparacion/roc_test_orig.png}
\caption{Curvas ROC macro en los originales de prueba.}
\label{fig:roc}
\end{figure}

\begin{figure}[htbp]
\centering
\includegraphics[width=\linewidth]{comparacion/barras_test.png}
\caption{Exactitud y F1 macro por modelo en la prueba de originales y en la prueba completa.}
\label{fig:barras}
\end{figure}

\paragraph{Lectura de los resultados.}
\begin{itemize}
  \item \textbf{Ningún modelo alcanza el 89\,\%.} La ResNet-50 acierta 683 de 768 imágenes (88,9\,\%; IC de Wilson al 95\,\%: 86,5--91,0\,\%). Las diferencias entre los cuatro modelos están dentro de ese margen.
  \item \textbf{La exactitud está favorecida por la clase mayoritaria.} Decir siempre Protrusión daría 73,2\,\%. Por eso se usa el F1 macro como criterio.
  \item \textbf{Degeneración es la clase débil:} como mucho se detecta la mitad de los 42 casos, y muchos se confunden con Protrusión.
  \item \textbf{Hay sobreajuste:} los modelos superan el 97\,\% en entrenamiento y quedan cerca del 89\,\% en prueba. En la prueba completa, que pesa igual las tres clases, la exactitud baja a 74--78\,\%.
  \item \textbf{Posible atajo por nivel:} la ResNet-50 acierta 97,4\,\% en L4-L5 y 78,4\,\% en L1-L2, en coherencia con el sesgo de la Tabla~\ref{tab:nivel}.
  \item \textbf{El ensamble no reemplaza a la ResNet-50.} Mejora ROC-AUC y MSE, pero detecta 17 de 42 casos de Degeneración frente a 21, y su F1 macro de validación es 0,822 frente a 0,824.
\end{itemize}

\subsection{Publicación de los modelos}
Los pesos (\arch{yolov8s_cls.pt}, \arch{resnet50.pt}, \arch{svm.joblib} y \arch{random_forest.joblib}, unos 121~MB) están en la \emph{Release} \arch{modelos-v1.0} de GitHub. El script \arch{scripts/descargar_modelos.py} los descarga y verifica su SHA-256. El detector de discos se versiona con Git LFS.

% =====================================================================
% Seccion: MVP de demostracion (Hito 2)
% Fragmento SIN preambulo: % =====================================================================
% Seccion: MVP de demostracion (Hito 2)
% Fragmento SIN preambulo: % =====================================================================
% Seccion: MVP de demostracion (Hito 2)
% Fragmento SIN preambulo: \input{secciones/mvp}
% Capturas tomadas en el commit 58b1a9e, ultimo con el MVP disponible.
% =====================================================================
\section{MVP de demostración}
\label{sec:mvp}

Para el Hito~2 se construyó un MVP local, que se clona y se ejecuta con un solo comando: un backend FastAPI con los cuatro modelos y una interfaz estática en \arch{/mvp}. Su objetivo era demostrar la simulación de resultados y la trazabilidad de una predicción. Las capturas de esta sección se tomaron en el commit \arch{58b1a9e}. Tras presentar el hito, el MVP se retiró del proyecto (Sección~\ref{sec:arreglos}); su código sigue disponible en el historial de Git.

El MVP tenía cinco pestañas: caso a caso, subir imagen, simulación, modelos y métricas, y estudio de resonancia. Esta última la agregó Juan Manuel Tengonó y fue la base de la aplicación final.

\paragraph{Caso a caso y trazabilidad.} Se elige uno de los 45 casos de prueba (15 por clase) y se ejecutan los cuatro modelos. La Figura~\ref{fig:mvpcaso} muestra un error real: en \arch{caso_04}, con etiqueta Degeneración, los cuatro modelos predicen Protrusión. La Figura~\ref{fig:mvptraza} muestra el recorrido de esa imagen con valores de la ejecución: la ResNet-50 se calcula una sola vez y sus características pasan a los modelos 2, 3 y 4.

\begin{figure}[htbp]
\centering
\includegraphics[width=\linewidth]{mvp_caso.jpg}
\caption{Pestaña Caso a caso del MVP: resultado de los cuatro modelos sobre un caso de prueba.}
\label{fig:mvpcaso}
\end{figure}

\begin{figure}[htbp]
\centering
\includegraphics[width=\linewidth]{mvp_trazabilidad.jpg}
\caption{Trazabilidad de una imagen por los cuatro modelos, con tiempos y valores intermedios.}
\label{fig:mvptraza}
\end{figure}

La pestaña \emph{Simulación} se describe aparte, en la Sección~\ref{sec:simulacion}.

\paragraph{Modelos y métricas.} La pestaña reunía las métricas oficiales, las matrices de confusión, las curvas ROC y el SHA-256 de cada archivo de pesos (Figuras~\ref{fig:mvpmetricas} y~\ref{fig:mvproc}). Esta información ya no está en la aplicación: se conserva en la documentación y en la Sección~\ref{sec:resultados}.

\begin{figure}[htbp]
\centering
\includegraphics[width=\linewidth]{mvp_metricas.jpg}
\caption{Pestaña Modelos y métricas del MVP.}
\label{fig:mvpmetricas}
\end{figure}

\begin{figure}[htbp]
\centering
\includegraphics[width=\linewidth]{mvp_roc.jpg}
\caption{Curvas ROC y barras de exactitud y F1 macro en el MVP.}
\label{fig:mvproc}
\end{figure}

% Capturas tomadas en el commit 58b1a9e, ultimo con el MVP disponible.
% =====================================================================
\section{MVP de demostración}
\label{sec:mvp}

Para el Hito~2 se construyó un MVP local, que se clona y se ejecuta con un solo comando: un backend FastAPI con los cuatro modelos y una interfaz estática en \arch{/mvp}. Su objetivo era demostrar la simulación de resultados y la trazabilidad de una predicción. Las capturas de esta sección se tomaron en el commit \arch{58b1a9e}. Tras presentar el hito, el MVP se retiró del proyecto (Sección~\ref{sec:arreglos}); su código sigue disponible en el historial de Git.

El MVP tenía cinco pestañas: caso a caso, subir imagen, simulación, modelos y métricas, y estudio de resonancia. Esta última la agregó Juan Manuel Tengonó y fue la base de la aplicación final.

\paragraph{Caso a caso y trazabilidad.} Se elige uno de los 45 casos de prueba (15 por clase) y se ejecutan los cuatro modelos. La Figura~\ref{fig:mvpcaso} muestra un error real: en \arch{caso_04}, con etiqueta Degeneración, los cuatro modelos predicen Protrusión. La Figura~\ref{fig:mvptraza} muestra el recorrido de esa imagen con valores de la ejecución: la ResNet-50 se calcula una sola vez y sus características pasan a los modelos 2, 3 y 4.

\begin{figure}[htbp]
\centering
\includegraphics[width=\linewidth]{mvp_caso.jpg}
\caption{Pestaña Caso a caso del MVP: resultado de los cuatro modelos sobre un caso de prueba.}
\label{fig:mvpcaso}
\end{figure}

\begin{figure}[htbp]
\centering
\includegraphics[width=\linewidth]{mvp_trazabilidad.jpg}
\caption{Trazabilidad de una imagen por los cuatro modelos, con tiempos y valores intermedios.}
\label{fig:mvptraza}
\end{figure}

La pestaña \emph{Simulación} se describe aparte, en la Sección~\ref{sec:simulacion}.

\paragraph{Modelos y métricas.} La pestaña reunía las métricas oficiales, las matrices de confusión, las curvas ROC y el SHA-256 de cada archivo de pesos (Figuras~\ref{fig:mvpmetricas} y~\ref{fig:mvproc}). Esta información ya no está en la aplicación: se conserva en la documentación y en la Sección~\ref{sec:resultados}.

\begin{figure}[htbp]
\centering
\includegraphics[width=\linewidth]{mvp_metricas.jpg}
\caption{Pestaña Modelos y métricas del MVP.}
\label{fig:mvpmetricas}
\end{figure}

\begin{figure}[htbp]
\centering
\includegraphics[width=\linewidth]{mvp_roc.jpg}
\caption{Curvas ROC y barras de exactitud y F1 macro en el MVP.}
\label{fig:mvproc}
\end{figure}

% Capturas tomadas en el commit 58b1a9e, ultimo con el MVP disponible.
% =====================================================================
\section{MVP de demostración}
\label{sec:mvp}

Para el Hito~2 se construyó un MVP local, que se clona y se ejecuta con un solo comando: un backend FastAPI con los cuatro modelos y una interfaz estática en \arch{/mvp}. Su objetivo era demostrar la simulación de resultados y la trazabilidad de una predicción. Las capturas de esta sección se tomaron en el commit \arch{58b1a9e}. Tras presentar el hito, el MVP se retiró del proyecto (Sección~\ref{sec:arreglos}); su código sigue disponible en el historial de Git.

El MVP tenía cinco pestañas: caso a caso, subir imagen, simulación, modelos y métricas, y estudio de resonancia. Esta última la agregó Juan Manuel Tengonó y fue la base de la aplicación final.

\paragraph{Caso a caso y trazabilidad.} Se elige uno de los 45 casos de prueba (15 por clase) y se ejecutan los cuatro modelos. La Figura~\ref{fig:mvpcaso} muestra un error real: en \arch{caso_04}, con etiqueta Degeneración, los cuatro modelos predicen Protrusión. La Figura~\ref{fig:mvptraza} muestra el recorrido de esa imagen con valores de la ejecución: la ResNet-50 se calcula una sola vez y sus características pasan a los modelos 2, 3 y 4.

\begin{figure}[htbp]
\centering
\includegraphics[width=\linewidth]{mvp_caso.jpg}
\caption{Pestaña Caso a caso del MVP: resultado de los cuatro modelos sobre un caso de prueba.}
\label{fig:mvpcaso}
\end{figure}

\begin{figure}[htbp]
\centering
\includegraphics[width=\linewidth]{mvp_trazabilidad.jpg}
\caption{Trazabilidad de una imagen por los cuatro modelos, con tiempos y valores intermedios.}
\label{fig:mvptraza}
\end{figure}

La pestaña \emph{Simulación} se describe aparte, en la Sección~\ref{sec:simulacion}.

\paragraph{Modelos y métricas.} La pestaña reunía las métricas oficiales, las matrices de confusión, las curvas ROC y el SHA-256 de cada archivo de pesos (Figuras~\ref{fig:mvpmetricas} y~\ref{fig:mvproc}). Esta información ya no está en la aplicación: se conserva en la documentación y en la Sección~\ref{sec:resultados}.

\begin{figure}[htbp]
\centering
\includegraphics[width=\linewidth]{mvp_metricas.jpg}
\caption{Pestaña Modelos y métricas del MVP.}
\label{fig:mvpmetricas}
\end{figure}

\begin{figure}[htbp]
\centering
\includegraphics[width=\linewidth]{mvp_roc.jpg}
\caption{Curvas ROC y barras de exactitud y F1 macro en el MVP.}
\label{fig:mvproc}
\end{figure}

% =====================================================================
% Seccion: Simulacion de resultados
% Fragmento SIN preambulo: % =====================================================================
% Seccion: Simulacion de resultados
% Fragmento SIN preambulo: % =====================================================================
% Seccion: Simulacion de resultados
% Fragmento SIN preambulo: \input{secciones/simulacion}
% Casos: data/demo/casos.csv del commit 58b1a9e. Cifras: nueva corrida con
% la version final de la API (/api/predecir), 2026-10-30.
% =====================================================================
\section{Simulación de resultados}
\label{sec:simulacion}

La simulación comprueba de punta a punta que el sistema desplegado se comporta igual que los modelos evaluados. Toma un conjunto fijo de casos, los pasa uno por uno por los cuatro modelos y compara cada predicción con la que quedó registrada en la evaluación oficial (\arch{puntajes.npz}, en \arch{docs/resultados/}).

\subsection{Casos y procedimiento}
Se usaron 45 imágenes originales de la partición de prueba, 15 por clase, repartidas entre los cinco niveles y elegidas con la semilla 42 por el script \arch{seleccionar_casos_demo.py}. Ningún modelo las vio al entrenar. Por cada caso se ejecutan los cuatro modelos y el ensamble, se actualiza la matriz de confusión y se marca si cada predicción coincide con la registrada. En el MVP esto se veía en vivo (Figuras~\ref{fig:mvpsim} y~\ref{fig:mvpsimfin}).

\begin{figure}[htbp]
\centering
\includegraphics[width=\linewidth]{mvp_simulacion.jpg}
\caption{Simulación en curso en el MVP (29 de 45 casos).}
\label{fig:mvpsim}
\end{figure}

\begin{figure}[htbp]
\centering
\includegraphics[width=\linewidth]{mvp_simulacion_fin.jpg}
\caption{Últimas filas de la simulación en el MVP: cada caso reproduce el registro oficial en los cuatro modelos (4/4).}
\label{fig:mvpsimfin}
\end{figure}

\subsection{Resultados}
Aunque el MVP se retiró, la simulación se repitió con la versión final de la API, sobre los mismos 45 casos, en CPU (unos 11~s en total). La Tabla~\ref{tab:simulacion} recoge el resultado.

\begin{table}[htbp]
\centering\small
\caption{Simulación sobre 45 casos de prueba (15 por clase). Cada celda indica cuántos casos de la clase real recibieron cada predicción, en el orden Degeneración / Normal / Protrusión.}
\label{tab:simulacion}
\begin{tabular}{lcccrr}
\toprule
 & \multicolumn{3}{c}{Clase real} & & \\
\cmidrule(lr){2-4}
Modelo & Degeneración & Normal & Protrusión & Aciertos & Reproduce \\
\midrule
1. YOLOv8s-cls   & 3 / 2 / 10 & 0 / 8 / 7  & 1 / 2 / 12 & 23 de 45 & 45 de 45 \\
2. ResNet-50     & 4 / 2 / 9  & 0 / 10 / 5 & 0 / 2 / 13 & 27 de 45 & 45 de 45 \\
3. SVM (RBF)     & 3 / 2 / 10 & 0 / 10 / 5 & 0 / 2 / 13 & 26 de 45 & 45 de 45 \\
4. Random Forest & 3 / 2 / 10 & 0 / 10 / 5 & 0 / 1 / 14 & 27 de 45 & 45 de 45 \\
\midrule
Extra: ensamble  & 3 / 2 / 10 & 0 / 10 / 5 & 0 / 2 / 13 & 26 de 45 & -- \\
\bottomrule
\end{tabular}
\end{table}

\paragraph{Lectura.}
\begin{itemize}
  \item \textbf{Reproducibilidad.} Las 45 predicciones de cada modelo coinciden con la evaluación oficial, es decir, 180 de 180. La aplicación ejecuta los mismos modelos que se evaluaron.
  \item \textbf{La exactitud de la simulación no es la oficial.} Los aciertos (51--60\,\%) no se comparan con el 88,9\,\% de la Tabla~\ref{tab:resultados}: aquí cada clase pesa un tercio, mientras que en la prueba de originales el 73\,\% es Protrusión. Estos casos sirven para verificar el sistema, no para medir su desempeño.
  \item \textbf{Se repite el patrón de error.} Como máximo se aciertan 4 de 15 casos de Degeneración, y entre 9 y 10 se confunden con Protrusión. Es la misma debilidad de la Sección~\ref{sec:resultados}, más visible porque la clase pesa más.
  \item \textbf{Acuerdo entre modelos.} Los cuatro coinciden en 35 de los 45 casos, lo que confirma que no son opiniones independientes: tres de ellos comparten las características de la ResNet-50.
\end{itemize}

% Casos: data/demo/casos.csv del commit 58b1a9e. Cifras: nueva corrida con
% la version final de la API (/api/predecir), 2026-10-30.
% =====================================================================
\section{Simulación de resultados}
\label{sec:simulacion}

La simulación comprueba de punta a punta que el sistema desplegado se comporta igual que los modelos evaluados. Toma un conjunto fijo de casos, los pasa uno por uno por los cuatro modelos y compara cada predicción con la que quedó registrada en la evaluación oficial (\arch{puntajes.npz}, en \arch{docs/resultados/}).

\subsection{Casos y procedimiento}
Se usaron 45 imágenes originales de la partición de prueba, 15 por clase, repartidas entre los cinco niveles y elegidas con la semilla 42 por el script \arch{seleccionar_casos_demo.py}. Ningún modelo las vio al entrenar. Por cada caso se ejecutan los cuatro modelos y el ensamble, se actualiza la matriz de confusión y se marca si cada predicción coincide con la registrada. En el MVP esto se veía en vivo (Figuras~\ref{fig:mvpsim} y~\ref{fig:mvpsimfin}).

\begin{figure}[htbp]
\centering
\includegraphics[width=\linewidth]{mvp_simulacion.jpg}
\caption{Simulación en curso en el MVP (29 de 45 casos).}
\label{fig:mvpsim}
\end{figure}

\begin{figure}[htbp]
\centering
\includegraphics[width=\linewidth]{mvp_simulacion_fin.jpg}
\caption{Últimas filas de la simulación en el MVP: cada caso reproduce el registro oficial en los cuatro modelos (4/4).}
\label{fig:mvpsimfin}
\end{figure}

\subsection{Resultados}
Aunque el MVP se retiró, la simulación se repitió con la versión final de la API, sobre los mismos 45 casos, en CPU (unos 11~s en total). La Tabla~\ref{tab:simulacion} recoge el resultado.

\begin{table}[htbp]
\centering\small
\caption{Simulación sobre 45 casos de prueba (15 por clase). Cada celda indica cuántos casos de la clase real recibieron cada predicción, en el orden Degeneración / Normal / Protrusión.}
\label{tab:simulacion}
\begin{tabular}{lcccrr}
\toprule
 & \multicolumn{3}{c}{Clase real} & & \\
\cmidrule(lr){2-4}
Modelo & Degeneración & Normal & Protrusión & Aciertos & Reproduce \\
\midrule
1. YOLOv8s-cls   & 3 / 2 / 10 & 0 / 8 / 7  & 1 / 2 / 12 & 23 de 45 & 45 de 45 \\
2. ResNet-50     & 4 / 2 / 9  & 0 / 10 / 5 & 0 / 2 / 13 & 27 de 45 & 45 de 45 \\
3. SVM (RBF)     & 3 / 2 / 10 & 0 / 10 / 5 & 0 / 2 / 13 & 26 de 45 & 45 de 45 \\
4. Random Forest & 3 / 2 / 10 & 0 / 10 / 5 & 0 / 1 / 14 & 27 de 45 & 45 de 45 \\
\midrule
Extra: ensamble  & 3 / 2 / 10 & 0 / 10 / 5 & 0 / 2 / 13 & 26 de 45 & -- \\
\bottomrule
\end{tabular}
\end{table}

\paragraph{Lectura.}
\begin{itemize}
  \item \textbf{Reproducibilidad.} Las 45 predicciones de cada modelo coinciden con la evaluación oficial, es decir, 180 de 180. La aplicación ejecuta los mismos modelos que se evaluaron.
  \item \textbf{La exactitud de la simulación no es la oficial.} Los aciertos (51--60\,\%) no se comparan con el 88,9\,\% de la Tabla~\ref{tab:resultados}: aquí cada clase pesa un tercio, mientras que en la prueba de originales el 73\,\% es Protrusión. Estos casos sirven para verificar el sistema, no para medir su desempeño.
  \item \textbf{Se repite el patrón de error.} Como máximo se aciertan 4 de 15 casos de Degeneración, y entre 9 y 10 se confunden con Protrusión. Es la misma debilidad de la Sección~\ref{sec:resultados}, más visible porque la clase pesa más.
  \item \textbf{Acuerdo entre modelos.} Los cuatro coinciden en 35 de los 45 casos, lo que confirma que no son opiniones independientes: tres de ellos comparten las características de la ResNet-50.
\end{itemize}

% Casos: data/demo/casos.csv del commit 58b1a9e. Cifras: nueva corrida con
% la version final de la API (/api/predecir), 2026-10-30.
% =====================================================================
\section{Simulación de resultados}
\label{sec:simulacion}

La simulación comprueba de punta a punta que el sistema desplegado se comporta igual que los modelos evaluados. Toma un conjunto fijo de casos, los pasa uno por uno por los cuatro modelos y compara cada predicción con la que quedó registrada en la evaluación oficial (\arch{puntajes.npz}, en \arch{docs/resultados/}).

\subsection{Casos y procedimiento}
Se usaron 45 imágenes originales de la partición de prueba, 15 por clase, repartidas entre los cinco niveles y elegidas con la semilla 42 por el script \arch{seleccionar_casos_demo.py}. Ningún modelo las vio al entrenar. Por cada caso se ejecutan los cuatro modelos y el ensamble, se actualiza la matriz de confusión y se marca si cada predicción coincide con la registrada. En el MVP esto se veía en vivo (Figuras~\ref{fig:mvpsim} y~\ref{fig:mvpsimfin}).

\begin{figure}[htbp]
\centering
\includegraphics[width=\linewidth]{mvp_simulacion.jpg}
\caption{Simulación en curso en el MVP (29 de 45 casos).}
\label{fig:mvpsim}
\end{figure}

\begin{figure}[htbp]
\centering
\includegraphics[width=\linewidth]{mvp_simulacion_fin.jpg}
\caption{Últimas filas de la simulación en el MVP: cada caso reproduce el registro oficial en los cuatro modelos (4/4).}
\label{fig:mvpsimfin}
\end{figure}

\subsection{Resultados}
Aunque el MVP se retiró, la simulación se repitió con la versión final de la API, sobre los mismos 45 casos, en CPU (unos 11~s en total). La Tabla~\ref{tab:simulacion} recoge el resultado.

\begin{table}[htbp]
\centering\small
\caption{Simulación sobre 45 casos de prueba (15 por clase). Cada celda indica cuántos casos de la clase real recibieron cada predicción, en el orden Degeneración / Normal / Protrusión.}
\label{tab:simulacion}
\begin{tabular}{lcccrr}
\toprule
 & \multicolumn{3}{c}{Clase real} & & \\
\cmidrule(lr){2-4}
Modelo & Degeneración & Normal & Protrusión & Aciertos & Reproduce \\
\midrule
1. YOLOv8s-cls   & 3 / 2 / 10 & 0 / 8 / 7  & 1 / 2 / 12 & 23 de 45 & 45 de 45 \\
2. ResNet-50     & 4 / 2 / 9  & 0 / 10 / 5 & 0 / 2 / 13 & 27 de 45 & 45 de 45 \\
3. SVM (RBF)     & 3 / 2 / 10 & 0 / 10 / 5 & 0 / 2 / 13 & 26 de 45 & 45 de 45 \\
4. Random Forest & 3 / 2 / 10 & 0 / 10 / 5 & 0 / 1 / 14 & 27 de 45 & 45 de 45 \\
\midrule
Extra: ensamble  & 3 / 2 / 10 & 0 / 10 / 5 & 0 / 2 / 13 & 26 de 45 & -- \\
\bottomrule
\end{tabular}
\end{table}

\paragraph{Lectura.}
\begin{itemize}
  \item \textbf{Reproducibilidad.} Las 45 predicciones de cada modelo coinciden con la evaluación oficial, es decir, 180 de 180. La aplicación ejecuta los mismos modelos que se evaluaron.
  \item \textbf{La exactitud de la simulación no es la oficial.} Los aciertos (51--60\,\%) no se comparan con el 88,9\,\% de la Tabla~\ref{tab:resultados}: aquí cada clase pesa un tercio, mientras que en la prueba de originales el 73\,\% es Protrusión. Estos casos sirven para verificar el sistema, no para medir su desempeño.
  \item \textbf{Se repite el patrón de error.} Como máximo se aciertan 4 de 15 casos de Degeneración, y entre 9 y 10 se confunden con Protrusión. Es la misma debilidad de la Sección~\ref{sec:resultados}, más visible porque la clase pesa más.
  \item \textbf{Acuerdo entre modelos.} Los cuatro coinciden en 35 de los 45 casos, lo que confirma que no son opiniones independientes: tres de ellos comparten las características de la ResNet-50.
\end{itemize}

% =====================================================================
% Seccion: Aplicacion web final y correcciones
% Fragmento SIN preambulo: % =====================================================================
% Seccion: Aplicacion web final y correcciones
% Fragmento SIN preambulo: % =====================================================================
% Seccion: Aplicacion web final y correcciones
% Fragmento SIN preambulo: \input{secciones/aplicacion_web}
% Commits: rama feature/julian-webapp-final, integrada en main en faa42ff
% =====================================================================
\section{Aplicación web para el radiólogo}
\label{sec:web}

Sobre el MVP, Juan Manuel Tengonó desarrolló la aplicación final en \arch{webapp/frontend/clinica/}. Se ejecuta en local, en CPU y sin internet después de la primera instalación: FastAPI sirve la API y una interfaz en HTML, CSS y JavaScript sin compilación. Se arranca con \arch{webapp/iniciar_app.bat} en Windows o con \arch{webapp/iniciar_app.sh} en Linux y macOS. Se prefirió esta tecnología frente a React porque no requiere compilar nada y mantiene el arranque de un solo comando.

\subsection{Flujo de trabajo}
\begin{enumerate}
  \item \textbf{Cargar.} Recibe archivos \arch{.IMA} o \arch{.dcm}, la carpeta del resonador o un \arch{.zip} (Figura~\ref{fig:appinicio}). Convierte a PNG con el mismo convertidor del dataset e identifica la serie T2 sagital.
  \item \textbf{Cortes.} Visor con brillo, contraste, zoom y cajas del detector. El detector recomienda los cortes donde mejor ve los cinco discos, y el usuario puede cambiarlos (Figura~\ref{fig:appcortes}).
  \item \textbf{Resultados.} Conclusión en lenguaje natural y una tarjeta por disco con la clase del ensamble, en cuántos cortes se repite y cuántos votos coinciden. En cada disco el radiólogo marca si concuerda o discrepa.
  \item \textbf{Informe.} Hallazgos, impresión, recomendación y firma. Se imprime o se guarda como PDF (Figura~\ref{fig:informe}).
\end{enumerate}
El atajo \emph{Analizar una imagen} clasifica un recorte de disco o un corte completo suelto.

\begin{figure}[htbp]
\centering
\includegraphics[width=\linewidth]{app_inicio.png}
\caption{Pantalla inicial de la versión final de la aplicación.}
\label{fig:appinicio}
\end{figure}

\subsection{Rechazo de imágenes y privacidad}
Los cuatro modelos siempre eligen una de las tres clases, así que antes de predecir se filtran las imágenes. La capa~1 rechaza imágenes en color, diminutas o de proporciones raras. La capa~2 distingue recorte de corte completo, y un corte solo se acepta si el detector halla al menos tres niveles en orden. La capa~3, una distancia de Mahalanobis sobre las características, \textbf{no está activa}: requiere calibrar un umbral con \arch{scripts/calibrar_rechazo.py}, y mientras tanto la aplicación lo advierte.

Los DICOM se borran al convertirse, y el nombre, la fecha de nacimiento y los identificadores nunca se devuelven ni se registran. Los PNG se borran tras 30 minutos sin uso o al detener el servidor.

\subsection{Correcciones a la aplicación}
\label{sec:arreglos}
Al revisar la aplicación frente a lo definido en el proyecto se hallaron fallos y restos del MVP. La Tabla~\ref{tab:arreglos} los resume, y todos los cambios se integraron en \arch{main} con el commit \arch{faa42ff}. La Figura~\ref{fig:confirmar} muestra la confirmación dentro de la página, y la comparación de las Figuras~\ref{fig:apprechazo} y~\ref{fig:appinicio} muestra el aviso reforzado.

\begin{table}[htbp]
\centering\small
\caption{Correcciones hechas a la aplicación y commit de cada una.}
\label{tab:arreglos}
\begin{tabular}{p{0.36\linewidth}p{0.40\linewidth}p{0.13\linewidth}}
\toprule
Problema & Corrección & Commit \\
\midrule
Repetir el análisis podía duplicar resultados o mezclar uno anterior & Un contador cancela los análisis atrasados y el cambio de serie reinicia el estado & \arch{b5e848f} \\
Los estudios vencidos solo se borraban al arrancar & Limpieza periódica cada minuto y borrado total al detener el servidor & \arch{35ab7e5} \\
Los avisos no explicaban los límites de la herramienta & Avisos sobre etiquetas, confianza orientativa y filtro sin calibrar & \arch{13d05f3} \\
Convivían el MVP, código sin uso y rutas que la aplicación no llama & Retiro del MVP, los casos de demostración, CORS y las rutas sobrantes & \arch{1360698} \\
Scripts y entorno seguían con el nombre del MVP & Paso a \arch{iniciar_app} y \arch{.venv-app} & \arch{be46285} \\
Borrar un estudio abría la ventana nativa del navegador & Diálogo de confirmación dentro de la página & \arch{8ba1c36} \\
El cierre atrasado de un diálogo consumía la confirmación siguiente y el estudio no se borraba & Se ignora ese cierre mientras el diálogo sigue abierto & \arch{173f44f} \\
Una confirmación cancelada quedaba pendiente y provocaba un doble borrado & Una sola confirmación pendiente a la vez & \arch{2e87449} \\
Documentación y carpetas desactualizadas & README al día; retiro de \arch{assets/} y \arch{data/external/} & \arch{0cef456}, \arch{cefe20c} \\
\bottomrule
\end{tabular}
\end{table}

\begin{figure}[htbp]
\centering
\includegraphics[width=0.8\linewidth]{app_confirmar.png}
\caption{Confirmación dentro de la página, en lugar de la ventana del navegador (versión final).}
\label{fig:confirmar}
\end{figure}

\paragraph{Verificación.} Tras las correcciones, las 30 pruebas automáticas de \arch{tests/test_webapp.py} pasan sobre la versión final. El diálogo se probó además en un navegador sin interfaz durante seis rondas seguidas de borrado: cada una hizo una sola petición \arch{DELETE}, sin ventanas nativas ni excepciones.

% Commits: rama feature/julian-webapp-final, integrada en main en faa42ff
% =====================================================================
\section{Aplicación web para el radiólogo}
\label{sec:web}

Sobre el MVP, Juan Manuel Tengonó desarrolló la aplicación final en \arch{webapp/frontend/clinica/}. Se ejecuta en local, en CPU y sin internet después de la primera instalación: FastAPI sirve la API y una interfaz en HTML, CSS y JavaScript sin compilación. Se arranca con \arch{webapp/iniciar_app.bat} en Windows o con \arch{webapp/iniciar_app.sh} en Linux y macOS. Se prefirió esta tecnología frente a React porque no requiere compilar nada y mantiene el arranque de un solo comando.

\subsection{Flujo de trabajo}
\begin{enumerate}
  \item \textbf{Cargar.} Recibe archivos \arch{.IMA} o \arch{.dcm}, la carpeta del resonador o un \arch{.zip} (Figura~\ref{fig:appinicio}). Convierte a PNG con el mismo convertidor del dataset e identifica la serie T2 sagital.
  \item \textbf{Cortes.} Visor con brillo, contraste, zoom y cajas del detector. El detector recomienda los cortes donde mejor ve los cinco discos, y el usuario puede cambiarlos (Figura~\ref{fig:appcortes}).
  \item \textbf{Resultados.} Conclusión en lenguaje natural y una tarjeta por disco con la clase del ensamble, en cuántos cortes se repite y cuántos votos coinciden. En cada disco el radiólogo marca si concuerda o discrepa.
  \item \textbf{Informe.} Hallazgos, impresión, recomendación y firma. Se imprime o se guarda como PDF (Figura~\ref{fig:informe}).
\end{enumerate}
El atajo \emph{Analizar una imagen} clasifica un recorte de disco o un corte completo suelto.

\begin{figure}[htbp]
\centering
\includegraphics[width=\linewidth]{app_inicio.png}
\caption{Pantalla inicial de la versión final de la aplicación.}
\label{fig:appinicio}
\end{figure}

\subsection{Rechazo de imágenes y privacidad}
Los cuatro modelos siempre eligen una de las tres clases, así que antes de predecir se filtran las imágenes. La capa~1 rechaza imágenes en color, diminutas o de proporciones raras. La capa~2 distingue recorte de corte completo, y un corte solo se acepta si el detector halla al menos tres niveles en orden. La capa~3, una distancia de Mahalanobis sobre las características, \textbf{no está activa}: requiere calibrar un umbral con \arch{scripts/calibrar_rechazo.py}, y mientras tanto la aplicación lo advierte.

Los DICOM se borran al convertirse, y el nombre, la fecha de nacimiento y los identificadores nunca se devuelven ni se registran. Los PNG se borran tras 30 minutos sin uso o al detener el servidor.

\subsection{Correcciones a la aplicación}
\label{sec:arreglos}
Al revisar la aplicación frente a lo definido en el proyecto se hallaron fallos y restos del MVP. La Tabla~\ref{tab:arreglos} los resume, y todos los cambios se integraron en \arch{main} con el commit \arch{faa42ff}. La Figura~\ref{fig:confirmar} muestra la confirmación dentro de la página, y la comparación de las Figuras~\ref{fig:apprechazo} y~\ref{fig:appinicio} muestra el aviso reforzado.

\begin{table}[htbp]
\centering\small
\caption{Correcciones hechas a la aplicación y commit de cada una.}
\label{tab:arreglos}
\begin{tabular}{p{0.36\linewidth}p{0.40\linewidth}p{0.13\linewidth}}
\toprule
Problema & Corrección & Commit \\
\midrule
Repetir el análisis podía duplicar resultados o mezclar uno anterior & Un contador cancela los análisis atrasados y el cambio de serie reinicia el estado & \arch{b5e848f} \\
Los estudios vencidos solo se borraban al arrancar & Limpieza periódica cada minuto y borrado total al detener el servidor & \arch{35ab7e5} \\
Los avisos no explicaban los límites de la herramienta & Avisos sobre etiquetas, confianza orientativa y filtro sin calibrar & \arch{13d05f3} \\
Convivían el MVP, código sin uso y rutas que la aplicación no llama & Retiro del MVP, los casos de demostración, CORS y las rutas sobrantes & \arch{1360698} \\
Scripts y entorno seguían con el nombre del MVP & Paso a \arch{iniciar_app} y \arch{.venv-app} & \arch{be46285} \\
Borrar un estudio abría la ventana nativa del navegador & Diálogo de confirmación dentro de la página & \arch{8ba1c36} \\
El cierre atrasado de un diálogo consumía la confirmación siguiente y el estudio no se borraba & Se ignora ese cierre mientras el diálogo sigue abierto & \arch{173f44f} \\
Una confirmación cancelada quedaba pendiente y provocaba un doble borrado & Una sola confirmación pendiente a la vez & \arch{2e87449} \\
Documentación y carpetas desactualizadas & README al día; retiro de \arch{assets/} y \arch{data/external/} & \arch{0cef456}, \arch{cefe20c} \\
\bottomrule
\end{tabular}
\end{table}

\begin{figure}[htbp]
\centering
\includegraphics[width=0.8\linewidth]{app_confirmar.png}
\caption{Confirmación dentro de la página, en lugar de la ventana del navegador (versión final).}
\label{fig:confirmar}
\end{figure}

\paragraph{Verificación.} Tras las correcciones, las 30 pruebas automáticas de \arch{tests/test_webapp.py} pasan sobre la versión final. El diálogo se probó además en un navegador sin interfaz durante seis rondas seguidas de borrado: cada una hizo una sola petición \arch{DELETE}, sin ventanas nativas ni excepciones.

% Commits: rama feature/julian-webapp-final, integrada en main en faa42ff
% =====================================================================
\section{Aplicación web para el radiólogo}
\label{sec:web}

Sobre el MVP, Juan Manuel Tengonó desarrolló la aplicación final en \arch{webapp/frontend/clinica/}. Se ejecuta en local, en CPU y sin internet después de la primera instalación: FastAPI sirve la API y una interfaz en HTML, CSS y JavaScript sin compilación. Se arranca con \arch{webapp/iniciar_app.bat} en Windows o con \arch{webapp/iniciar_app.sh} en Linux y macOS. Se prefirió esta tecnología frente a React porque no requiere compilar nada y mantiene el arranque de un solo comando.

\subsection{Flujo de trabajo}
\begin{enumerate}
  \item \textbf{Cargar.} Recibe archivos \arch{.IMA} o \arch{.dcm}, la carpeta del resonador o un \arch{.zip} (Figura~\ref{fig:appinicio}). Convierte a PNG con el mismo convertidor del dataset e identifica la serie T2 sagital.
  \item \textbf{Cortes.} Visor con brillo, contraste, zoom y cajas del detector. El detector recomienda los cortes donde mejor ve los cinco discos, y el usuario puede cambiarlos (Figura~\ref{fig:appcortes}).
  \item \textbf{Resultados.} Conclusión en lenguaje natural y una tarjeta por disco con la clase del ensamble, en cuántos cortes se repite y cuántos votos coinciden. En cada disco el radiólogo marca si concuerda o discrepa.
  \item \textbf{Informe.} Hallazgos, impresión, recomendación y firma. Se imprime o se guarda como PDF (Figura~\ref{fig:informe}).
\end{enumerate}
El atajo \emph{Analizar una imagen} clasifica un recorte de disco o un corte completo suelto.

\begin{figure}[htbp]
\centering
\includegraphics[width=\linewidth]{app_inicio.png}
\caption{Pantalla inicial de la versión final de la aplicación.}
\label{fig:appinicio}
\end{figure}

\subsection{Rechazo de imágenes y privacidad}
Los cuatro modelos siempre eligen una de las tres clases, así que antes de predecir se filtran las imágenes. La capa~1 rechaza imágenes en color, diminutas o de proporciones raras. La capa~2 distingue recorte de corte completo, y un corte solo se acepta si el detector halla al menos tres niveles en orden. La capa~3, una distancia de Mahalanobis sobre las características, \textbf{no está activa}: requiere calibrar un umbral con \arch{scripts/calibrar_rechazo.py}, y mientras tanto la aplicación lo advierte.

Los DICOM se borran al convertirse, y el nombre, la fecha de nacimiento y los identificadores nunca se devuelven ni se registran. Los PNG se borran tras 30 minutos sin uso o al detener el servidor.

\subsection{Correcciones a la aplicación}
\label{sec:arreglos}
Al revisar la aplicación frente a lo definido en el proyecto se hallaron fallos y restos del MVP. La Tabla~\ref{tab:arreglos} los resume, y todos los cambios se integraron en \arch{main} con el commit \arch{faa42ff}. La Figura~\ref{fig:confirmar} muestra la confirmación dentro de la página, y la comparación de las Figuras~\ref{fig:apprechazo} y~\ref{fig:appinicio} muestra el aviso reforzado.

\begin{table}[htbp]
\centering\small
\caption{Correcciones hechas a la aplicación y commit de cada una.}
\label{tab:arreglos}
\begin{tabular}{p{0.36\linewidth}p{0.40\linewidth}p{0.13\linewidth}}
\toprule
Problema & Corrección & Commit \\
\midrule
Repetir el análisis podía duplicar resultados o mezclar uno anterior & Un contador cancela los análisis atrasados y el cambio de serie reinicia el estado & \arch{b5e848f} \\
Los estudios vencidos solo se borraban al arrancar & Limpieza periódica cada minuto y borrado total al detener el servidor & \arch{35ab7e5} \\
Los avisos no explicaban los límites de la herramienta & Avisos sobre etiquetas, confianza orientativa y filtro sin calibrar & \arch{13d05f3} \\
Convivían el MVP, código sin uso y rutas que la aplicación no llama & Retiro del MVP, los casos de demostración, CORS y las rutas sobrantes & \arch{1360698} \\
Scripts y entorno seguían con el nombre del MVP & Paso a \arch{iniciar_app} y \arch{.venv-app} & \arch{be46285} \\
Borrar un estudio abría la ventana nativa del navegador & Diálogo de confirmación dentro de la página & \arch{8ba1c36} \\
El cierre atrasado de un diálogo consumía la confirmación siguiente y el estudio no se borraba & Se ignora ese cierre mientras el diálogo sigue abierto & \arch{173f44f} \\
Una confirmación cancelada quedaba pendiente y provocaba un doble borrado & Una sola confirmación pendiente a la vez & \arch{2e87449} \\
Documentación y carpetas desactualizadas & README al día; retiro de \arch{assets/} y \arch{data/external/} & \arch{0cef456}, \arch{cefe20c} \\
\bottomrule
\end{tabular}
\end{table}

\begin{figure}[htbp]
\centering
\includegraphics[width=0.8\linewidth]{app_confirmar.png}
\caption{Confirmación dentro de la página, en lugar de la ventana del navegador (versión final).}
\label{fig:confirmar}
\end{figure}

\paragraph{Verificación.} Tras las correcciones, las 30 pruebas automáticas de \arch{tests/test_webapp.py} pasan sobre la versión final. El diálogo se probó además en un navegador sin interfaz durante seis rondas seguidas de borrado: cada una hizo una sola petición \arch{DELETE}, sin ventanas nativas ni excepciones.

% =====================================================================
% Seccion: Pruebas de la aplicacion con imagenes de distinto origen
% Fragmento SIN preambulo: % =====================================================================
% Seccion: Pruebas de la aplicacion con imagenes de distinto origen
% Fragmento SIN preambulo: % =====================================================================
% Seccion: Pruebas de la aplicacion con imagenes de distinto origen
% Fragmento SIN preambulo: \input{secciones/pruebas_aplicacion}
% El nombre del paciente se retiro de informe_generado.png a nivel de pixel.
% =====================================================================
\section{Pruebas de la aplicación y de los modelos}
\label{sec:pruebas}

Se probó el sistema completo con tres tipos de imagen: una ajena a la temática, material del conjunto de entrenamiento y un estudio real externo al dataset. Salvo que se indique lo contrario, las capturas corresponden a la versión previa a las correcciones (commit \arch{58b1a9e}); el comportamiento de clasificación no cambió, porque no se modificaron los modelos.

\subsection{Imagen ajena a la temática}
Al subir una imagen sin relación con la columna, la aplicación no la clasifica: el detector halla 0 de 5 discos y la rechaza con un mensaje claro (Figura~\ref{fig:apprechazo}). Esta capa cubre las imágenes que no son cortes; un recorte pequeño y en grises que no sea un disco pasaría este filtro y solo lo detendría la capa~3, hoy inactiva.

\begin{figure}[htbp]
\centering
\includegraphics[width=\linewidth]{app_rechazo.jpg}
\caption{Rechazo de una imagen ajena a la temática: el detector halló 0 de 5 discos.}
\label{fig:apprechazo}
\end{figure}

\subsection{Imágenes del conjunto de entrenamiento}
Primero se analizó un corte completo del conjunto de entrenamiento con el atajo \emph{Analizar una imagen}. El detector halló los 5 discos y los cuatro modelos coincidieron en cada uno (Figura~\ref{fig:apptrain}).

\begin{figure}[htbp]
\centering
\includegraphics[width=\linewidth]{app_imagen_train.jpg}
\caption{Corte completo del conjunto de entrenamiento: 5 de 5 discos detectados y cuatro modelos de acuerdo.}
\label{fig:apptrain}
\end{figure}

Después, con la versión final, se pasaron los cinco ROI originales del corte 008 del paciente 0153, que es el primer paciente del manifiesto y pertenece a entrenamiento. Los cinco coinciden con su etiqueta en los cuatro modelos (Tabla~\ref{tab:train} y Figura~\ref{fig:approi}). El resultado es el esperado, porque los modelos ya vieron estas imágenes, y por eso no mide la capacidad de generalizar: solo comprueba que la aplicación carga los modelos y conserva su comportamiento. Las cifras de generalización son las de la Sección~\ref{sec:resultados}.

\begin{table}[htbp]
\centering\small
\caption{ROI de entrenamiento (corte 0153\_008) analizados con \arch{POST /api/predecir}. Puntaje de la clase elegida; en el SVM es un puntaje normalizado con softmax, no una probabilidad.}
\label{tab:train}
\begin{tabular}{llrrrrl}
\toprule
Nivel & Etiqueta & YOLOv8s-cls & ResNet-50 & SVM & R. Forest & Ensamble \\
\midrule
L1-L2 & Normal     & 0,999 & 0,939 & 0,746 & 0,954 & Normal \\
L2-L3 & Normal     & 0,974 & 0,956 & 0,752 & 0,979 & Normal \\
L3-L4 & Normal     & 0,822 & 0,965 & 0,746 & 0,947 & Normal \\
L4-L5 & Protrusión & 1,000 & 0,940 & 0,754 & 0,956 & Protrusión \\
L5-S1 & Protrusión & 1,000 & 0,930 & 0,752 & 0,977 & Protrusión \\
\bottomrule
\end{tabular}
\end{table}

\begin{figure}[htbp]
\centering
\includegraphics[width=0.85\linewidth]{app_roi_train.png}
\caption{ROI L4-L5 de entrenamiento en la versión final: los cuatro modelos eligen Protrusión.}
\label{fig:approi}
\end{figure}

\subsection{Estudio externo al dataset}
Se cargó un estudio real que no pertenece al dataset, adquirido con un equipo GE Signa Explorer de 1,5~T, mientras que los modelos se entrenaron con un Siemens. La aplicación identificó la serie T2 sagital de 12 cortes y advirtió de esa diferencia de equipo (Figura~\ref{fig:appcortes}). El detector ubicó los cinco discos pese al cambio de fabricante. Se analizaron los cortes 6 y 7, y la persona revisora completó el informe de la Figura~\ref{fig:informe}. El nombre del paciente se retiró de la imagen.

\begin{figure}[htbp]
\centering
\includegraphics[width=\linewidth]{app_cortes_externo.jpg}
\caption{Paso Cortes con un estudio externo: cajas del detector, recomendación y aviso del equipo distinto.}
\label{fig:appcortes}
\end{figure}

\begin{figure}[htbp]
\centering
\includegraphics[width=0.9\linewidth]{informe_generado.png}
\caption{Informe generado por la aplicación para el estudio externo, con el nombre del paciente retirado.}
\label{fig:informe}
\end{figure}

La revisión concordó en 2 de los 5 discos, L4-L5 y L5-S1, y discrepó en 3. En los tres casos marcó Normal: en L1-L2 y L3-L4 el ensamble dijo Protrusión, y en L2-L3 no fue concluyente. Las discrepancias van en la dirección que anticipan los datos: una clase Protrusión mayoritaria, concentrada en los niveles bajos, y un equipo distinto al de entrenamiento. Es un solo estudio, así que no permite estimar el desempeño externo, pero muestra el uso previsto: el sistema propone, la revisión humana decide y el desacuerdo queda registrado en el informe.

\subsection{Pruebas automáticas}
Las 30 pruebas de \arch{tests/test_webapp.py} pasan sobre la versión final (\arch{30 passed}). Con los modelos reales, la API reportó estado correcto con CPU, los cuatro modelos y el detector disponibles.

\begin{table}[htbp]
\centering\small
\caption{Resumen de las pruebas.}
\label{tab:pruebas}
\begin{tabular}{p{0.34\linewidth}p{0.56\linewidth}}
\toprule
Entrada & Resultado \\
\midrule
Imagen ajena a la temática & Rechazada: 0 de 5 discos \\
Corte de entrenamiento & 5 de 5 discos detectados, cuatro modelos de acuerdo \\
5 ROI de entrenamiento & 5 de 5 coinciden con la etiqueta en los cuatro modelos \\
Estudio externo (GE, 1,5~T) & 5 discos detectados, aviso de equipo, 2 de 5 concordancias con la revisión \\
Pruebas automáticas & 30 de 30 pasan \\
\bottomrule
\end{tabular}
\end{table}

% El nombre del paciente se retiro de informe_generado.png a nivel de pixel.
% =====================================================================
\section{Pruebas de la aplicación y de los modelos}
\label{sec:pruebas}

Se probó el sistema completo con tres tipos de imagen: una ajena a la temática, material del conjunto de entrenamiento y un estudio real externo al dataset. Salvo que se indique lo contrario, las capturas corresponden a la versión previa a las correcciones (commit \arch{58b1a9e}); el comportamiento de clasificación no cambió, porque no se modificaron los modelos.

\subsection{Imagen ajena a la temática}
Al subir una imagen sin relación con la columna, la aplicación no la clasifica: el detector halla 0 de 5 discos y la rechaza con un mensaje claro (Figura~\ref{fig:apprechazo}). Esta capa cubre las imágenes que no son cortes; un recorte pequeño y en grises que no sea un disco pasaría este filtro y solo lo detendría la capa~3, hoy inactiva.

\begin{figure}[htbp]
\centering
\includegraphics[width=\linewidth]{app_rechazo.jpg}
\caption{Rechazo de una imagen ajena a la temática: el detector halló 0 de 5 discos.}
\label{fig:apprechazo}
\end{figure}

\subsection{Imágenes del conjunto de entrenamiento}
Primero se analizó un corte completo del conjunto de entrenamiento con el atajo \emph{Analizar una imagen}. El detector halló los 5 discos y los cuatro modelos coincidieron en cada uno (Figura~\ref{fig:apptrain}).

\begin{figure}[htbp]
\centering
\includegraphics[width=\linewidth]{app_imagen_train.jpg}
\caption{Corte completo del conjunto de entrenamiento: 5 de 5 discos detectados y cuatro modelos de acuerdo.}
\label{fig:apptrain}
\end{figure}

Después, con la versión final, se pasaron los cinco ROI originales del corte 008 del paciente 0153, que es el primer paciente del manifiesto y pertenece a entrenamiento. Los cinco coinciden con su etiqueta en los cuatro modelos (Tabla~\ref{tab:train} y Figura~\ref{fig:approi}). El resultado es el esperado, porque los modelos ya vieron estas imágenes, y por eso no mide la capacidad de generalizar: solo comprueba que la aplicación carga los modelos y conserva su comportamiento. Las cifras de generalización son las de la Sección~\ref{sec:resultados}.

\begin{table}[htbp]
\centering\small
\caption{ROI de entrenamiento (corte 0153\_008) analizados con \arch{POST /api/predecir}. Puntaje de la clase elegida; en el SVM es un puntaje normalizado con softmax, no una probabilidad.}
\label{tab:train}
\begin{tabular}{llrrrrl}
\toprule
Nivel & Etiqueta & YOLOv8s-cls & ResNet-50 & SVM & R. Forest & Ensamble \\
\midrule
L1-L2 & Normal     & 0,999 & 0,939 & 0,746 & 0,954 & Normal \\
L2-L3 & Normal     & 0,974 & 0,956 & 0,752 & 0,979 & Normal \\
L3-L4 & Normal     & 0,822 & 0,965 & 0,746 & 0,947 & Normal \\
L4-L5 & Protrusión & 1,000 & 0,940 & 0,754 & 0,956 & Protrusión \\
L5-S1 & Protrusión & 1,000 & 0,930 & 0,752 & 0,977 & Protrusión \\
\bottomrule
\end{tabular}
\end{table}

\begin{figure}[htbp]
\centering
\includegraphics[width=0.85\linewidth]{app_roi_train.png}
\caption{ROI L4-L5 de entrenamiento en la versión final: los cuatro modelos eligen Protrusión.}
\label{fig:approi}
\end{figure}

\subsection{Estudio externo al dataset}
Se cargó un estudio real que no pertenece al dataset, adquirido con un equipo GE Signa Explorer de 1,5~T, mientras que los modelos se entrenaron con un Siemens. La aplicación identificó la serie T2 sagital de 12 cortes y advirtió de esa diferencia de equipo (Figura~\ref{fig:appcortes}). El detector ubicó los cinco discos pese al cambio de fabricante. Se analizaron los cortes 6 y 7, y la persona revisora completó el informe de la Figura~\ref{fig:informe}. El nombre del paciente se retiró de la imagen.

\begin{figure}[htbp]
\centering
\includegraphics[width=\linewidth]{app_cortes_externo.jpg}
\caption{Paso Cortes con un estudio externo: cajas del detector, recomendación y aviso del equipo distinto.}
\label{fig:appcortes}
\end{figure}

\begin{figure}[htbp]
\centering
\includegraphics[width=0.9\linewidth]{informe_generado.png}
\caption{Informe generado por la aplicación para el estudio externo, con el nombre del paciente retirado.}
\label{fig:informe}
\end{figure}

La revisión concordó en 2 de los 5 discos, L4-L5 y L5-S1, y discrepó en 3. En los tres casos marcó Normal: en L1-L2 y L3-L4 el ensamble dijo Protrusión, y en L2-L3 no fue concluyente. Las discrepancias van en la dirección que anticipan los datos: una clase Protrusión mayoritaria, concentrada en los niveles bajos, y un equipo distinto al de entrenamiento. Es un solo estudio, así que no permite estimar el desempeño externo, pero muestra el uso previsto: el sistema propone, la revisión humana decide y el desacuerdo queda registrado en el informe.

\subsection{Pruebas automáticas}
Las 30 pruebas de \arch{tests/test_webapp.py} pasan sobre la versión final (\arch{30 passed}). Con los modelos reales, la API reportó estado correcto con CPU, los cuatro modelos y el detector disponibles.

\begin{table}[htbp]
\centering\small
\caption{Resumen de las pruebas.}
\label{tab:pruebas}
\begin{tabular}{p{0.34\linewidth}p{0.56\linewidth}}
\toprule
Entrada & Resultado \\
\midrule
Imagen ajena a la temática & Rechazada: 0 de 5 discos \\
Corte de entrenamiento & 5 de 5 discos detectados, cuatro modelos de acuerdo \\
5 ROI de entrenamiento & 5 de 5 coinciden con la etiqueta en los cuatro modelos \\
Estudio externo (GE, 1,5~T) & 5 discos detectados, aviso de equipo, 2 de 5 concordancias con la revisión \\
Pruebas automáticas & 30 de 30 pasan \\
\bottomrule
\end{tabular}
\end{table}

% El nombre del paciente se retiro de informe_generado.png a nivel de pixel.
% =====================================================================
\section{Pruebas de la aplicación y de los modelos}
\label{sec:pruebas}

Se probó el sistema completo con tres tipos de imagen: una ajena a la temática, material del conjunto de entrenamiento y un estudio real externo al dataset. Salvo que se indique lo contrario, las capturas corresponden a la versión previa a las correcciones (commit \arch{58b1a9e}); el comportamiento de clasificación no cambió, porque no se modificaron los modelos.

\subsection{Imagen ajena a la temática}
Al subir una imagen sin relación con la columna, la aplicación no la clasifica: el detector halla 0 de 5 discos y la rechaza con un mensaje claro (Figura~\ref{fig:apprechazo}). Esta capa cubre las imágenes que no son cortes; un recorte pequeño y en grises que no sea un disco pasaría este filtro y solo lo detendría la capa~3, hoy inactiva.

\begin{figure}[htbp]
\centering
\includegraphics[width=\linewidth]{app_rechazo.jpg}
\caption{Rechazo de una imagen ajena a la temática: el detector halló 0 de 5 discos.}
\label{fig:apprechazo}
\end{figure}

\subsection{Imágenes del conjunto de entrenamiento}
Primero se analizó un corte completo del conjunto de entrenamiento con el atajo \emph{Analizar una imagen}. El detector halló los 5 discos y los cuatro modelos coincidieron en cada uno (Figura~\ref{fig:apptrain}).

\begin{figure}[htbp]
\centering
\includegraphics[width=\linewidth]{app_imagen_train.jpg}
\caption{Corte completo del conjunto de entrenamiento: 5 de 5 discos detectados y cuatro modelos de acuerdo.}
\label{fig:apptrain}
\end{figure}

Después, con la versión final, se pasaron los cinco ROI originales del corte 008 del paciente 0153, que es el primer paciente del manifiesto y pertenece a entrenamiento. Los cinco coinciden con su etiqueta en los cuatro modelos (Tabla~\ref{tab:train} y Figura~\ref{fig:approi}). El resultado es el esperado, porque los modelos ya vieron estas imágenes, y por eso no mide la capacidad de generalizar: solo comprueba que la aplicación carga los modelos y conserva su comportamiento. Las cifras de generalización son las de la Sección~\ref{sec:resultados}.

\begin{table}[htbp]
\centering\small
\caption{ROI de entrenamiento (corte 0153\_008) analizados con \arch{POST /api/predecir}. Puntaje de la clase elegida; en el SVM es un puntaje normalizado con softmax, no una probabilidad.}
\label{tab:train}
\begin{tabular}{llrrrrl}
\toprule
Nivel & Etiqueta & YOLOv8s-cls & ResNet-50 & SVM & R. Forest & Ensamble \\
\midrule
L1-L2 & Normal     & 0,999 & 0,939 & 0,746 & 0,954 & Normal \\
L2-L3 & Normal     & 0,974 & 0,956 & 0,752 & 0,979 & Normal \\
L3-L4 & Normal     & 0,822 & 0,965 & 0,746 & 0,947 & Normal \\
L4-L5 & Protrusión & 1,000 & 0,940 & 0,754 & 0,956 & Protrusión \\
L5-S1 & Protrusión & 1,000 & 0,930 & 0,752 & 0,977 & Protrusión \\
\bottomrule
\end{tabular}
\end{table}

\begin{figure}[htbp]
\centering
\includegraphics[width=0.85\linewidth]{app_roi_train.png}
\caption{ROI L4-L5 de entrenamiento en la versión final: los cuatro modelos eligen Protrusión.}
\label{fig:approi}
\end{figure}

\subsection{Estudio externo al dataset}
Se cargó un estudio real que no pertenece al dataset, adquirido con un equipo GE Signa Explorer de 1,5~T, mientras que los modelos se entrenaron con un Siemens. La aplicación identificó la serie T2 sagital de 12 cortes y advirtió de esa diferencia de equipo (Figura~\ref{fig:appcortes}). El detector ubicó los cinco discos pese al cambio de fabricante. Se analizaron los cortes 6 y 7, y la persona revisora completó el informe de la Figura~\ref{fig:informe}. El nombre del paciente se retiró de la imagen.

\begin{figure}[htbp]
\centering
\includegraphics[width=\linewidth]{app_cortes_externo.jpg}
\caption{Paso Cortes con un estudio externo: cajas del detector, recomendación y aviso del equipo distinto.}
\label{fig:appcortes}
\end{figure}

\begin{figure}[htbp]
\centering
\includegraphics[width=0.9\linewidth]{informe_generado.png}
\caption{Informe generado por la aplicación para el estudio externo, con el nombre del paciente retirado.}
\label{fig:informe}
\end{figure}

La revisión concordó en 2 de los 5 discos, L4-L5 y L5-S1, y discrepó en 3. En los tres casos marcó Normal: en L1-L2 y L3-L4 el ensamble dijo Protrusión, y en L2-L3 no fue concluyente. Las discrepancias van en la dirección que anticipan los datos: una clase Protrusión mayoritaria, concentrada en los niveles bajos, y un equipo distinto al de entrenamiento. Es un solo estudio, así que no permite estimar el desempeño externo, pero muestra el uso previsto: el sistema propone, la revisión humana decide y el desacuerdo queda registrado en el informe.

\subsection{Pruebas automáticas}
Las 30 pruebas de \arch{tests/test_webapp.py} pasan sobre la versión final (\arch{30 passed}). Con los modelos reales, la API reportó estado correcto con CPU, los cuatro modelos y el detector disponibles.

\begin{table}[htbp]
\centering\small
\caption{Resumen de las pruebas.}
\label{tab:pruebas}
\begin{tabular}{p{0.34\linewidth}p{0.56\linewidth}}
\toprule
Entrada & Resultado \\
\midrule
Imagen ajena a la temática & Rechazada: 0 de 5 discos \\
Corte de entrenamiento & 5 de 5 discos detectados, cuatro modelos de acuerdo \\
5 ROI de entrenamiento & 5 de 5 coinciden con la etiqueta en los cuatro modelos \\
Estudio externo (GE, 1,5~T) & 5 discos detectados, aviso de equipo, 2 de 5 concordancias con la revisión \\
Pruebas automáticas & 30 de 30 pasan \\
\bottomrule
\end{tabular}
\end{table}

% =====================================================================
% Seccion: Conclusiones
% Fragmento SIN preambulo: % =====================================================================
% Seccion: Conclusiones
% Fragmento SIN preambulo: % =====================================================================
% Seccion: Conclusiones
% Fragmento SIN preambulo: \input{secciones/conclusiones}
% =====================================================================
\section{Conclusiones y trabajo futuro}
\label{sec:conclusiones}

Se construyó un flujo completo: de 575 carpetas públicas a un dataset de 35\,007 imágenes dividido por paciente, cuatro modelos comparados con el mismo protocolo y una aplicación local que va del DICOM al informe firmado. La ResNet-50 ajustada es el mejor modelo por validación, con 88,9\,\% de exactitud y F1 macro de 0,769 en los originales de prueba. No alcanza el 89\,\% de la rúbrica y las diferencias entre los modelos están dentro del margen de error, por lo que la elección se apoya en la validación y no en esa diferencia.

Las principales limitaciones son las etiquetas de estudiantes sin radiólogo, un único equipo de entrenamiento, la baja detección de Degeneración, el sobreajuste y el posible atajo por nivel. La prueba con un estudio externo, aunque es un solo caso, refuerza la necesidad de la revisión humana.

Como trabajo futuro se propone activar la capa~3 del filtro de rechazo, obtener etiquetas validadas por un radiólogo, evaluar con estudios de otros equipos e incorporar cortes axiales para volver a separar Hernia y Abombamiento.

% =====================================================================
\section{Conclusiones y trabajo futuro}
\label{sec:conclusiones}

Se construyó un flujo completo: de 575 carpetas públicas a un dataset de 35\,007 imágenes dividido por paciente, cuatro modelos comparados con el mismo protocolo y una aplicación local que va del DICOM al informe firmado. La ResNet-50 ajustada es el mejor modelo por validación, con 88,9\,\% de exactitud y F1 macro de 0,769 en los originales de prueba. No alcanza el 89\,\% de la rúbrica y las diferencias entre los modelos están dentro del margen de error, por lo que la elección se apoya en la validación y no en esa diferencia.

Las principales limitaciones son las etiquetas de estudiantes sin radiólogo, un único equipo de entrenamiento, la baja detección de Degeneración, el sobreajuste y el posible atajo por nivel. La prueba con un estudio externo, aunque es un solo caso, refuerza la necesidad de la revisión humana.

Como trabajo futuro se propone activar la capa~3 del filtro de rechazo, obtener etiquetas validadas por un radiólogo, evaluar con estudios de otros equipos e incorporar cortes axiales para volver a separar Hernia y Abombamiento.

% =====================================================================
\section{Conclusiones y trabajo futuro}
\label{sec:conclusiones}

Se construyó un flujo completo: de 575 carpetas públicas a un dataset de 35\,007 imágenes dividido por paciente, cuatro modelos comparados con el mismo protocolo y una aplicación local que va del DICOM al informe firmado. La ResNet-50 ajustada es el mejor modelo por validación, con 88,9\,\% de exactitud y F1 macro de 0,769 en los originales de prueba. No alcanza el 89\,\% de la rúbrica y las diferencias entre los modelos están dentro del margen de error, por lo que la elección se apoya en la validación y no en esa diferencia.

Las principales limitaciones son las etiquetas de estudiantes sin radiólogo, un único equipo de entrenamiento, la baja detección de Degeneración, el sobreajuste y el posible atajo por nivel. La prueba con un estudio externo, aunque es un solo caso, refuerza la necesidad de la revisión humana.

Como trabajo futuro se propone activar la capa~3 del filtro de rechazo, obtener etiquetas validadas por un radiólogo, evaluar con estudios de otros equipos e incorporar cortes axiales para volver a separar Hernia y Abombamiento.


\bibliographystyle{IEEEtran}
\bibliography{referencias_dataset}
\end{document}
